% Les constituants de la matière — PCT 3ème — Cours signé OnBuch+
% Programme officiel MINESEC. Compiler avec : tectonic N01-constituants-de-la-matiere.tex
\newcommand{\DOCMATIERE}{PCT}
\newcommand{\DOCNIVEAU}{3ème}
\newcommand{\DOCMODULE}{Module 1 : Chimie}
\newcommand{\DOCLECON}{N01}
\newcommand{\DOCTITRE}{Les constituants de la matière}
% OnBuch+ — préambule « cours signés » ENRICHI (compilé avec tectonic / XeLaTeX)
% Système visuel « Pop » d'OnBuch+ : fond crème (jamais blanc), bordures encre
% épaisses, ombres « dures » décalées, titres Archivo Black en capitales.
% Version MATHÉMATIQUES : graphes (pgfplots/tikz), boîtes pédagogiques
% (définition, propriété, méthode, attention, le savais-tu, exercice résolu,
% exercice, TP, bilan), page de garde et signature OnBuch+ ancrée en bas.
%
% Macros attendues dans main.tex AVANT \input{preamble} :
%   \DOCMATIERE  \DOCNIVEAU  \DOCMODULE  \DOCLECON (numéro)  \DOCTITRE
\documentclass[11pt]{article}

\usepackage{fontspec}
\usepackage[french]{babel}
\usepackage[a4paper,margin=2cm,top=3.4cm,bottom=2.8cm,headheight=1.4cm,footskip=1.3cm]{geometry}
\usepackage{amsmath,amssymb,amsfonts}
\usepackage{mathtools} % \xmapsto, \coloneqq... souvent utilisés par les modèles
\usepackage{cancel} % simplifications barrées (bilans de forces)
% \abs{z} (module, valeur absolue) et \norm{u} (norme), utilisés par les modèles
\providecommand{\abs}{}\let\abs\relax\DeclarePairedDelimiter{\abs}{\lvert}{\rvert}
\providecommand{\norm}{}\let\norm\relax\DeclarePairedDelimiter{\norm}{\lVert}{\rVert}
\usepackage[table]{xcolor} % [table] : fournit aussi \rowcolors (alternance de lignes)
\usepackage{pagecolor}
\usepackage{tikz}
\usetikzlibrary{arrows.meta,positioning,calc,shapes.geometric,shapes.misc,shapes.arrows,decorations.pathmorphing,decorations.pathreplacing,decorations.markings,patterns,shadows,mindmap,trees,fit,backgrounds,matrix,intersections,angles,quotes}
\usepackage{pgfplots}
\usepackage[version=4]{mhchem} % équations nucléaires \ce{^{235}_{92}U}
\usepackage[european,siunitx]{circuitikz} % schémas électriques
\pgfplotsset{compat=1.18}
\usepgfplotslibrary{fillbetween,groupplots} % aires entre courbes (calcul intégral)
\usepackage{enumitem}
\usepackage{fancyhdr}
% [most] charge toutes les bibliothèques tcolorbox (listings, minted, raster,
% documentation...) et gonfle inutilement la mémoire TeX sur un document long
% (~300 boîtes) : on ne charge que ce qui est réellement utilisé.
\usepackage[skins,breakable]{tcolorbox}
\usepackage{graphicx}
\usepackage{colortbl}
\usepackage{booktabs}
\usepackage{tabularx}
\usepackage{multirow}
\usepackage{diagbox} % case d'en-tête diagonale des tableaux à double entrée
% Colonnes extensibles centrée / gauche / droite (C, L, R), souvent utilisées par les modèles
\newcolumntype{C}{>{\centering\arraybackslash}X}
\newcolumntype{L}{>{\raggedright\arraybackslash}X}
\newcolumntype{R}{>{\raggedleft\arraybackslash}X}
\usepackage{array}
\usepackage{multicol}
\usepackage{siunitx}
% parse-numbers=false : accepte aussi \SI{2,5 \times 10^{-3}}{...}
\sisetup{locale=FR,output-decimal-marker={,},group-minimum-digits=4,parse-numbers=false}
\DeclareSIUnit\ohm{\ensuremath{\symup{\Omega}}} % U+2126 absent de la police
\DeclareSIUnit\division{div} % réglages d'oscilloscope (ms/div, V/div)
\DeclareSIUnit\tr{tr} % tours (vitesse de rotation, ex. \SI{1500}{\tr\per\minute})
\DeclareSIUnit\molar{M} % concentration molaire (ex. \SI{3}{\molar}, \SI{1}{\micro\molar})
\DeclareSIUnit\an{an} % année (le modèle confond parfois avec \year, un compteur TeX) — ex. \SI{5}{\kilo\meter\per\an}
\DeclareSIUnit\FCFA{FCFA} % franc CFA (ex. \SI{500}{\FCFA\per\kilo\gram})
\DeclareSIUnit\degreeBrix{\textdegree Brix} % teneur en sucre d'un sirop/jus (réfractométrie)
\DeclareSIUnit\month{mois} % le modèle utilise parfois \month, un compteur TeX, comme unité de durée
\DeclareSIUnit\microliter{\micro\liter} % le modèle écrit parfois \microliter en un seul token
\DeclareSIUnit\million{millions} % dénombrement (ex. \SI{15}{\million\per\mL} spermatozoïdes)
\usepackage{qrcode}
% \degree : commande fréquemment utilisée par les modèles pour un angle ou une
% température en dehors de \SI{}{} (non fournie par les paquets chargés).
\providecommand{\degree}{^{\circ}}
% \qed (fin de démonstration), utilisé par les modèles sans amsthm
\providecommand{\qed}{\hfill$\square$}
% \barre : double barre des valeurs interdites dans un tableau de variations (tkz-tab)
\providecommand{\barre}{\ensuremath{\|}}
% environnement proof (amsthm non chargé) utilisé par les modèles
\ifdefined\proof\else\newenvironment{proof}[1][Démonstration]{\par\noindent\textit{#1.}\ }{\qed\par}\fi
\usepackage{lastpage}
\usepackage{hyperref}
\hypersetup{hidelinks}

% ============================================================
% Palette « Pop » OnBuch+ — mêmes hex que la classe P (plus_ui.dart)
% ============================================================
\definecolor{poporange}{HTML}{F59321}
\definecolor{popdark}{HTML}{E07A0C}
\definecolor{popcream}{HTML}{FFF6E8}
\definecolor{popink}{HTML}{1C1714}
\definecolor{poppurple}{HTML}{7A5AE0}
\definecolor{popgreen}{HTML}{2E9E6A}
\definecolor{poppink}{HTML}{E0457B}
\definecolor{popblue}{HTML}{2F7DE1}
\definecolor{popgold}{HTML}{C98A12}
\definecolor{popmuted}{HTML}{8A7F76}
\definecolor{popline}{HTML}{EADFD2}
\definecolor{popgray}{HTML}{8A7F76}
\colorlet{popgrey}{popgray}
% Alias tolérants (le modèle nomme parfois les couleurs différemment)
\colorlet{popwhite}{white}
\colorlet{popblack}{popink}
\colorlet{popred}{poppink}
\colorlet{popyellow}{popgold}
% Teintes claires (fonds de tableaux/encadrés) — le modèle nomme parfois
% les couleurs avec un suffixe L (ex. popblueL) sans les avoir définies.
\colorlet{poporangeL}{poporange!20}
\colorlet{popdarkL}{popdark!20}
\colorlet{poppurpleL}{poppurple!20}
\colorlet{popgreenL}{popgreen!20}
\colorlet{poppinkL}{poppink!20}
\colorlet{popblueL}{popblue!20}
\colorlet{popgoldL}{popgold!20}
\colorlet{popredL}{popred!20}
\colorlet{popgrayL}{popgray!20}
% Teintes claires pour fonds de tableaux / zones
\colorlet{popblueL}{popblue!10!white}
\colorlet{poporangeL}{poporange!14!white}
\colorlet{popgreenL}{popgreen!12!white}
\colorlet{poppurpleL}{poppurple!12!white}
\colorlet{poppinkL}{poppink!12!white}
\colorlet{popgoldL}{popgold!14!white}

\pagecolor{popcream}

% ============================================================
% Typographie
% ============================================================
\defaultfontfeatures{Ligatures=TeX}
\setmainfont{PlusJakartaSans-Medium.ttf}[Path=../../../fonts/,BoldFont=PlusJakartaSans-Bold.ttf]
% Maths en sans-serif (Fira Math) avec chiffres et lettres droites de Plus
% Jakarta, pour que les formules se fondent dans le texte.
\usepackage[math-style=ISO,bold-style=ISO]{unicode-math}
\setmathfont{FiraMath-Regular.otf}[Path=../../../fonts/]
\setmathfont{PlusJakartaSans-Medium.ttf}[Path=../../../fonts/,range={up/num,up/{Latin,latin},\mathup}]
\usepackage{icomma} % virgule décimale sans espace en mode math
% Caractères mathématiques Unicode absents des polices de texte (Plus Jakarta,
% Archivo Black) : rendus par leur équivalent mathématique, en texte comme en
% formule — sinon ils s'affichent en carré vide (ex. « Arithmétique dans ℤ »).
\usepackage{newunicodechar}
\newunicodechar{ℝ}{\ensuremath{\mathbb{R}}}
\newunicodechar{ℤ}{\ensuremath{\mathbb{Z}}}
\newunicodechar{ℂ}{\ensuremath{\mathbb{C}}}
\newunicodechar{ℕ}{\ensuremath{\mathbb{N}}}
\newunicodechar{ℚ}{\ensuremath{\mathbb{Q}}}
\newunicodechar{𝔻}{\ensuremath{\mathbb{D}}}
\newunicodechar{α}{\ensuremath{\alpha}}
\newunicodechar{β}{\ensuremath{\beta}}
\newunicodechar{γ}{\ensuremath{\gamma}}
\newunicodechar{δ}{\ensuremath{\delta}}
\newunicodechar{ε}{\ensuremath{\varepsilon}}
\newunicodechar{θ}{\ensuremath{\theta}}
\newunicodechar{λ}{\ensuremath{\lambda}}
\newunicodechar{μ}{\ensuremath{\mu}}
\newunicodechar{σ}{\ensuremath{\sigma}}
\newunicodechar{τ}{\ensuremath{\tau}}
\newunicodechar{φ}{\ensuremath{\varphi}}
\newunicodechar{ψ}{\ensuremath{\psi}}
\newunicodechar{ω}{\ensuremath{\omega}}
\newunicodechar{Σ}{\ensuremath{\Sigma}}
% majuscules produites par \MakeUppercase dans les titres (α-aminés -> Α)
\newunicodechar{Α}{\ensuremath{\alpha}}
\newunicodechar{Β}{\ensuremath{\beta}}
\newunicodechar{Γ}{\ensuremath{\Gamma}}
\newunicodechar{Δ}{\ensuremath{\Delta}}
\newunicodechar{Ε}{\ensuremath{\varepsilon}}
\newunicodechar{Θ}{\ensuremath{\Theta}}
\newunicodechar{Λ}{\ensuremath{\Lambda}}
\newunicodechar{Μ}{\ensuremath{\mu}}
\newunicodechar{Τ}{\ensuremath{\tau}}
\newunicodechar{Φ}{\ensuremath{\Phi}}
\newunicodechar{Ψ}{\ensuremath{\Psi}}
\newunicodechar{Ω}{\ensuremath{\Omega}}
\newunicodechar{↦}{\ensuremath{\mapsto}}
\newunicodechar{∈}{\ensuremath{\in}}
\newunicodechar{∉}{\ensuremath{\notin}}
\newunicodechar{∧}{\ensuremath{\wedge}}
\newunicodechar{⇔}{\ensuremath{\Leftrightarrow}}
\newunicodechar{⟺}{\ensuremath{\Longleftrightarrow}}
\newunicodechar{⇒}{\ensuremath{\Rightarrow}}
\newunicodechar{⟹}{\ensuremath{\Longrightarrow}}
\newunicodechar{∘}{\ensuremath{\circ}}
\newunicodechar{∪}{\ensuremath{\cup}}
\newunicodechar{∩}{\ensuremath{\cap}}
\newunicodechar{≡}{\ensuremath{\equiv}}
\newunicodechar{⊂}{\ensuremath{\subset}}
\newunicodechar{⊥}{\ensuremath{\perp}}
\newunicodechar{∃}{\ensuremath{\exists}}
\newunicodechar{∀}{\ensuremath{\forall}}
\newunicodechar{∛}{\ensuremath{\sqrt[3]{\,}}}
\newunicodechar{∜}{\ensuremath{\sqrt[4]{\,}}}
\newunicodechar{⃗}{\ensuremath{{}^{\to}}}
\newunicodechar{ⁿ}{\ifmmode^{n}\else\textsuperscript{\ensuremath{n}}\fi}
\newunicodechar{ˣ}{\ifmmode^{x}\else\textsuperscript{\ensuremath{x}}\fi}
\newunicodechar{ᵇ}{\ifmmode^{b}\else\textsuperscript{\ensuremath{b}}\fi}
\newunicodechar{ʳ}{\ifmmode^{r}\else\textsuperscript{\ensuremath{r}}\fi}
\newunicodechar{ᵘ}{\ifmmode^{u}\else\textsuperscript{\ensuremath{u}}\fi}
\newunicodechar{ʸ}{\ifmmode^{y}\else\textsuperscript{\ensuremath{y}}\fi}
\newunicodechar{ᵃ}{\ifmmode^{a}\else\textsuperscript{\ensuremath{a}}\fi}
\newunicodechar{ᵉ}{\ifmmode^{e}\else\textsuperscript{\ensuremath{e}}\fi}
\newunicodechar{ᵏ}{\ifmmode^{k}\else\textsuperscript{\ensuremath{k}}\fi}
\newunicodechar{ᵖ}{\ifmmode^{p}\else\textsuperscript{\ensuremath{p}}\fi}
\newunicodechar{ᴺ}{\ifmmode^{N}\else\textsuperscript{\ensuremath{N}}\fi}
\newunicodechar{ᶜ}{\ifmmode^{c}\else\textsuperscript{\ensuremath{c}}\fi}
\newunicodechar{ʰ}{\ifmmode^{h}\else\textsuperscript{\ensuremath{h}}\fi}
\newunicodechar{ᵠ}{\ifmmode^{q}\else\textsuperscript{\ensuremath{q}}\fi}
\newunicodechar{ₙ}{\ifmmode_{n}\else\textsubscript{\ensuremath{n}}\fi}
\newunicodechar{ₐ}{\ifmmode_{a}\else\textsubscript{\ensuremath{a}}\fi}
\newunicodechar{ᵢ}{\ifmmode_{i}\else\textsubscript{\ensuremath{i}}\fi}
\newunicodechar{ᵣ}{\ifmmode_{r}\else\textsubscript{\ensuremath{r}}\fi}
\newunicodechar{ₖ}{\ifmmode_{k}\else\textsubscript{\ensuremath{k}}\fi}
\newunicodechar{ᵦ}{\ifmmode_{\beta}\else\textsubscript{\ensuremath{\beta}}\fi}
\newunicodechar{ₓ}{\ifmmode_{x}\else\textsubscript{\ensuremath{x}}\fi}
\newfontfamily\popdisplay{ArchivoBlack-Regular.ttf}[Path=../../../fonts/]
\newfontfamily\poplogo{SpaceGrotesk-Bold.ttf}[Path=../../../fonts/]
\newfontfamily\popsemi{PlusJakartaSans-SemiBold.ttf}[Path=../../../fonts/]
\newfontfamily\popxbold{PlusJakartaSans-ExtraBold.ttf}[Path=../../../fonts/]
\newfontfamily\popxboldls{PlusJakartaSans-ExtraBold.ttf}[Path=../../../fonts/,LetterSpace=3.0]

\setlist[enumerate]{leftmargin=1.7em,itemsep=5pt,topsep=4pt}
\setlist[itemize]{leftmargin=1.7em,itemsep=4pt,topsep=4pt,label={\color{poporange}\textbullet}}
\setlength{\parindent}{0pt}
\setlength{\parskip}{5pt}
\renewcommand{\arraystretch}{1.25}

% pgfplots : style Pop par défaut
\pgfplotsset{
  every axis/.append style={axis line style={popink,line width=1pt},tick style={popink},
    grid style={popline,line width=0.5pt},label style={font=\small},tick label style={font=\footnotesize}},
  popcurve/.style={poporange,line width=1.8pt,no markers,smooth},
}
% Même style disponible dans un \draw TikZ simple (hors pgfplots)
\tikzset{popcurve/.style={poporange,line width=1.8pt}}
\tikzset{popgrid/.style={popline,very thin}}
\colorlet{popcurve}{poporange} % le nom du style est parfois utilisé comme couleur

% ============================================================
% Pill / squiggle
% ============================================================
\newcommand{\poppill}[3]{%
  \tikz[baseline=-0.55ex]\node[draw=popink,line width=1.2pt,rounded corners=2.6pt,%
    fill=#2,inner xsep=6.5pt,inner ysep=3pt]{%
    \popxboldls\fontsize{7}{7}\selectfont\color{#3}\MakeUppercase{#1}};%
}
\newcommand{\popsquiggle}[1]{%
  \tikz[baseline=-0.4ex]{
    \draw[#1,line width=2.6pt,line cap=round]
      plot [smooth] coordinates {(0,0.09) (0.34,0.01) (0.68,0.21) (1.02,0.01) (1.36,0.10)};
  }%
}
% Mot-clé surligné (marqueur orange sous le texte)
\newcommand{\cle}[1]{{\popxbold\color{popdark}#1}}

% ============================================================
% En-tête / pied
% ============================================================
\pagestyle{fancy}
\fancyhf{}
\renewcommand{\headrulewidth}{0pt}
\renewcommand{\footrulewidth}{0pt}
\lhead{\raisebox{-0.15ex}{\poplogo\fontsize{13}{13}\selectfont\color{popink}OnBuch\color{poporange}+}}
\rhead{\poppill{\DOCMATIERE\ \textbullet\ \DOCNIVEAU\ \textbullet\ Leçon \DOCLECON}{white}{popink}}
\lfoot{\popxbold\fontsize{7.6}{7.6}\selectfont\color{popmuted}\MakeUppercase{Cours signé OnBuch+}\ \textbullet\ {\popsemi\DOCTITRE}}
\rfoot{\tikz[baseline=-0.6ex]\node[rounded corners=8.5pt,fill=popink,minimum height=17pt,minimum width=17pt,inner xsep=5pt,inner sep=0]{\popxbold\fontsize{8}{8}\selectfont\color{popcream}\thepage\,/\,\pageref*{LastPage}};}

% ============================================================
% Titres
% ============================================================
\newcommand{\chaptitre}[2]{%
  \begin{tcolorbox}[enhanced,colback=poporange,colframe=popink,boxrule=2.6pt,arc=11pt,
      drop shadow=popink,left=15pt,right=15pt,top=12pt,bottom=11pt,before skip=3pt,after skip=18pt]
    \poppill{#2}{popink}{popcream}\par\vspace{9pt}
    {\raggedright\hyphenpenalty=10000\exhyphenpenalty=10000\popdisplay\fontsize{22}{24}\selectfont\color{popcream}\MakeUppercase{#1}\par}
  \end{tcolorbox}
}
% Section numérotée : pastille orange avec le numéro + titre Archivo
\newcounter{popsec}
\newcommand{\coursec}[1]{%
  \stepcounter{popsec}\setcounter{popsub}{0}%
  \par\vspace{16pt}\noindent
  \begin{minipage}{\linewidth}
  \tikz[baseline=-0.7ex]\node[draw=popink,line width=1.6pt,fill=poporange,rounded corners=4pt,minimum size=22pt,inner sep=0]{\popdisplay\fontsize{12}{12}\selectfont\color{popcream}\thepopsec};%
  \hspace{8pt}{\popdisplay\fontsize{15}{17}\selectfont\color{popink}\MakeUppercase{#1}}\par
  \vspace{4pt}\noindent{\color{poporange}\rule{36pt}{3.6pt}}
  \end{minipage}\par\nopagebreak\vspace{8pt}
}
\newcounter{popsub}
\newcommand{\courssub}[1]{%
  \stepcounter{popsub}%
  \par\vspace{9pt}\noindent{\popxbold\fontsize{11.5}{13}\selectfont\color{popink}{\color{poporange}\thepopsec.\thepopsub}\hspace{6pt}#1}\par\nopagebreak\vspace{4pt}
}

% ============================================================
% Boîtes pédagogiques — même recette : fond blanc, bordure épaisse colorée,
% ombre dure encre, badge plein. Titre optionnel [#1].
% ============================================================
\tcbset{popbase/.style={breakable,enhanced,colback=white,boxrule=2.3pt,arc=9pt,
  shadow={3pt}{-3pt}{0pt}{popink},left=14pt,right=14pt,top=11pt,bottom=11pt,before skip=11pt,after skip=15pt,
  fontupper=\popsemi\fontsize{10.6}{14.5}\selectfont\color{popink},
  fontlower=\popsemi\fontsize{10.6}{14.5}\selectfont\color{popink}}}
% Coche dessinée (le glyphe ✓ n'existe pas dans les polices du document)
\newcommand{\popcheck}[1]{\tikz[baseline=-0.5ex]\draw[#1,line width=1.6pt,line cap=round,line join=round](-0.13,0.0)--(-0.04,-0.09)--(0.13,0.1);}
\providecommand{\coche}{\popcheck{popgreen}}
\providecommand{\resultat}[1]{\par\smallskip\noindent\fcolorbox{popgreen}{popgreenL}{\parbox{\dimexpr\linewidth-2\fboxsep-2\fboxrule\relax}{\textbf{Résultat.} #1}}\par\smallskip}
\newcommand{\popboxhead}[3]{\poppill{#1}{#2}{white}\ifx\relax#3\relax\else\hspace{6pt}{\popxbold\fontsize{10.6}{12}\selectfont\color{popink}#3}\fi\par\vspace{7pt}}

\newtcolorbox{aretenir}[1][]{popbase,colframe=popgreen,before upper={\popboxhead{À retenir}{popgreen}{#1}}}
\newtcolorbox{exemplebox}[1][]{popbase,colframe=poporange,before upper={\popboxhead{Exemple}{poporange}{#1}}}
\newtcolorbox{definition}[1][]{popbase,colframe=popblue,before upper={\popboxhead{Définition}{popblue}{#1}}}
\newtcolorbox{propriete}[1][]{popbase,colframe=poppurple,before upper={\popboxhead{Propriété}{poppurple}{#1}}}
\newtcolorbox{methode}[1][]{popbase,colframe=popgold,before upper={\popboxhead{Méthode}{popgold}{#1}}}
\newtcolorbox{attention}[1][]{popbase,colframe=poppink,before upper={\popboxhead{Attention}{poppink}{#1}}}
\newtcolorbox{experience}[1][]{popbase,colframe=popblue,colback=popblueL,before upper={\popboxhead{Expérience}{popblue}{#1}}}
% « Le savais-tu ? » : carte violette pleine (contexte, culture, Cameroun)
\newtcolorbox{savaistu}[1][]{popbase,colback=poppurple,colframe=popink,
  fontupper=\popsemi\fontsize{10.4}{14.2}\selectfont\color{white},
  before upper={\poppill{Le savais-tu\,?}{white}{poppurple}\ifx\relax#1\relax\else\hspace{6pt}{\popxbold\color{white}#1}\fi\par\vspace{7pt}}}
% Exercice résolu : énoncé en haut, \tcblower puis la solution sur fond crème
\newcounter{popexo}\newcounter{popexr}
\newtcolorbox{exoresolu}[1][]{popbase,colframe=poporange,colbacklower=popcream,
  segmentation style={popink,line width=1.2pt,dashed},
  before upper={\stepcounter{popexr}\popboxhead{Exercice résolu \thepopexr}{poporange}{#1}},
  before lower={\poppill{Solution}{popink}{popcream}\par\vspace{6pt}}}
% Exercice (non résolu en place) — la correction est dans la partie Corrigés
\newtcolorbox{exercice}[1][]{popbase,colframe=popink,
  before upper={\stepcounter{popexo}\popboxhead{Exercice \thepopexo}{popink}{#1}}}
% Corrigé
\newtcolorbox{corrige}[1][]{popbase,colframe=popgreen,colback=popgreenL,
  before upper={\popboxhead{Corrigé}{popgreen}{#1}}}
% Objectifs / prérequis / bilan
\newtcolorbox{objectifs}{popbase,colframe=popink,colback=poporangeL,
  before upper={\popboxhead{Objectifs de la leçon}{popink}{}}}
\newtcolorbox{prerequis}{popbase,colframe=popink,colback=popblueL,
  before upper={\popboxhead{Prérequis}{popblue}{}}}
\newtcolorbox{bilan}{popbase,colframe=popink,colback=white,boxrule=2.8pt,
  before upper={\popboxhead{Fiche bilan}{poporange}{L'essentiel à connaître pour l'examen}}}
% Figure encadrée avec légende
\newtcolorbox{popfigure}[1][]{enhanced,breakable=false,colback=white,colframe=popline,boxrule=1.4pt,arc=8pt,
  left=10pt,right=10pt,top=10pt,bottom=8pt,before skip=10pt,after skip=12pt,halign=center,
  lower separated=false,halign lower=center,
  fontlower=\popsemi\fontsize{9}{11.5}\selectfont\color{popmuted},
  #1}
\newcommand{\legende}[1]{\par\vspace{6pt}{\popsemi\fontsize{9}{11.5}\selectfont\color{popmuted}#1}}

% ============================================================
% Signature OnBuch+
% ============================================================
\newcommand{\poplogoinline}[1]{{\poplogo\fontsize{#1}{#1}\selectfont\color{popink}OnBuch\color{poporange}+}}
\newcommand{\signatureonbuch}{%
  \begin{tcolorbox}[enhanced,colback=popink,colframe=popink,boxrule=2.2pt,arc=10pt,
      drop shadow=poporange,left=14pt,right=14pt,top=10pt,bottom=10pt,before skip=0pt,after skip=0pt]
    \tikz[baseline=-0.7ex]\node[circle,fill=popgreen,draw=popcream,line width=1.3pt,minimum size=22pt,inner sep=0]{\popcheck{white}};%
    \hspace{9pt}\begin{minipage}[c]{0.62\linewidth}
      {\popxbold\fontsize{12}{13}\selectfont\color{popcream}Signé\ \textbullet\ {\poplogo OnBuch\color{poporange}+}}\par\vspace{2pt}
      {\raggedright\popsemi\fontsize{8}{10.5}\selectfont\color{popline}\MakeUppercase{Équipe pédagogique OnBuch+}\\\MakeUppercase{Conforme au programme officiel MINESEC}\par}
    \end{minipage}\hfill
    {\poplogo\fontsize{22}{22}\selectfont\color{popcream}OnBuch\color{poporange}+}
  \end{tcolorbox}%
}

% ============================================================
% Page de garde
% #1 titre · #2 matière · #3 niveau · #4 module · #5 n° leçon · #6 durée ·
% #7 description / objectifs (texte ou itemize)
% ============================================================
\newcommand{\pagegarde}[7]{%
  \thispagestyle{empty}
  \newgeometry{a4paper,margin=1.7cm,top=1.6cm,bottom=1.4cm}%
  \begin{tikzpicture}[remember picture,overlay]
    \fill[poporange!18] ([xshift=-3.2cm,yshift=-3.6cm]current page.north east) circle (4.6cm);
    \fill[poppurple!14] ([xshift=2.4cm,yshift=7.6cm]current page.south west) circle (3.8cm);
  \end{tikzpicture}%
  {\poplogo\fontsize{30}{30}\selectfont\color{popink}OnBuch\color{poporange}+}\hfill
  \raisebox{0.6ex}{\poppill{Cours signé \textbullet\ Premium}{popink}{popcream}}\par
  \vspace{1pt}\popsquiggle{poppurple}\par
  \vspace{8pt}
  \begin{tcolorbox}[enhanced,colback=poporange,colframe=popink,boxrule=2.8pt,arc=13pt,
      drop shadow=popink,left=18pt,right=18pt,top=12pt,bottom=13pt,after skip=10pt]
    \poppill{#2 \textbullet\ #3}{popink}{popcream}\hspace{5pt}\poppill{#4}{white}{popink}\par\vspace{8pt}
    {\popdisplay\fontsize{14}{14}\selectfont\color{popink}LEÇON #5}\par\vspace{4pt}
    {\raggedright\hyphenpenalty=10000\exhyphenpenalty=10000\popdisplay\fontsize{25}{27}\selectfont\color{popcream}\MakeUppercase{#1}\par}
    \vspace{8pt}
    {\popxbold\fontsize{9.5}{9.5}\selectfont\color{popink}\MakeUppercase{Durée conseillée : #6}}
  \end{tcolorbox}
  \begin{tcolorbox}[enhanced,colback=white,colframe=popink,boxrule=2.2pt,arc=10pt,
      drop shadow=popink,left=16pt,right=16pt,top=10pt,bottom=10pt,after skip=8pt]
    \poppill{Dans cette leçon}{poppurple}{white}\par\vspace{5pt}
    \popsemi\fontsize{9.3}{12.6}\selectfont\color{popink}#7
  \end{tcolorbox}
  \noindent\begin{minipage}[t]{0.487\linewidth}
    \begin{tcolorbox}[enhanced,colback=popgreen,colframe=popink,boxrule=2.2pt,arc=10pt,
        drop shadow=popink,left=11pt,right=11pt,top=10pt,bottom=10pt,height=3.1cm,valign=center]
      \tcbox[colback=white,colframe=popink,boxrule=1.3pt,arc=5pt,boxsep=0pt,nobeforeafter,
        left=3pt,right=3pt,top=3pt,bottom=3pt,valign=center]{\qrcode[height=1.9cm]{https://play.google.com/store/apps/details?id=com.luvvix.onbuch&hl=fr}}%
      \hspace{8pt}\begin{minipage}[c]{0.5\linewidth}
        {\popdisplay\fontsize{11}{12.5}\selectfont\color{white}\MakeUppercase{Télécharge OnBuch}}\par\vspace{4pt}
        {\raggedright\popsemi\fontsize{8.4}{11}\selectfont\color{white}Gratuit \textbullet\ Google Play\\Tuteur IA, annales, concours, résultats\par}
      \end{minipage}
    \end{tcolorbox}
  \end{minipage}\hfill
  \begin{minipage}[t]{0.487\linewidth}
    \begin{tcolorbox}[enhanced,colback=poppurple,colframe=popink,boxrule=2.2pt,arc=10pt,
        drop shadow=popink,left=11pt,right=11pt,top=10pt,bottom=10pt,height=3.1cm,valign=center]
      \tikz[baseline=-0.7ex]\node[circle,fill=white,draw=popink,line width=1.3pt,minimum size=1.6cm,inner sep=0]{\popdisplay\fontsize{24}{24}\selectfont\color{poppurple}+};%
      \hspace{8pt}\begin{minipage}[c]{0.6\linewidth}
        {\popdisplay\fontsize{11}{12.5}\selectfont\color{white}\MakeUppercase{Passe à OnBuch+}}\par\vspace{4pt}
        {\raggedright\popsemi\fontsize{8.4}{11}\selectfont\color{white}Tous les cours signés, Tuteur IA illimité, fiches PDF — onglet OnBuch+ de l'app\par}
      \end{minipage}
    \end{tcolorbox}
  \end{minipage}
  \vspace{8pt}
  \popcontenu
  \vfill
  \signatureonbuch
  \restoregeometry
  \newpage
}

% Rangée « Ce cours contient » (page de garde)
\newcommand{\poptile}[3]{% #1 couleur #2 gros texte #3 légende
  \begin{tcolorbox}[enhanced,colback=#1,colframe=popink,boxrule=2pt,arc=9pt,shadow={2.5pt}{-2.5pt}{0pt}{popink},
      left=6pt,right=6pt,top=9pt,bottom=9pt,halign=center,valign=center,height=1.75cm,width=\linewidth,nobeforeafter]
    {\popdisplay\fontsize{12.5}{13}\selectfont\color{white}\MakeUppercase{#2}}\par\vspace{3pt}
    {\centering\popsemi\fontsize{7.8}{9.5}\selectfont\color{white}#3\par}
  \end{tcolorbox}}
\newcommand{\popcontenu}{%
  \par\noindent{\popxboldls\fontsize{7.5}{7.5}\selectfont\color{popmuted}CE COURS SIGNÉ CONTIENT}\par\vspace{7pt}
  \noindent\begin{minipage}[t]{0.235\linewidth}\poptile{popblue}{Cours}{complet, illustré, pas à pas}\end{minipage}\hfill
  \begin{minipage}[t]{0.235\linewidth}\poptile{poporange}{Méthodes}{et exercices résolus}\end{minipage}\hfill
  \begin{minipage}[t]{0.235\linewidth}\poptile{poppink}{Exercices}{type examen + corrigés}\end{minipage}\hfill
  \begin{minipage}[t]{0.235\linewidth}\poptile{popgreen}{Fiche bilan}{l'essentiel à retenir}\end{minipage}\par}

% Sommaire
\newtcolorbox{sommaire}{enhanced,colback=white,colframe=popink,boxrule=2.2pt,arc=10pt,
  drop shadow=popink,left=18pt,right=18pt,top=13pt,bottom=13pt,before skip=6pt,after skip=20pt,
  before upper={\poppill{Sommaire}{poppurple}{white}\par\vspace{9pt}\popsemi\fontsize{10.6}{14.5}\selectfont\color{popink}}}

% Signature de fin de cours — poussée tout en bas de la dernière page
\newcommand{\findecours}{%
  \par\vfill\nopagebreak
  \begin{center}\popsquiggle{poporange}\end{center}\vspace{4pt}
  \signatureonbuch
}
\AtBeginDocument{\def\llbracket{\mathopen{[\mkern-3.5mu[}}\def\rrbracket{\mathclose{]\mkern-3.5mu]}}} % intervalles d'entiers (glyphe ⟦ absent de la police)
% Glyphes absents de Fira Math : redéfinis à partir de glyphes disponibles
\AtBeginDocument{%
  \def\setminus{\mathbin{\backslash}}\let\smallsetminus\setminus
  \def\vdots{\mathord{\vbox{\baselineskip3.2pt\lineskiplimit0pt\kern4pt\hbox{.}\hbox{.}\hbox{.}}}}%
  \def\ddots{\mathinner{\mkern1mu\raise6.5pt\hbox{.}\mkern2mu\raise3.6pt\hbox{.}\mkern2mu\raise0.7pt\hbox{.}\mkern1mu}}%
  \def\triangle{\mathord{\scalebox{0.9}{$\symup{\Delta}$}}}%
  \def\nabla{\mathord{\raisebox{\depth}{\scalebox{1}[-1]{$\symup{\Delta}$}}}}%
  \def\checkmark{\popcheck{popdark}}%
  \def\star{\mathbin{\ast}}\def\bigstar{\mathord{\ast}}%
  \def\diamond{\mathbin{\scalebox{0.6}{$\Diamond$}}}%
  \def\bigcup{\mathop{\mathchoice{\raisebox{-0.3em}{\scalebox{1.7}{$\cup$}}}{\scalebox{1.25}{$\cup$}}{\cup}{\cup}}\displaylimits}%
  \def\bigcap{\mathop{\mathchoice{\raisebox{-0.3em}{\scalebox{1.7}{$\cap$}}}{\scalebox{1.25}{$\cap$}}{\cap}{\cap}}\displaylimits}%
  \def\lesseqqgtr{\mathrel{\lessgtr}}\def\bigoplus{\mathop{\oplus}\displaylimits}%
}
\providecommand{\ssi}{si et seulement si} % abréviation fréquente
\AtBeginDocument{\providecommand{\wideparen}{\overparen}} % arc de cercle

\begin{document}

\pagegarde{Les constituants de la matière}{PCT}{3ème}{Module 1 : Chimie}{N01}{2 séances (1 h chacune)}{Cette leçon d'ouverture de l'année de 3ème fait découvrir les constituants microscopiques de la matière : l'atome (noyau et électrons, modèle de Bohr), puis les édifices qu'il forme : molécules et ions. Elle fournit les outils indispensables (symboles, formules) pour toute la chimie de l'année et du BEPC.
\vspace{4pt}
\begin{multicols}{2}\small
\begin{itemize}
\item À la fin de cette leçon, je sais décrire la constitution d'un atome (noyau, protons, neutrons, électrons) et citer ses ordres de grandeur.
\item Je sais interpréter les expériences historiques (Crookes, Rutherford, électrolyse de l'eau) pour établir l'existence des constituants de la matière.
\item Je sais écrire le symbole d'un élément chimique et représenter un atome simple par le modèle de Bohr.
\item Je sais définir une molécule, écrire et interpréter sa formule chimique (H₂O, CO₂, O₂...).
\item Je sais définir un ion, distinguer cation et anion, et écrire la formule des ions simples (Na⁺, Cl⁻, Ca²⁺) et polyatomiques (SO₄²⁻, OH⁻).
\item Je sais vérifier l'électroneutralité d'un composé ionique à partir des formules des ions.
\end{itemize}\end{multicols}}

\begin{sommaire}
\begin{multicols}{2}
\begin{enumerate}
\item L'atome : brique de la matière
\item Symboles des éléments et modèle de Bohr
\item Les molécules
\item Les ions : ions simples
\item Ions polyatomiques et composés ioniques
\item Synthèse : de l'atome aux matériaux du quotidien
\item Activité d'intégration
\item Méthodes et exercices résolus
\item Exercices
\item Corrigés des exercices
\item Fiche bilan
\end{enumerate}
\end{multicols}
\end{sommaire}

\begin{objectifs}
\begin{itemize}
\item À la fin de cette leçon, je sais décrire la constitution d'un atome (noyau, protons, neutrons, électrons) et citer ses ordres de grandeur.
\item Je sais interpréter les expériences historiques (Crookes, Rutherford, électrolyse de l'eau) pour établir l'existence des constituants de la matière.
\item Je sais écrire le symbole d'un élément chimique et représenter un atome simple par le modèle de Bohr.
\item Je sais définir une molécule, écrire et interpréter sa formule chimique (H₂O, CO₂, O₂...).
\item Je sais définir un ion, distinguer cation et anion, et écrire la formule des ions simples (Na⁺, Cl⁻, Ca²⁺) et polyatomiques (SO₄²⁻, OH⁻).
\item Je sais vérifier l'électroneutralité d'un composé ionique à partir des formules des ions.
\item Je sais identifier les constituants de la matière dans des situations concrètes du quotidien camerounais (sel de cuisine, eau minérale, soluté de réhydratation).
\item Je sais appliquer les consignes de sécurité lors de la manipulation de produits chimiques et d'appareils électriques en salle de TP.
\end{itemize}
\end{objectifs}

\begin{prerequis}
\begin{itemize}
\item Notion de corps pur et de mélange (classe de 4ème).
\item Éléments chimiques courants et leurs symboles rencontrés en 4ème (H, O, C, Fe...).
\item Notion de courant électrique et de circuit simple (physique, 4ème).
\item Écriture et lecture de formules simples rencontrées au quotidien (H₂O, CO₂).
\item Consignes générales de sécurité en salle de sciences.
\end{itemize}
\end{prerequis}

\newpage

\coursec{L'atome : brique de la matière}

Au marché de Mokolo, une maman dissout du sel dans la sauce : le sel « disparaît » mais le goût reste. À l'hôpital de district, on soigne une déshydratation avec une solution de sels minéraux (soluté de réhydratation orale). D'où viennent ces « sels » ? De quoi sont faits l'eau, le sel, le fer de la marmite ? Toute la matière qui nous entoure est construite à partir de briques invisibles : les atomes. Comment ces briques s'assemblent-elles pour former les molécules et les ions que l'on retrouve dans nos aliments et nos médicaments ?

\courssub{La divisibilité de la matière : une limite infranchissable}

Prenons un morceau de craie. Tu peux le casser en deux, puis en quatre, le broyer en poudre fine. Chaque grain de poudre est encore de la craie. Si tu continues à broyer, tu arrives à des particules que tu ne vois plus à l'œil nu. Mais si tu dissous du sel dans l'eau, le sel semble avoir disparu : la solution est transparente. Pourtant, l'eau a un goût salé. Le sel n'a pas été détruit, il a été divisé jusqu'à des entités trop petites pour être vues : les \cle{atomes} et les \cle{molécules}.

\begin{experience}[Divisibilité de la matière : du visible à l'invisible]
\textbf{Matériel :} morceau de craie, mortier et pilon, sel, eau, bécher, baguette de verre.
\textbf{Protocole :}
\begin{enumerate}
    \item Broie le morceau de craie jusqu'à obtenir une poudre fine. Observe.
    \item Verse de l'eau dans le bécher, ajoute une cuillère de sel, mélange jusqu'à dissolution complète. Observe et goûte (si l'eau est propre).
\end{enumerate}
\textbf{Observations :} La craie se divise en grains de plus en plus petits. Le sel « disparaît » dans l'eau mais le goût persiste.
\textbf{Interprétation :} La matière est divisible. Il existe une limite à cette division : la plus petite particule d'un élément chimique qui conserve les propriétés de cet élément est l'\textbf{atome}. Le sel (chlorure de sodium) est composé d'atomes de sodium et d'atomes de chlore assemblés.
\end{experience}

\begin{definition}[Atome]
L'\textbf{atome} est la plus petite particule d'un élément chimique, électriquement neutre, constituée d'un noyau (protons et neutrons) et d'électrons qui gravitent autour.
\end{definition}

\courssub{La découverte de l'électron : l'expérience de Crookes}

À la fin du XIX\textsuperscript{e} siècle, les physiciens étudient les \textbf{rayons cathodiques} dans des tubes à vide (tubes de Crookes). Ils observent un rayon lumineux partant de la cathode (électrode négative) vers l'anode (électrode positive).

\begin{experience}[Tube de Crookes : mise en évidence de l'électron]
\textbf{Dispositif :} Tube de Crookes (tube de verre sous vide muni de deux électrodes), source de haute tension, aimant.
\textbf{Protocole :}
\begin{enumerate}
    \item Applique une haute tension entre la cathode et l'anode. Un rayon lumineux (rayon cathodique) apparaît, allant de la cathode vers l'anode.
    \item Approche un aimant du tube. Observe la déviation du rayon.
\end{enumerate}
\textbf{Observations :} Le rayon cathodique est dévié par le champ magnétique. La direction de la déviation montre que les particules portent une charge \textbf{négative}.
\textbf{Interprétation (J.J. Thomson, 1897) :} Les rayons cathodiques sont un flux de particules chargées négativement, identiques quel que soit le gaz ou le métal des électrodes. Ces particules sont des constituants universels de la matière : les \textbf{électrons}.
\end{experience}

\begin{methode}[Interpréter une expérience historique]
Pour comprendre un modèle scientifique issu de l'histoire :
\begin{enumerate}
    \item \textbf{Observation :} Que voit-on ? (Déviation du rayon).
    \item \textbf{Hypothèse :} De quoi s'agit-il ? (Particules chargées).
    \item \textbf{Expérience décisive :} Mesure de la rapport charge/masse ($e/m$) par Thomson.
    \item \textbf{Conclusion :} Existence de l'électron, particule légère, charge négative élémentaire.
\end{enumerate}
\end{methode}

\courssub{La découverte du noyau : l'expérience de Rutherford}

Le modèle de Thomson (« pudding aux prunes » : sphère positive avec électrons incrustés) est testé par Ernest Rutherford en 1911. Il bombarde une feuille d'or très fine avec des \textbf{particules alpha} (noyaux d'hélium, charge $+2e$, massives) émises par une source radioactive.

\begin{popfigure}
\begin{tikzpicture}[scale=0.8, >=Stealth]
    % Source
    \draw[fill=poporange!80] (-4,0) rectangle (-3.5,0.5);
    \node at (-3.75, -0.5) {Source $\alpha$};
    \draw[poporange, thick, ->] (-3.5,0.25) -- (-2,0.25);
    % Feuille d'or
    \draw[fill=popgold!60, thick] (-1.5,-1.5) rectangle (-1.3,1.5);
    \node[rotate=90] at (-1.6,0) {Feuille d'or};
    % Trajectoires
    % 1. Traversée (majorité)
    \draw[popblue, thick, ->] (-3.5,1.2) -- (-1.5,1.2) -- (3,1.2);
    \node[popblue] at (3.5,1.2) {Traversée (99,9\%)};
    % 2. Légère déviation
    \draw[popgreen, thick, ->] (-3.5,0.5) -- (-1.5,0.5) .. controls (-1,0.6) and (0,0.8) .. (3,1.0);
    \node[popgreen] at (3.5,1.0) {Faible déviation};
    % 3. Grand déviation
    \draw[poppink, thick, ->] (-3.5,-0.2) -- (-1.5,-0.2) .. controls (-1,0.5) and (0,1.5) .. (1,2.5);
    \node[poppink] at (1.5,3) {Grande déviation};
    % 4. Rebond (retour)
    \draw[popred, thick, ->] (-3.5,-1.0) -- (-1.5,-1.0) .. controls (-1.2,-0.2) and (-2,0.5) .. (-3.5,1.5);
    \node[popred] at (-4,2) {Rebond (rare)};
    % Noyau symbolique au centre de la feuille
    \draw[fill=popdark, opacity=0.5] (-1.4,0) circle (0.15);
    \node[popdark, font=\tiny] at (-1.4, -0.5) {Noyau};
\end{tikzpicture}
\legende{Expérience de Rutherford (1911) : bombardement d'une feuille d'or par des particules $\alpha$. La plupart traversent, quelques-unes sont déviées, de rares rebroussent chemin.}
\end{popfigure}

\begin{experience}[Expérience de Rutherford : interprétation guidée]
\textbf{Résultats expérimentaux :}
\begin{itemize}
    \item La grande majorité des particules $\alpha$ traversent la feuille sans déviation.
    \item Une faible proportion est déviée à de grands angles.
    \item Une infime minorité (environ 1 sur 20 000) rebrousse presque chemin (angle $> 90^\circ$).
\end{itemize}
\textbf{Interprétation (Rutherford) :}
\begin{enumerate}
    \item \textbf{La matière est essentiellement vide :} Les particules traversent sans rencontrer d'obstacle.
    \item \textbf{La charge positive et la masse sont concentrées dans un volume minuscule :} Le \textbf{noyau}. Seule une collision frontale ou rapprochée avec ce noyau massif et positif peut faire rebrousser chemin à une particule $\alpha$ massive et positive.
    \item Les électrons, très légers, gravitent loin du noyau et n'influencent pas la trajectoire des $\alpha$.
\end{enumerate}
\end{experience}

\begin{attention}[Piège : l'atome n'est pas une bille pleine !]
Erreur fréquente : dessiner l'atome comme une sphère compacte colorée.
\textbf{Réalité :} L'atome est \textbf{essentiellement du vide}. Si le noyau était un grain de mil, l'atome ferait la taille d'un stade de football. Les électrons sont des « nuages » de probabilité de présence, pas des billes sur des rails fixes.
\end{attention}

\courssub{Structure de l'atome : noyau et électrons}

L'atome est composé de trois particules subatomiques principales :

\begin{center}
\begin{tabularx}{\linewidth}{|l|X|X|X|}
\hline
\rowcolor{popblueL}
\textbf{Particule} & \textbf{Symbole} & \textbf{Charge électrique} & \textbf{Masse approximative} \\
\hline
Proton & $p^+$ ou $p$ & $+e = +\SI{1,6 \times 10^{-19}}{C}$ & $m_p \approx \SI{1,67 \times 10^{-27}}{kg}$ \\
Neutron & $n^0$ ou $n$ & $0$ (neutre) & $m_n \approx \SI{1,67 \times 10^{-27}}{kg}$ \\
Électron & $e^-$ ou $e$ & $-e = -\SI{1,6 \times 10^{-19}}{C}$ & $m_e \approx \SI{9,11 \times 10^{-31}}{kg}$ \\
\hline
\end{tabularx}
\end{center}

\begin{propriete}[Constitution de l'atome]
\begin{itemize}
    \item Le \textbf{noyau} contient les protons (chargés $+$) et les neutrons (neutres). Il concentre \textbf{presque toute la masse} de l'atome ($m_p \approx 1836 \times m_e$).
    \item Les \textbf{électrons} (chargés $-$) gravitent autour du noyau dans une région appelée \textbf{électrosphère}.
    \item L'atome est \textbf{électriquement neutre} : \textbf{nombre d'électrons = nombre de protons}.
\end{itemize}
\end{propriete}

\begin{aretenir}[Ordres de grandeur à connaître]
\begin{itemize}
    \item Diamètre de l'atome : $\approx \SI{10^{-10}}{m}$ (\SI{0,1}{\nano\meter} = \SI{100}{\pico\meter}).
    \item Diamètre du noyau : $\approx \SI{10^{-15}}{m}$ (\SI{1}{\femto\meter}).
    \item Rapport des diamètres : $\dfrac{d_{\text{atome}}}{d_{\text{noyau}}} \approx 100\,000$.
    \item Rapport des volumes : $\approx 10^{15}$ (le noyau occupe $10^{-13}\%$ du volume).
\end{itemize}
\end{aretenir}

\courssub{Analogie de taille : le stade Ahmadou Ahidjo et le grain de mil}

Pour visualiser ce vide incroyable, imaginons l'atome agrandi à la taille du stade Ahmadou Ahidjo de Yaoundé (environ \SI{100}{m} de diamètre).

\begin{popfigure}
\begin{tikzpicture}[scale=0.6]
    % Stade (cercle grand)
    \draw[popblue, thick] (0,0) circle (6cm);
    \node[popblue, font=\bfseries] at (0, 6.8) {Atome : Stade Ahmadou Ahidjo ($\varnothing \approx \SI{100}{m}$)};
    % Tribunes suggérées
    \foreach \angle in {0,30,...,330} {
        \draw[popblue!50] (\angle:5.5cm) -- (\angle:6cm);
    }
    % Noyau (point minuscule au centre)
    \draw[fill=popred, popred] (0,0) circle (0.03cm); % 0.3mm sur papier = 1.5e-5 * 100m = 1.5mm réel? Non, échelle graphique.
    % Légende noyau
    \draw[popred, <->] (0.03,0) -- (0.5,0);
    \node[popred, right] at (0.5,0) {Noyau : Grain de mil ($\varnothing \approx \SI{1}{mm}$)};
    % Flèche électron
    \draw[popgreen, ->, thick] (0,0) -- (5.5, 2);
    \node[popgreen, above right] at (5.5, 2) {Électrons : Moustiques volant dans les tribunes};
\end{tikzpicture}
\legende{Analogie d'échelle : si l'atome = stade Ahmadou Ahidjo (\SI{100}{m}), le noyau = grain de mil (\SI{1}{mm}) au centre. Les électrons sont comme des moustiques dans les gradins.}
\end{popfigure}

\begin{savaistu}[Le savais-tu ?]
Si l'atome d'hydrogène (le plus simple) était agrandi jusqu'à faire la taille de la Terre ($\varnothing \approx \SI{12\,700}{km}$), son noyau (un seul proton) ferait la taille d'un stade de football (\SI{100}{m}) et l'électron serait une bille de \SI{1}{cm} quelque part dans l'atmosphère ! Cela montre à quel point la matière est « vide » à l'échelle atomique.
\end{savaistu}

\courssub{Masse de l'atome et numéros caractéristiques}

La masse de l'atome est quasi entièrement due aux nucléons (protons + neutrons). La masse de l'électron est négligeable devant celle du nucléon ($m_p \approx 1836 \, m_e$).

\begin{definition}[Numéro atomique $Z$ et nombre de masse $A$]
\begin{itemize}
    \item \textbf{Numéro atomique $Z$} = nombre de protons dans le noyau. Il caractérise l'élément chimique (classement dans le tableau périodique).
    \item \textbf{Nombre de masse $A$} = nombre total de nucléons = nombre de protons ($Z$) + nombre de neutrons ($N$). $A = Z + N$.
    \item Notation d'un noyau : $\ce{^{A}_{Z}X}$ où X est le symbole de l'élément.
\end{itemize}
\end{definition}

\begin{exemplebox}[Exemples : Hydrogène et Carbone]
\begin{itemize}
    \item \textbf{Hydrogène} (le plus abondant dans l'univers, carburant des étoiles, présent dans l'eau $\ce{H2O}$) : $\ce{^{1}_{1}H}$. $Z=1$, $A=1$. \textbf{1 proton, 0 neutron, 1 électron}.
    \item \textbf{Carbone} (base de la vie, du charbon, du graphite, du diamant, du $\ce{CO2}$) : $\ce{^{12}_{6}C}$. $Z=6$, $A=12$. \textbf{6 protons, 6 neutrons, 6 électrons}.
    \item \textbf{Oxygène} (air, eau, combustion) : $\ce{^{16}_{8}O}$. $Z=8$, $A=16$. \textbf{8 protons, 8 neutrons, 8 électrons}.
\end{itemize}
\end{exemplebox}

\begin{exoresolu}[Identifier la composition d'un atome]
\textbf{Énoncé :} Un atome est représenté par le symbole $\ce{^{27}_{13}Al}$. C'est de l'aluminium (marmite, câbles ENEO, canettes). Préciser sa composition en protons, neutrons et électrons. Rappeler pourquoi il est neutre.
\textbf{Solution détaillée :}
\begin{enumerate}
    \item \textbf{Numéro atomique $Z = 13$} $\rightarrow$ \textbf{13 protons} (charge $+13e$).
    \item \textbf{Nombre de masse $A = 27$} $\rightarrow$ Nombre de neutrons $N = A - Z = 27 - 13 = \textbf{14 neutrons}$.
    \item \textbf{Atome neutre} $\rightarrow$ Nombre d'électrons = Nombre de protons = \textbf{13 électrons} (charge $-13e$).
    \item \textbf{Bilan de charge :} $+13e + (-13e) = 0$. L'atome est électriquement neutre.
\end{enumerate}
\end{exoresolu}

\begin{exercice}[Entraînement : composition atomique]
Donner la composition (protons, neutrons, électrons) des atomes suivants, présents dans notre quotidien :
\begin{enumerate}
    \item Fer (marmite, armature béton) : $\ce{^{56}_{26}Fe}$.
    \item Chlore (eau de Javel, sel de cuisine) : $\ce{^{35}_{17}Cl}$.
    \item Azote (air, engrais) : $\ce{^{14}_{7}N}$.
\end{enumerate}
\end{exercice}


\begin{corrige}[Exercice n°1]
\begin{enumerate}
    \item $\ce{Fe}$ : $Z=26 \rightarrow 26\,p^+,\ 26\,e^-$. $N=56-26=30 \rightarrow 30\,n^0$.
    \item $\ce{Cl}$ : $Z=17 \rightarrow 17\,p^+,\ 17\,e^-$. $N=35-17=18 \rightarrow 18\,n^0$.
    \item $\ce{N}$ : $Z=7 \rightarrow 7\,p^+,\ 7\,e^-$. $N=14-7=7 \rightarrow 7\,n^0$.
\end{enumerate}
\end{corrige}

Nous avons maintenant la « brique » de base : l'atome, neutre, avec son noyau massif et ses électrons légers. Mais dans la nature, les atomes s'assemblent rarement seuls. Ils forment des \textbf{molécules} (gaz de la cuisinière, eau, plastique) ou des \textbf{ions} (sel dissous, batterie de moto, crampes musculaires). C'est l'objet de la suite de cette leçon.

\coursec{Symboles des éléments et modèle de Bohr}

\courssub{Du nom de l'élément à son symbole : un langage universel}

Dans la section précédente, nous avons découvert que la matière est constituée d'atomes et que chaque \cle{élément chimique} correspond à une espèce d'atomes bien précise : tous les atomes d'un même élément possèdent le \textbf{même nombre de protons} dans leur noyau (c'est le \cle{numéro atomique}, noté $Z$). Il existe une centaine d'éléments connus, présents dans la nature ou créés en laboratoire. Pour communiquer sans ambiguïté entre chimistes du monde entier — du laboratoire de l'Université de Yaoundé I aux raffineries de la SONARA à Limbé — il faut un code court, unique et standardisé : le \cle{symbole chimique}.

\begin{definition}[Symbole d'un élément chimique]
Le \textbf{symbole d'un élément chimique} est l'abréviation internationale de son nom. Il s'écrit toujours avec une \textbf{majuscule} pour la première lettre, suivie éventuellement d'une \textbf{minuscule} pour la seconde. Il ne prend jamais de point abréviatif.
\end{definition}

\begin{exemplebox}[Exemples de symboles à connaître]
\begin{itemize}
    \item \ce{H} : Hydrogène ; \ce{O} : Oxygène ; \ce{C} : Carbone ; \ce{N} : Azote.
    \item \ce{Fe} : Fer (du latin \emph{ferrum}) ; \ce{Na} : Sodium (du latin \emph{natrium}).
    \item \ce{Cl} : Chlore ; \ce{Ca} : Calcium ; \ce{Cu} : Cuivre (du latin \emph{cuprum}).
    \item \ce{Al} : Aluminium ; \ce{Mg} : Magnésium ; \ce{S} : Soufre ; \ce{K} : Potassium.
\end{itemize}
\end{exemplebox}

\begin{savaistu}[Pourquoi Fe et pas Fr pour le Fer ?]
Beaucoup de symboles proviennent du nom latin de l'élément, héritage de l'alchimie et de la chimie ancienne. C'est le cas pour le Fer (\ce{Fe}, \emph{ferrum}), le Sodium (\ce{Na}, \emph{natrium}), le Potassium (\ce{K}, \emph{kalium}), le Cuivre (\ce{Cu}, \emph{cuprum}), l'Argent (\ce{Ag}, \emph{argentum}), l'Or (\ce{Au}, \emph{aurum}), le Plomb (\ce{Pb}, \emph{plumbum}), l'Étain (\ce{Sn}, \emph{stannum}) et le Tungstène (\ce{W}, \emph{wolfram}). Connaître ces origines évite bien des confusions !
\end{savaistu}

\begin{attention}[Piège fréquent : majuscule ou minuscule ?]
\ce{Co} est le symbole du \textbf{Cobalt}. \ce{CO} est la formule du \textbf{monoxyde de carbone} (un atome de Carbone \ce{C} lié à un atome d'Oxygène \ce{O}), un gaz toxique dégagé par les groupes électrogènes mal ventilés. La casse (majuscule/minuscule) change tout le sens ! De même, \ce{Cl} (Chlore) est un élément, alors que \ce{CL} n'existe pas. Respecte scrupuleusement la règle : \textbf{une majuscule, puis une minuscule si besoin}.
\end{attention}

\courssub{Le modèle de Bohr : visualiser l'atome}

Si le symbole nous dit \emph{quel} atome c'est, le \textbf{modèle de Bohr} (version simplifiée pour la 3ème) nous montre \emph{comment} ses électrons sont organisés autour du noyau. Rappelons-nous : l'atome est globalement \textbf{neutre}. Il possède un noyau (protons chargés positivement + neutrons sans charge) et des électrons chargés négativement qui gravitent autour. Dans le modèle de Bohr, les électrons se répartissent sur des \cle{couches électroniques} circulaires, numérotées de l'intérieur vers l'extérieur : K, L, M, N... (ou 1, 2, 3, 4...).

\begin{propriete}[Règle de remplissage des couches (Programme 3ème)]
Dans le modèle simplifié enseigné en 3ème, les couches se remplissent dans l'ordre, chacune ayant une capacité maximale :
\begin{itemize}
    \item \textbf{1ère couche (K) :} 2 électrons maximum.
    \item \textbf{2ème couche (L) :} 8 électrons maximum.
    \item \textbf{3ème couche (M) :} 8 électrons maximum (pour les éléments jusqu'au Calcium \ce{Ca}, $Z = 20$, étudiés au programme).
\end{itemize}
On note la \textbf{répartition électronique} sous la forme $(n_1, n_2, n_3)$ où $n_i$ est le nombre d'électrons sur la couche $i$.
\end{propriete}

\begin{methode}[Dessiner le modèle de Bohr d'un atome]
Pour représenter un atome de numéro atomique $Z$ (donc $Z$ électrons, car l'atome est neutre) :
\begin{enumerate}
    \item \textbf{Compter les électrons :} $Z$ électrons à placer.
    \item \textbf{Remplir la 1ère couche (K) :} y placer 2 électrons (ou $Z$ si $Z < 2$).
    \item \textbf{Remplir la 2ème couche (L) :} y placer jusqu'à 8 électrons avec le reste.
    \item \textbf{Remplir la 3ème couche (M) :} y placer le reste (jusqu'à 8 pour nos éléments).
    \item \textbf{Tracer le schéma :} dessiner le noyau au centre, puis des cercles concentriques pour les couches utilisées, et placer les électrons (petits points) régulièrement espacés sur chaque cercle.
\end{enumerate}
\end{methode}

\courssub{Exemples tracés : de l'Hydrogène au Sodium}

Appliquons la méthode aux éléments les plus courants de notre environnement : l'Hydrogène (eau, carburant), le Carbone (charbon, matière vivante), l'Oxygène (air, combustion), le Sodium (sel de cuisine, \ce{NaCl}).

\begin{exemplebox}[Répartition électronique des atomes H, C, O, Na]
\begin{center}
\begin{tabularx}{\linewidth}{|c|c|c|X|}\hline
\textbf{Élément} & \textbf{Z} & \textbf{Répartition} & \textbf{Commentaire} \\ \hline
Hydrogène \ce{H} & 1 & (1) & 1 seul électron sur la couche K. \\ \hline
Carbone \ce{C} & 6 & (2, 4) & Couche K pleine (2), 4 électrons sur L. \\ \hline
Oxygène \ce{O} & 8 & (2, 6) & Couche K pleine (2), 6 électrons sur L. \\ \hline
Sodium \ce{Na} & 11 & (2, 8, 1) & K et L pleines ($2 + 8 = 10$), 1 électron sur M. \\ \hline
\end{tabularx}
\end{center}
\end{exemplebox}

\begin{popfigure}
\begin{tikzpicture}[scale=0.85, every node/.style={font=\small}]
    \tikzset{
        noyau/.style={circle, draw=popdark, thick, fill=poporange!80, minimum size=1.1cm, align=center},
        electron/.style={circle, fill=popblue, inner sep=1.6pt},
        couche/.style={draw=popmuted, dashed, thick},
        titre/.style={font=\bfseries\small, text=popdark}
    }

    % --- Hydrogene ---
    \begin{scope}[shift={(0,0)}]
        \node[noyau] (nH) {1p};
        \draw[couche] (nH) circle (1.2cm);
        \node[electron] at (90:1.2) {};
        \node[titre] at (0,-1.8) {Hydrogène \ce{H}};
        \node[font=\scriptsize] at (0,-2.3) {$Z=1$ : (1)};
    \end{scope}

    % --- Carbone ---
    \begin{scope}[shift={(4.5,0)}]
        \node[noyau, align=center] (nC){6p\\6n};
        \draw[couche] (nC) circle (0.8cm);
        \draw[couche] (nC) circle (1.6cm);
        \node[electron] at ($(nC)+(90:0.8)$) {};
        \node[electron] at ($(nC)+(270:0.8)$) {};
        \foreach \a in {45,135,225,315} \node[electron] at ($(nC)+(\a:1.6)$) {};
        \node[titre] at (0,-2.2) {Carbone \ce{C}};
        \node[font=\scriptsize] at (0,-2.7) {$Z=6$ : (2, 4)};
    \end{scope}

    % --- Oxygene ---
    \begin{scope}[shift={(9,0)}]
        \node[noyau, align=center] (nO){8p\\8n};
        \draw[couche] (nO) circle (0.8cm);
        \draw[couche] (nO) circle (1.6cm);
        \node[electron] at ($(nO)+(90:0.8)$) {};
        \node[electron] at ($(nO)+(270:0.8)$) {};
        \foreach \a in {30,90,150,210,270,330} \node[electron] at ($(nO)+(\a:1.6)$) {};
        \node[titre] at (0,-2.2) {Oxygène \ce{O}};
        \node[font=\scriptsize] at (0,-2.7) {$Z=8$ : (2, 6)};
    \end{scope}

    % --- Sodium ---
    \begin{scope}[shift={(13.5,0)}]
        \node[noyau, align=center] (nNa){11p\\12n};
        \draw[couche] (nNa) circle (0.7cm);
        \draw[couche] (nNa) circle (1.3cm);
        \draw[couche] (nNa) circle (1.9cm);
        \node[electron] at ($(nNa)+(90:0.7)$) {};
        \node[electron] at ($(nNa)+(270:0.7)$) {};
        \foreach \a in {0,45,90,135,180,225,270,315} \node[electron] at ($(nNa)+(\a:1.3)$) {};
        \node[electron] at ($(nNa)+(90:1.9)$) {};
        \node[titre] at (0,-2.5) {Sodium \ce{Na}};
        \node[font=\scriptsize] at (0,-3.0) {$Z=11$ : (2, 8, 1)};
    \end{scope}
\end{tikzpicture}
\legende{Modèles de Bohr simplifiés de l'Hydrogène, du Carbone, de l'Oxygène et du Sodium. Le noyau (orange) contient les protons (p) et les neutrons (n). Les cercles pointillés sont les couches K, L, M. Les points bleus sont les électrons.}
\end{popfigure}

\begin{exoresolu}[Déterminer la répartition électronique de l'atome de Fluor]
\textbf{Énoncé :} Le fluor est un gaz très réactif (on retrouve l'élément Fluor dans les fluorures de certains dentifrices). Son numéro atomique est $Z = 9$. Donne la répartition électronique de l'atome de Fluor et décris son modèle de Bohr.
\tcblower
\textbf{Solution détaillée :}
\begin{enumerate}
    \item \textbf{Nombre d'électrons :} l'atome est neutre, donc il possède autant d'électrons que de protons : $9$ électrons.
    \item \textbf{Remplissage des couches :}
    \begin{itemize}
        \item Couche K (maximum 2) : on place 2 électrons. Il reste $9 - 2 = 7$ électrons.
        \item Couche L (maximum 8) : on y place les 7 électrons restants. Il reste 0 électron.
    \end{itemize}
    \item \textbf{Répartition électronique :} \textbf{(2, 7)}.
    \item \textbf{Schéma :} noyau central (9 protons), cercle K avec 2 électrons (placés à l'opposé l'un de l'autre), cercle L avec 7 électrons régulièrement espacés.
\end{enumerate}
\end{exoresolu}

\courssub{Repérer les éléments : extrait du tableau périodique}

Tous ces éléments s'organisent dans le \textbf{tableau périodique}. Les lignes (appelées \cle{périodes}) correspondent au nombre de couches électroniques occupées. Les colonnes (appelées \cle{familles}) regroupent des éléments ayant le même nombre d'électrons sur leur couche externe : ils ont donc des propriétés chimiques voisines. Voici les 20 premiers éléments, essentiels pour la suite de l'année.

\begin{popfigure}
\begin{tikzpicture}[scale=0.72, every node/.style={font=\scriptsize}]
    \colorlet{alkali}{poporange!80}
    \colorlet{alkaline}{popgold!80}
    \colorlet{metalloid}{popgreen!40}
    \colorlet{nonmetal}{poppink!40}
    \colorlet{noble}{poppurple!40}
    \colorlet{othermetal}{popblue!20}

    \tikzset{
        cell/.style={draw=popline, thick, minimum width=1.7cm, minimum height=1.3cm, align=center, inner sep=1pt}
    }

    % Periode 1
    \node[cell, fill=nonmetal, align=center] at (0,0){\textbf{1}\\[1pt]{\large\bfseries \ce{H}}\\[1pt]Hydrogène};
    \node[cell, fill=noble, align=center] at (17,0){\textbf{2}\\[1pt]{\large\bfseries \ce{He}}\\[1pt]Hélium};

    % Periode 2
    \node[cell, fill=alkali, align=center] at (0,-1.5){\textbf{3}\\[1pt]{\large\bfseries \ce{Li}}\\[1pt]Lithium};
    \node[cell, fill=alkaline, align=center] at (1.7,-1.5){\textbf{4}\\[1pt]{\large\bfseries \ce{Be}}\\[1pt]Béryllium};
    \node[cell, fill=metalloid, align=center] at (10.2,-1.5){\textbf{5}\\[1pt]{\large\bfseries \ce{B}}\\[1pt]Bore};
    \node[cell, fill=nonmetal, align=center] at (11.9,-1.5){\textbf{6}\\[1pt]{\large\bfseries \ce{C}}\\[1pt]Carbone};
    \node[cell, fill=nonmetal, align=center] at (13.6,-1.5){\textbf{7}\\[1pt]{\large\bfseries \ce{N}}\\[1pt]Azote};
    \node[cell, fill=nonmetal, align=center] at (15.3,-1.5){\textbf{8}\\[1pt]{\large\bfseries \ce{O}}\\[1pt]Oxygène};
    \node[cell, fill=nonmetal, align=center] at (17,-1.5){\textbf{9}\\[1pt]{\large\bfseries \ce{F}}\\[1pt]Fluor};
    \node[cell, fill=noble, align=center] at (18.7,-1.5){\textbf{10}\\[1pt]{\large\bfseries \ce{Ne}}\\[1pt]Néon};

    % Periode 3
    \node[cell, fill=alkali, align=center] at (0,-3){\textbf{11}\\[1pt]{\large\bfseries \ce{Na}}\\[1pt]Sodium};
    \node[cell, fill=alkaline, align=center] at (1.7,-3){\textbf{12}\\[1pt]{\large\bfseries \ce{Mg}}\\[1pt]Magnésium};
    \node[cell, fill=othermetal, align=center] at (10.2,-3){\textbf{13}\\[1pt]{\large\bfseries \ce{Al}}\\[1pt]Aluminium};
    \node[cell, fill=metalloid, align=center] at (11.9,-3){\textbf{14}\\[1pt]{\large\bfseries \ce{Si}}\\[1pt]Silicium};
    \node[cell, fill=nonmetal, align=center] at (13.6,-3){\textbf{15}\\[1pt]{\large\bfseries \ce{P}}\\[1pt]Phosphore};
    \node[cell, fill=nonmetal, align=center] at (15.3,-3){\textbf{16}\\[1pt]{\large\bfseries \ce{S}}\\[1pt]Soufre};
    \node[cell, fill=nonmetal, align=center] at (17,-3){\textbf{17}\\[1pt]{\large\bfseries \ce{Cl}}\\[1pt]Chlore};
    \node[cell, fill=noble, align=center] at (18.7,-3){\textbf{18}\\[1pt]{\large\bfseries \ce{Ar}}\\[1pt]Argon};

    % Periode 4 (debut)
    \node[cell, fill=alkali, align=center] at (0,-4.5){\textbf{19}\\[1pt]{\large\bfseries \ce{K}}\\[1pt]Potassium};
    \node[cell, fill=alkaline, align=center] at (1.7,-4.5){\textbf{20}\\[1pt]{\large\bfseries \ce{Ca}}\\[1pt]Calcium};

    % Legende
    \node[cell, fill=alkali, minimum width=0.6cm, minimum height=0.4cm] at (0,-6) {};
    \node[anchor=west] at (0.5,-6) {Métaux alcalins (1 électron externe)};
    \node[cell, fill=alkaline, minimum width=0.6cm, minimum height=0.4cm] at (0,-6.7) {};
    \node[anchor=west] at (0.5,-6.7) {Métaux alcalino-terreux (2 électrons externes)};
    \node[cell, fill=nonmetal, minimum width=0.6cm, minimum height=0.4cm] at (0,-7.4) {};
    \node[anchor=west] at (0.5,-7.4) {Non-métaux};
    \node[cell, fill=noble, minimum width=0.6cm, minimum height=0.4cm] at (0,-8.1) {};
    \node[anchor=west] at (0.5,-8.1) {Gaz nobles (couche externe pleine)};
    \node[cell, fill=metalloid, minimum width=0.6cm, minimum height=0.4cm] at (0,-8.8) {};
    \node[anchor=west] at (0.5,-8.8) {Métalloïdes};
    \node[cell, fill=othermetal, minimum width=0.6cm, minimum height=0.4cm] at (0,-9.5) {};
    \node[anchor=west] at (0.5,-9.5) {Autres métaux};
\end{tikzpicture}
\legende{Extrait du tableau périodique (éléments 1 à 20). Dans chaque case : le numéro atomique $Z$ (en haut), le symbole (au centre) et le nom (en bas). Les couleurs indiquent les familles chimiques. Les colonnes des métaux de transition sont omises pour plus de clarté en 3ème.}
\end{popfigure}

\begin{attention}[Lecture du tableau : ce qui est exigible]
Au BEPC, on ne te demande pas de mémoriser les masses des atomes. Tu dois savoir : retrouver le \textbf{symbole} et le \textbf{nom} d'un élément à partir de son numéro atomique $Z$ (et inversement), situer sa \textbf{ligne} (nombre de couches) et sa \textbf{colonne} (nombre d'électrons externes) pour les 20 premiers éléments.
\end{attention}

\courssub{La couche externe : clé de la réactivité chimique}

Regarde bien les schémas de Bohr ci-dessus et le tableau périodique.
\begin{itemize}
    \item L'Hélium \ce{He} (2) et le Néon \ce{Ne} (2, 8) ont leur couche externe \textbf{pleine}. Ce sont des \textbf{gaz nobles}, très peu réactifs (on les utilise dans les tubes lumineux publicitaires ou pour protéger les soudures).
    \item Le Sodium \ce{Na} (2, 8, 1) a \textbf{1 seul électron} sur sa couche externe (M). Il a tendance à le \emph{perdre} pour obtenir la structure stable du Néon.
    \item Le Chlore \ce{Cl} (2, 8, 7) a \textbf{7 électrons} sur sa couche externe (M). Il a tendance à \emph{gagner} 1 électron pour compléter sa couche à 8 et obtenir la structure stable de l'Argon.
\end{itemize}

\begin{aretenir}[Règle de l'octet (savoir admis en 3ème)]
Les atomes cherchent à acquérir la structure électronique stable du gaz noble le plus proche, c'est-à-dire \textbf{8 électrons sur leur couche externe} (2 pour les plus petits atomes comme l'Hydrogène, couche K). C'est ce qui guide la formation des \textbf{molécules} (liaisons entre atomes) et des \textbf{ions} (perte ou gain d'électrons), que nous étudierons dans les sections suivantes.
\end{aretenir}

\begin{attention}[Modèle de Bohr $\neq$ réalité exacte]
Le modèle de Bohr (électrons sur des orbites circulaires fixes, comme des planètes autour du Soleil) est une \textbf{image simplifiée}, commode pour compter les électrons et comprendre la réactivité en 3ème. En réalité, les électrons ne suivent pas des trajectoires précises : on les décrit par des « nuages de probabilité ». Mais pour écrire les formules des molécules et des ions au BEPC, le modèle de Bohr et la règle de remplissage (2, 8, 8) sont \textbf{exactement ce qu'il faut maîtriser}.
\end{attention}

\courssub{Exercice d'application immédiate}

\begin{exercice}[Repérage et représentation]
\begin{enumerate}
    \item Donne le symbole et le nom des éléments de numéro atomique $Z = 4$, $Z = 13$ et $Z = 19$.
    \item Écris la répartition électronique des atomes suivants : Magnésium \ce{Mg} ($Z = 12$), Aluminium \ce{Al} ($Z = 13$), Phosphore \ce{P} ($Z = 15$), Chlore \ce{Cl} ($Z = 17$), Potassium \ce{K} ($Z = 19$).
    \item Parmi les éléments de numéro atomique compris entre 1 et 20, lesquels ont une couche externe complète ? Quelle famille chimique forment-ils ?
    \item Dessine le modèle de Bohr de l'atome de Magnésium \ce{Mg}.
\end{enumerate}
\end{exercice}

\begin{corrige}[Exercice : Repérage et représentation]
\begin{enumerate}
    \item $Z = 4$ : \ce{Be} (Béryllium) ; $Z = 13$ : \ce{Al} (Aluminium) ; $Z = 19$ : \ce{K} (Potassium).
    \item \ce{Mg} ($Z = 12$) : (2, 8, 2) ; \ce{Al} ($Z = 13$) : (2, 8, 3) ; \ce{P} ($Z = 15$) : (2, 8, 5) ; \ce{Cl} ($Z = 17$) : (2, 8, 7) ; \ce{K} ($Z = 19$) : (2, 8, 8, 1).
    \item Les éléments dont la couche externe est complète sont : \ce{He} (2), \ce{Ne} (2, 8) et \ce{Ar} (2, 8, 8). Ils forment la famille des \textbf{gaz nobles} (dernière colonne du tableau).
    \item \textbf{Magnésium \ce{Mg} ($Z = 12$) :} noyau central (12 protons), couche K avec 2 électrons, couche L avec 8 électrons régulièrement espacés, couche M avec 2 électrons placés à l'opposé l'un de l'autre.
\end{enumerate}
\end{corrige}

\courssub{Bilan de la section}

Tu sais maintenant :
\begin{itemize}
    \item Identifier un élément par son \textbf{symbole} (une majuscule, suivie éventuellement d'une minuscule) et le situer dans le \textbf{tableau périodique} (20 premiers éléments).
    \item Utiliser le \textbf{modèle de Bohr simplifié} : noyau central, couches K, L, M...
    \item Appliquer la \textbf{règle de remplissage} (2, 8, 8) pour écrire la \textbf{répartition électronique} $(n_K, n_L, n_M)$.
    \item Dessiner le schéma de Bohr des atomes de H à Ca.
    \item Comprendre que ce sont les \textbf{électrons de la couche externe} qui dictent le comportement chimique : perte, gain ou mise en commun d'électrons pour atteindre la stabilité (règle de l'octet).
\end{itemize}

Dans la prochaine section, nous verrons comment ces atomes s'assemblent pour former des \textbf{molécules} ou se transforment en \textbf{ions} (perte ou gain d'électrons), donnant naissance aux substances qui nous entourent : l'eau que tu bois, le sel que tu manges, le dioxyde de carbone que tu expires, le calcaire de nos roches...

\coursec{Les molécules}

Dans la section précédente, nous avons découvert l'atome : cette brique minuscule, formée d'un noyau et d'électrons, que l'on représente par un symbole et un modèle de Bohr. Mais regarde autour de toi : l'eau que tu bois, l'air que tu respires, le gaz de la bonbonne de cuisine... Aucune de ces matières n'est faite d'atomes isolés ! Les atomes s'associent les uns aux autres pour former des édifices plus grands : les \cle{molécules}. C'est ce que nous allons découvrir maintenant, en commençant par une expérience célèbre qui va nous révéler de quoi l'eau est faite.

\courssub{Une expérience fondatrice : l'électrolyse de l'eau}

\textbf{Question de départ.} L'eau est partout : dans les rivières du Cameroun, dans la marmite, dans notre corps. Mais de quoi est-elle constituée ? Est-ce une matière « purement élémentaire », ou cache-t-elle des ingrédients plus simples ? Pour le savoir, les scientifiques ont eu une idée : \emph{décomposer} l'eau à l'aide du courant électrique et observer ce qui apparaît.

\begin{experience}[L'électrolyse de l'eau au voltamètre]
\textbf{Matériel :} un voltamètre (cuve contenant de l'eau rendue conductrice par quelques gouttes d'acide sulfurique dilué ou de soude), deux électrodes (souvent en platine ou en graphite) reliées à un générateur de courant continu, deux tubes gradués renversés sur les électrodes, une allumette, une braise.

\textbf{Protocole :}
\begin{enumerate}
\item On remplit le voltamètre d'eau acidifiée (ou basique) et on plonge les deux électrodes. On remplit aussi les tubes gradués qu'on retourne délicatement sur chaque électrode.
\item On ferme le circuit : le courant électrique traverse la solution.
\item On observe pendant quelques minutes la montée des gaz dans les tubes, puis on teste les gaz recueillis : on approche une allumette enflammée du tube contenant le plus petit volume de gaz (anode), et une braise (allumette qui se consume sans flamme) du tube contenant le plus grand volume (cathode).
\end{enumerate}

\textbf{Observations :}
\begin{itemize}
\item Des bulles de gaz se forment sur chaque électrode et s'accumulent en haut des tubes gradués.
\item Le volume de gaz recueilli à la \textbf{cathode} (électrode reliée au pôle $-$ du générateur) est \textbf{deux fois plus grand} que celui recueilli à l'\textbf{anode} (électrode reliée au pôle $+$) : le rapport des volumes est de 2 pour 1.
\item Le gaz le plus abondant (cathode) brûle avec un petit bruit sec (un « aboiement ») : c'est le \cle{dihydrogène} \ce{H2}.
\item L'autre gaz (anode) ravive vivement la braise de l'allumette : c'est le \cle{dioxygène} \ce{O2}.
\end{itemize}

\textbf{Interprétation :} l'eau s'est décomposée en deux corps simples, le dihydrogène et le dioxygène, toujours dans le même rapport de volumes $V(\ce{H2}) = 2 \times V(\ce{O2})$. Cela signifie que dans l'eau, les atomes d'hydrogène et d'oxygène sont présents dans le rapport de 2 atomes d'hydrogène pour 1 atome d'oxygène. L'équation de la réaction s'écrit :
\[ \ce{2H2O -> 2H2 + O2} \]

\textbf{Conclusion :} l'eau n'est pas une matière élémentaire. Elle est constituée de particules qui contiennent à la fois de l'\textbf{hydrogène} et de l'\textbf{oxygène}, dans la proportion de 2 atomes d'hydrogène pour 1 atome d'oxygène. Ces particules sont des \textbf{molécules}.
\end{experience}

\begin{attention}[Sécurité lors de l'électrolyse]
Le dihydrogène est un gaz \textbf{très inflammable} : mélangé à l'air, il forme un mélange explosif (le « gaz détonnant ») au contact d'une flamme. Ne jamais approcher une flamme d'un dégagement de dihydrogène sans contrôle strict. Le test de l'aboiement ne se fait \textbf{que} sur le petit volume recueilli dans le tube gradué, \textbf{sous la surveillance de l'enseignant}. On ne manipule jamais seul le générateur ni les produits chimiques ajoutés à l'eau (acide ou soude).
\end{attention}

\begin{popfigure}
\begin{tikzpicture}[scale=0.9, every node/.style={font=\small}]
% Générateur
\draw[fill=popgoldL, thick] (-1.1,-2.6) rectangle (1.1,-1.6);
\node at (0,-2.1) {\textbf{Générateur}};
\node at (-0.6,-1.35) {$-$};
\node at (0.6,-1.35) {$+$};
% Cuve
\draw[thick, popdark] (-3.2,-1.2) -- (-3.2,1.6) -- (3.2,1.6) -- (3.2,-1.2);
% Eau
\fill[popblue!25] (-3.15,-1.15) rectangle (3.15,1.0);
\node[popblue] at (2.3,0.5) {eau acidifiée};
% Électrodes
\draw[thick, fill=popmuted] (-1.7,-1.15) rectangle (-1.5,0.4);
\draw[thick, fill=popmuted] (1.5,-1.15) rectangle (1.7,0.4);
% Fils
\draw[thick] (-1.6,-1.15) -- (-1.6,-1.6) -- (-0.6,-1.6);
\draw[thick] (1.6,-1.15) -- (1.6,-1.6) -- (0.6,-1.6);
% Tubes gradués renversés
\draw[thick] (-2.3,0.2) -- (-2.3,3.0) -- (-0.9,3.0) -- (-0.9,0.2);
\draw[thick] (0.9,0.2) -- (0.9,3.0) -- (2.3,3.0) -- (2.3,0.2);
% Graduations tubes
\foreach \y in {0.5,1.0,1.5,2.0,2.5} {
  \draw[thin] (-2.3,\y) -- (-2.1,\y);
  \draw[thin] (0.9,\y) -- (1.1,\y);
}
% Gaz H2 (grand volume)
\fill[poporangeL] (-2.25,0.9) rectangle (-0.95,2.95);
\node at (-1.6,1.9) {\ce{H2}};
% Gaz O2 (petit volume)
\fill[popgreenL] (0.95,2.0) rectangle (2.25,2.95);
\node at (1.6,2.45) {\ce{O2}};
% Bulles
\foreach \y in {-0.6,-0.1,0.4} {\draw[popblue] (-1.6,\y) circle (0.06);}
\foreach \y in {-0.5,0.0,0.5} {\draw[popblue] (1.6,\y) circle (0.06);}
% Annotations volumes
\draw[<->, thick, poporange] (-2.6,0.9) -- (-2.6,2.95) node[midway, left] {$2V$};
\draw[<->, thick, popgreen] (2.6,2.0) -- (2.6,2.95) node[midway, right] {$V$};
% Labels électrodes
\node[below] at (-1.6,-1.2) {\footnotesize cathode ($-$)};
\node[below] at (1.6,-1.2) {\footnotesize anode ($+$)};
\end{tikzpicture}
\legende{Le voltamètre d'Hofmann : le courant électrique décompose l'eau. Le volume de dihydrogène recueilli à la cathode est le double du volume de dioxygène recueilli à l'anode.}
\end{popfigure}

\courssub{Définition de la molécule}

L'expérience précédente nous a montré que l'eau est faite de particules contenant de l'hydrogène et de l'oxygène \emph{liés ensemble}. Généralisons.

\begin{definition}[Molécule]
Une \cle{molécule} est un édifice \textbf{électriquement neutre} formé d'\textbf{atomes liés entre eux}. Sa \textbf{formule chimique} indique la nature et le nombre des atomes qui la constituent.
\end{definition}

Deux points de cette définition méritent ton attention :

\begin{itemize}
\item \textbf{« Liés entre eux »} : les atomes d'une molécule ne sont pas simplement mélangés, ils sont \emph{accrochés} les uns aux autres par des liaisons chimiques. On ne peut pas les séparer par simple décantation ou filtration : il faut une transformation chimique (comme l'électrolyse pour l'eau).
\item \textbf{« Électriquement neutre »} : une molécule ne porte aucune charge électrique globale. Les charges positives des noyaux compensent exactement les charges négatives des électrons.
\end{itemize}

\begin{aretenir}[L'échelle des molécules]
Une molécule est infiniment petite : de l'ordre de $10^{-10}$ à $10^{-9}$ mètre. Dans une seule goutte d'eau, il y a environ $10^{21}$ molécules d'eau, soit mille milliards de milliards ! C'est pourquoi on ne peut les « voir » qu'à travers des modèles.
\end{aretenir}

\courssub{La formule chimique d'une molécule}

Pour décrire une molécule, les chimistes utilisent une écriture condensée : la \cle{formule chimique}. Elle s'écrit avec les symboles des atomes (vus à la section précédente), chacun suivi, si besoin, d'un \textbf{indice} placé en bas à droite qui indique le nombre d'atomes de ce type.

\begin{exemplebox}[Quelques formules à connaître]
\begin{center}
\begin{tabularx}{\linewidth}{|l|X|}
\hline
\rowcolor{popblueL} \textbf{Formule} & \textbf{Nom et présence dans la vie courante} \\
\hline
\ce{H2O} & Eau (boisson, cuisine, pluie) \\
\hline
\ce{O2} & Dioxygène (air, respiration, combustion) \\
\hline
\ce{CO2} & Dioxyde de carbone (rejeté par la respiration, gaz des boissons gazeuses) \\
\hline
\ce{CH4} & Méthane (biogaz des biodigesteurs, gaz de marais, gaz naturel) \\
\hline
\ce{C6H12O6} & Glucose (sucre utilisé par nos cellules pour l'énergie) \\
\hline
\end{tabularx}
\end{center}
\end{exemplebox}

\begin{methode}[Lire et interpréter une formule chimique]
Pour dénombrer les atomes d'une molécule à partir de sa formule :
\begin{enumerate}
\item Repérer chaque \textbf{symbole} d'atome présent dans la formule (majuscule éventuellement suivie d'une minuscule).
\item Lire l'\textbf{indice} écrit en bas à droite de chaque symbole : c'est le nombre d'atomes de ce type.
\item Si aucun indice n'est écrit, cela signifie qu'il y a \textbf{un seul} atome de ce type : l'indice 1 ne s'écrit jamais.
\item Additionner tous les indices pour obtenir le nombre total d'atomes dans la molécule.
\end{enumerate}
\textbf{Exemple guidé :} \ce{H2O} se lit « 2 atomes d'hydrogène et 1 atome d'oxygène », soit 3 atomes au total.
\end{methode}

\begin{attention}[Le piège de l'indice]
Erreur très fréquente au BEPC : dans \ce{H2O}, le « 2 » ne porte \textbf{que sur l'hydrogène}, pas sur l'oxygène ! La molécule d'eau contient 2 atomes d'hydrogène et \textbf{1 seul} atome d'oxygène. De même, dans \ce{CO2}, c'est l'oxygène qui est en double (indice 2) et le carbone qui est unique (indice 1 non écrit). Un indice ne concerne toujours que le symbole placé \emph{juste avant lui}.
\end{attention}

\begin{exemplebox}[Exemple chiffré : le glucose]
Le glucose, sucre qui donne de l'énergie à nos muscles, a pour formule \ce{C6H12O6}. Comptons ses atomes :
\begin{itemize}
\item carbone C : indice 6 $\rightarrow$ 6 atomes de carbone ;
\item hydrogène H : indice 12 $\rightarrow$ 12 atomes d'hydrogène ;
\item oxygène O : indice 6 $\rightarrow$ 6 atomes d'oxygène.
\end{itemize}
Total : $6 + 12 + 6 = 24$ atomes dans une seule molécule de glucose.
\end{exemplebox}

\courssub{Atome ou molécule : ne pas confondre !}

Voici une distinction capitale. Le symbole \ce{O} désigne \textbf{un atome d'oxygène}, c'est-à-dire l'élément chimique oxygène. Mais l'oxygène que tu respires n'existe pas sous forme d'atomes isolés dans la nature : les atomes d'oxygène s'associent deux par deux pour former la molécule \ce{O2}, le \textbf{dioxygène}.

\begin{aretenir}[La différence à retenir]
\begin{itemize}
\item \ce{O} : un atome d'oxygène (une brique élémentaire).
\item \ce{O2} : une molécule de dioxygène (deux briques accrochées), le gaz de l'air.
\item \ce{O} $\neq$ \ce{O2} : écrire l'un à la place de l'autre est une erreur de chimie majeure !
\end{itemize}
\end{aretenir}

\courssub{Molécules simples et molécules composées}

On classe les molécules en deux familles, selon le nombre de \emph{types} d'atomes qu'elles contiennent.

\begin{definition}[Molécule simple et molécule composée]
\begin{itemize}
\item Une molécule est \cle{simple} si elle ne contient qu'\textbf{un seul type d'atome}. Exemples : \ce{O2} (dioxygène), \ce{H2} (dihydrogène), \ce{N2} (diazote, le gaz le plus abondant de l'air), \ce{Cl2} (dichlore).
\item Une molécule est \cle{composée} si elle contient \textbf{plusieurs types d'atomes}. Exemples : \ce{H2O} (eau), \ce{CO2} (dioxyde de carbone), \ce{CH4} (méthane), \ce{C6H12O6} (glucose).
\end{itemize}
\end{definition}

\begin{attention}[Simple ne veut pas dire « un seul atome »]
La molécule \ce{O2} est dite \emph{simple} alors qu'elle contient \textbf{deux} atomes ! « Simple » signifie : un seul \emph{type} d'atome (ici, uniquement de l'oxygène). Ne confonds pas le \textbf{nombre d'atomes} (indice) avec le \textbf{nombre de types d'atomes} (nombre de symboles différents).
\end{attention}

\courssub{Les modèles moléculaires : voir l'invisible}

Puisqu'aucun œil ne peut voir une molécule, les chimistes les représentent par des \cle{modèles moléculaires} : des boules colorées reliées entre elles, comme des perles sur un fil. La convention de couleurs internationale (CPK) est la suivante :

\begin{center}
\begin{tabularx}{\linewidth}{|l|X|}
\hline
\rowcolor{popblueL} \textbf{Atome} & \textbf{Couleur de la boule} \\
\hline
Hydrogène H & blanc (ou gris clair) \\
\hline
Oxygène O & rouge \\
\hline
Carbone C & noir (ou gris foncé) \\
\hline
Azote N & bleu \\
\hline
Chlore Cl & vert \\
\hline
\end{tabularx}
\end{center}

En classe, tu peux fabriquer ces modèles avec du matériel de fortune : des boules de pâte à modeler de couleurs différentes piquées avec des allumettes ou des cure-dents, ou des haricots de variétés différentes collés ensemble. C'est une excellente façon de mémoriser les formules !

\begin{popfigure}
\begin{tikzpicture}[scale=0.85, every node/.style={font=\small}]
% H2O
\begin{scope}[shift={(0,0)}]
\shade[ball color=red] (0,0) circle (0.55);
\shade[ball color=white] (-0.75,0.55) circle (0.32);
\draw[popdark, thin] (-0.75,0.55) circle (0.32); % contour pour H blanc
\shade[ball color=white] (0.75,0.55) circle (0.32);
\draw[popdark, thin] (0.75,0.55) circle (0.32);
\node[below] at (0,-0.8) {\ce{H2O} : eau};
\end{scope}
% O2
\begin{scope}[shift={(3.4,0.2)}]
\shade[ball color=red] (-0.5,0) circle (0.5);
\shade[ball color=red] (0.5,0) circle (0.5);
\node[below] at (0,-0.9) {\ce{O2} : dioxygène};
\end{scope}
% CO2
\begin{scope}[shift={(6.6,0.2)}]
\shade[ball color=red] (-1.1,0) circle (0.5);
\shade[ball color=black] (0,0) circle (0.55);
\shade[ball color=red] (1.1,0) circle (0.5);
\node[below] at (0,-0.9) {\ce{CO2} : dioxyde de carbone};
\end{scope}
% CH4
\begin{scope}[shift={(10.2,0)}]
\shade[ball color=black] (0,0) circle (0.55);
\shade[ball color=white] (-0.8,0.6) circle (0.3);
\draw[popdark, thin] (-0.8,0.6) circle (0.3);
\shade[ball color=white] (0.8,0.6) circle (0.3);
\draw[popdark, thin] (0.8,0.6) circle (0.3);
\shade[ball color=white] (-0.8,-0.6) circle (0.3);
\draw[popdark, thin] (-0.8,-0.6) circle (0.3);
\shade[ball color=white] (0.8,-0.6) circle (0.3);
\draw[popdark, thin] (0.8,-0.6) circle (0.3);
\node[below] at (0,-1.1) {\ce{CH4} : méthane};
\end{scope}
\end{tikzpicture}
\legende{Modèles moléculaires en boules colorées (convention CPK) : boule blanche = hydrogène (contour gris pour visibilité), boule rouge = oxygène, boule noire = carbone. Compte les boules et compare avec les indices des formules.}
\end{popfigure}

\begin{methode}[Construire un modèle moléculaire à partir d'une formule]
\begin{enumerate}
\item Lire la formule et dénombrer les atomes de chaque type (méthode vue plus haut).
\item Choisir une boule de la bonne couleur pour chaque atome (blanc pour H, rouge pour O, noir pour C, bleu pour N, vert pour Cl).
\item Assembler les boules en les reliant (allumettes, cure-dents) : chaque boule doit être accrochée à au moins une autre pour former un seul édifice.
\item Vérifier : le nombre de boules de chaque couleur doit correspondre exactement aux indices de la formule.
\end{enumerate}
\end{methode}

\begin{savaistu}[Le biogaz au Cameroun]
Dans plusieurs villages et lycées du Cameroun, on installe des \emph{biodigesteurs} : on y enfouit des déchets végétaux et des excréments qui, en fermentant sans air (méthanisation), produisent du \textbf{méthane} \ce{CH4}. Ce gaz sert ensuite à la cuisson, comme le gaz des bonbonnes (butane/propane), ou à faire tourner un groupe électrogène. Une molécule de méthane, c'est simplement 1 atome de carbone lié à 4 atomes d'hydrogène : comprendre les molécules, c'est aussi comprendre les technologies qui transforment notre quotidien et gèrent nos déchets !
\end{savaistu}

\begin{exoresolu}[Lecture de formules]
\textbf{Énoncé.} On donne les formules suivantes : \ce{O2}, \ce{CO2}, \ce{CH4}, \ce{N2}, \ce{C6H12O6}.
\begin{enumerate}
\item Pour chaque molécule, indique le nombre et la nature des atomes.
\item Classe ces molécules en molécules simples et molécules composées.
\end{enumerate}
\tcblower
\textbf{Solution détaillée.}

\textbf{1. Dénombrement des atomes} (on lit chaque symbole et son indice ; l'absence d'indice signifie 1 atome) :
\begin{itemize}
\item \ce{O2} : 2 atomes d'oxygène (2 atomes au total) ;
\item \ce{CO2} : 1 atome de carbone et 2 atomes d'oxygène (3 atomes au total) ;
\item \ce{CH4} : 1 atome de carbone et 4 atomes d'hydrogène (5 atomes au total) ;
\item \ce{N2} : 2 atomes d'azote (2 atomes au total) ;
\item \ce{C6H12O6} : 6 atomes de carbone, 12 atomes d'hydrogène et 6 atomes d'oxygène (24 atomes au total).
\end{itemize}

\textbf{2. Classement :}
\begin{itemize}
\item \textbf{Molécules simples} (un seul type d'atome) : \ce{O2} et \ce{N2}.
\item \textbf{Molécules composées} (plusieurs types d'atomes) : \ce{CO2}, \ce{CH4} et \ce{C6H12O6}.
\end{itemize}
\end{exoresolu}

\begin{exercice}[À toi de jouer]
L'acide acétique, qui donne son goût piquant au vinaigre utilisé en cuisine, a pour formule \ce{C2H4O2}.
\begin{enumerate}
\item Combien d'atomes de chaque type contient une molécule d'acide acétique ? Combien d'atomes au total ?
\item S'agit-il d'une molécule simple ou composée ? Justifie ta réponse.
\item Décris le modèle moléculaire (boules colorées) que tu fabriquerais avec de la pâte à modeler pour représenter cette molécule.
\end{enumerate}
\end{exercice}

\begin{corrige}[Exercice 1]
\begin{enumerate}
\item Carbone : 2 atomes ; hydrogène : 4 atomes ; oxygène : 2 atomes. Total : $2 + 4 + 2 = 8$ atomes.
\item C'est une molécule \textbf{composée}, car elle contient trois types d'atomes différents (C, H et O).
\item Il faut 2 boules noires (carbone), 4 boules blanches (hydrogène) et 2 boules rouges (oxygène), toutes reliées entre elles pour former un seul édifice.
\end{enumerate}
\end{corrige}

\begin{aretenir}[L'essentiel de la section]
\begin{itemize}
\item L'électrolyse de l'eau produit du dihydrogène et du dioxygène dans le rapport de volumes 2 pour 1 : l'eau est donc constituée de particules contenant de l'hydrogène et de l'oxygène.
\item Une \textbf{molécule} est un édifice électriquement neutre formé d'atomes liés entre eux.
\item Sa \textbf{formule} indique la nature et le nombre des atomes ; l'indice ne concerne que le symbole qui le précède, et l'indice 1 ne s'écrit pas.
\item Une molécule est \textbf{simple} (un seul type d'atome : \ce{O2}, \ce{H2}, \ce{N2}) ou \textbf{composée} (plusieurs types : \ce{H2O}, \ce{CO2}, \ce{CH4}).
\item Les modèles moléculaires (boules colorées : H blanc, O rouge, C noir, N bleu, Cl vert) permettent de représenter l'invisible.
\end{itemize}
\end{aretenir}

Mais toutes les particules de la matière ne sont pas électriquement neutres. Que se passe-t-il lorsqu'un atome gagne ou perd un ou plusieurs électrons ? C'est le monde fascinant des \emph{ions}, que nous explorerons dans la section suivante.

\coursec{Les ions : ions simples}

Après avoir découvert comment les atomes s'associent pour former des \cle{molécules} (comme l'eau \ce{H2O} ou le dioxyde de carbone \ce{CO2}), nous allons maintenant explorer une autre façon pour la matière de s'organiser. Pourquoi l'eau du robinet conduit-elle le courant électrique alors que l'eau distillée ne le fait pas ? Pourquoi le sel de cuisine (\ce{NaCl}) disparaît-il dans l'eau en la rendant conductrice ? La réponse se trouve dans l'existence de particules chargées : les \cle{ions}.

\courssub{Découverte expérimentale : la conductivité des solutions}

Commençons par une observation simple que tu peux reproduire à la maison ou au laboratoire. Prenons un circuit électrique basique : une pile, une lampe, un interrupteur et deux fils métalliques terminés par des électrodes (des lames de carbone ou des clous en fer). Nous allons plonger ces électrodes dans deux béchers différents.

\begin{experience}[Test de conductivité de l'eau distillée et de l'eau salée]
\textbf{Matériel :} Pile plate \SI{4,5}{V} ou générateur basse tension, lampe \SI{3,5}{V}, interrupteur, fils de connexion, deux électrodes (graphite ou fer), deux béchers, eau distillée, sel de cuisine (\ce{NaCl}), spatule.

\textbf{Protocole :}
\begin{enumerate}
    \item Monte le circuit en laissant les deux électrodes libres (circuit ouvert). Vérifie que la lampe s'allume quand tu touches les deux électrodes l'une contre l'autre (circuit fermé).
    \item Remplis le premier bécher d'eau distillée. Plonge les électrodes sans qu'elles se touchent. Ferme l'interrupteur. Observe.
    \item Remplis le second bécher d'eau distillée, ajoute une cuillère de sel, remue jusqu'à dissolution complète. Plonge les mêmes électrodes (rincées). Ferme l'interrupteur. Observe.
\end{enumerate}

\textbf{Observations :}
\begin{itemize}
    \item Dans l'eau distillée : la lampe \textbf{ne s'allume pas} (ou brille très faiblement). Le circuit reste ouvert.
    \item Dans l'eau salée : la lampe \textbf{s'allume} franchement. Le circuit est fermé par la solution.
\end{itemize}

\textbf{Interprétation :}
Pour que le courant passe dans un liquide, il faut des \cle{porteurs de charge mobiles}. L'eau distillée en contient très peu (seules quelques molécules d'eau se décomposent spontanément). En ajoutant du chlorure de sodium, nous avons introduit de nouvelles particules qui se déplacent librement dans l'eau et transportent l'électricité. Ces particules chargées sont les \textbf{ions}.
\legende{Le sel invisible à l'œil nu libère des particules chargées qui permettent le passage du courant.}
\end{experience}

\begin{popfigure}
\begin{tikzpicture}[scale=0.9, every node/.style={font=\small}]
    % Styles
    \tikzset{
        comp/.style={draw, thick, minimum width=1.2cm, minimum height=0.8cm, rounded corners=2pt},
        wire/.style={thick},
        electrode/.style={draw, thick, fill=popgray!30, minimum width=0.3cm, minimum height=1.5cm},
        beaker/.style={draw, thick, minimum width=2cm, minimum height=2.5cm, rounded corners=5pt},
        liquid/.style={fill=popblue!20, draw=none, minimum width=1.6cm, minimum height=1.5cm, rounded corners=3pt},
        bulb/.style={draw, thick, circle, minimum size=0.8cm, path picture={\draw[thick] (-0.2,0) -- (0.2,0); \draw[thick] (0,-0.2) -- (0,0.2); \draw[thick] (-0.15,-0.15) -- (0.15,0.15); \draw[thick] (-0.15,0.15) -- (0.15,-0.15);}}
    }

    % Circuit 1 : Eau distillée (Left)
    \begin{scope}[xshift=-5cm]
        \node[comp, label=above:Générateur] (gen1) at (0,2) {$\sim$ \SI{4.5}{V}};
        \node[comp, label=above:Interrupteur] (sw1) at (3,2) {$\backslash$ /};
        \node[bulb, label=above:Lampe] (lamp1) at (6,2) {};
        \node[electrode] (el1a) at (9,2) {};
        \node[electrode] (el1b) at (9,0) {};
        \node[beaker] (beak1) at (9,0.5) {};
        \node[liquid] (liq1) at (9,0.2) {};
        \node[align=center, text width=2.5cm] at (9,-1.5) {\textbf{Eau distillée}\\\textit{Lampe éteinte}};

        \draw[wire] (gen1.east) -- (sw1.west);
        \draw[wire] (sw1.east) -- (lamp1.west);
        \draw[wire] (lamp1.east) -- (el1a.west);
        \draw[wire] (el1b.west) -- ++(-0.5,0) |- (gen1.west);
        \draw[wire, popred] (el1a) -- (el1b); % No current effectively
    \end{scope}

    % Circuit 2 : Eau salée (Right)
    \begin{scope}[xshift=5cm]
        \node[comp, label=above:Générateur] (gen2) at (0,2) {$\sim$ \SI{4.5}{V}};
        \node[comp, label=above:Interrupteur] (sw2) at (3,2) {$\backslash$ /};
        \node[bulb, label=above:Lampe, fill=popgold!50] (lamp2) at (6,2) {};
        \node[electrode] (el2a) at (9,2) {};
        \node[electrode] (el2b) at (9,0) {};
        \node[beaker] (beak2) at (9,0.5) {};
        \node[liquid] (liq2) at (9,0.2) {};
        % Salt ions representation (fixed positions for reproducibility)
        \foreach \angle/\radius in {30/0.5, 90/0.7, 150/0.4, 210/0.6, 270/0.5, 330/0.7} {
            \pgfmathsetmacro{\rx}{9 + \radius*cos(\angle)}
            \pgfmathsetmacro{\ry}{0.5 + \radius*sin(\angle)}
            \node[circle, fill=poporange, minimum size=0.15cm] at (\rx,\ry) {}; % Cations
            \pgfmathsetmacro{\rxb}{\rx+0.1}
            \pgfmathsetmacro{\ryb}{\ry+0.1}
            \node[circle, fill=popgreen, minimum size=0.15cm] at (\rxb,\ryb) {}; % Anions
        }
        \node[align=center, text width=2.5cm] at (9,-1.5) {\textbf{Eau + Sel (\ce{NaCl})}\\\textit{Lampe allumée}};

        \draw[wire] (gen2.east) -- (sw2.west);
        \draw[wire] (sw2.east) -- (lamp2.west);
        \draw[wire] (lamp2.east) -- (el2a.west);
        \draw[wire] (el2b.west) -- ++(-0.5,0) |- (gen2.west);
        \draw[wire, popgreen!70!black, very thick, ->] (el2a) -- (el2b) node[midway, right] {$I$};
    \end{scope}
\end{tikzpicture}
\legende{Schéma du dispositif de test de conductivité. À gauche, l'eau distillée ne conduit pas (lampe éteinte). À droite, la solution de chlorure de sodium conduit le courant (lampe allumée) grâce à la présence d'ions mobiles (\textcolor{poporange}{$\bullet$} cations \ce{Na+}, \textcolor{popgreen}{$\bullet$} anions \ce{Cl-}).}
\end{popfigure}

\begin{propriete}[Électrolyte]
Une solution qui conduit le courant électrique grâce à la présence d'ions mobiles s'appelle une \cle{solution électrolytique} (ou simplement un \cle{électrolyte}). L'eau distillée n'est pas un électrolyte (ou très faiblement) ; l'eau salée, l'eau du robinet, les acides et les bases en sont.
\end{propriete}

\courssub{Définition et nature de l'ion}

L'expérience montre que le chlorure de sodium solide, qui ne conduit pas le courant, libère des particules chargées quand il est dissous dans l'eau. D'où viennent-elles ? Rappelons-nous la structure de l'atome : un noyau positif (protons + neutrons) et un nuage d'électrons négatifs. Un atome est \textbf{électriquement neutre} : nombre de protons = nombre d'électrons.

\begin{definition}[Ion]
Un \cle{ion} est une particule chargée électriquement, formée par un atome (ou un groupe d'atomes) qui a \textbf{gagné} ou \textbf{perdu} un ou plusieurs \textbf{électrons}.
\end{definition}

\textbf{Point fondamental :} Lors de la formation d'un ion, \textbf{le noyau ne change jamais}. Le nombre de protons (numéro atomique $Z$) reste identique. Seule la population électronique varie. L'identité chimique de l'élément (sa place dans le tableau périodique) est donc conservée.

\courssub{Les cations et les anions}

Selon qu'il a perdu ou gagné des électrons, l'ion porte une charge positive ou négative.

\begin{itemize}
    \item \textbf{Cation (ion positif) :} L'atome a \textbf{perdu} un ou plusieurs électrons. Il y a maintenant \textbf{plus de protons que d'électrons}. La charge est positive.
    \item \textbf{Anion (ion négatif) :} L'atome a \textbf{gagné} un ou plusieurs électrons. Il y a maintenant \textbf{plus d'électrons que de protons}. La charge est négative.
\end{itemize}

\begin{exemplebox}[Exemples classiques d'ions simples]
\begin{center}
\begin{tabularx}{\linewidth}{|l|X|X|X|}\hline
\rowcolor{popblueL}
\textbf{Élément (Atome neutre)} & \textbf{Évolution électronique} & \textbf{Ion formé} & \textbf{Type} \\ \hline
Sodium (Na) : 11 p$^+$, 11 e$^-$ & Perd 1 e$^-$ & \ce{Na+} (11 p$^+$, 10 e$^-$) & Cation \\ \hline
Calcium (Ca) : 20 p$^+$, 20 e$^-$ & Perd 2 e$^-$ & \ce{Ca^{2+}} (20 p$^+$, 18 e$^-$) & Cation \\ \hline
Aluminium (Al) : 13 p$^+$, 13 e$^-$ & Perd 3 e$^-$ & \ce{Al^{3+}} (13 p$^+$, 10 e$^-$) & Cation \\ \hline
Fer (Fe) : 26 p$^+$, 26 e$^-$ & Perd 2 e$^-$ (ou 3) & \ce{Fe^{2+}} (ou \ce{Fe^{3+}}) & Cation \\ \hline
Chlore (Cl) : 17 p$^+$, 17 e$^-$ & Gagne 1 e$^-$ & \ce{Cl-} (17 p$^+$, 18 e$^-$) & Anion \\ \hline
Oxygène (O) : 8 p$^+$, 8 e$^-$ & Gagne 2 e$^-$ & \ce{O^{2-}} (8 p$^+$, 10 e$^-$) & Anion \\ \hline
Soufre (S) : 16 p$^+$, 16 e$^-$ & Gagne 2 e$^-$ & \ce{S^{2-}} (16 p$^+$, 18 e$^-$) & Anion \\ \hline
\end{tabularx}
\end{center}
\legende{Les métaux (à gauche du tableau) tendent à perdre des électrons $\rightarrow$ cations. Les non-métaux (à droite) tendent à en gagner $\rightarrow$ anions.}
\end{exemplebox}

\courssub{Modèle de Bohr et formation des ions}

Le modèle de Bohr (couches électroniques K, L, M...) nous aide à visualiser pourquoi tel atome perd ou gagne précisément tel nombre d'électrons. La règle de stabilité (règle du duet pour la couche K, règle de l'octet pour les couches L, M...) guide ce processus : les atomes cherchent à atteindre la configuration électronique du gaz noble le plus proche.

Regardons le sodium (Na, $Z=11$) et le chlore (Cl, $Z=17$), les deux éléments du sel de cuisine.

\begin{popfigure}
\begin{tikzpicture}[scale=0.9, every node/.style={font=\small}]
    % Styles
    \tikzset{
        nucleus/.style={circle, draw, thick, minimum size=1cm, fill=poporange!80, label=center:\textbf{Noyau}},
        shell/.style={circle, draw, thick, dashed},
        electron/.style={circle, fill=popblue, minimum size=0.35cm, draw=white, thick},
        lostelectron/.style={circle, fill=popred, minimum size=0.35cm, draw=white, thick, opacity=0.5},
        arrow/.style={->, thick, popred, >=Stealth}
    }

    % ATOME DE SODIUM Na (2,8,1)
    \begin{scope}[xshift=-4cm]
        \node[nucleus, align=center] (nucNa) at (0,0){Na\\11p$^+$, 12n};
        \node[shell, minimum size=1.8cm] (K) at (0,0) {};
        \node[shell, minimum size=3.2cm] (L) at (0,0) {};
        \node[shell, minimum size=4.6cm] (M) at (0,0) {};
        % Electrons
        \foreach \angle in {90, 210, 330} \node[electron] at (K.\angle) {};
        \foreach \angle in {10, 55, 100, 145, 190, 235, 280, 325} \node[electron] at (L.\angle) {};
        \node[electron] (eM) at (M.90) {};
        \node[align=center] at (0,-3) {\textbf{Atome Na neutre}\\(2, 8, 1)\\11 p$^+$, 11 e$^-$};
    \end{scope}

    % FLÈCHE TRANSFORMATION
    \draw[arrow] (1.5,0) -- node[above] {\textbf{Perte de 1 e$^-$}} (3.5,0);

    % ION SODIUM Na+ (2,8)
    \begin{scope}[xshift=8cm]
        \node[nucleus, align=center] (nucNa2) at (0,0){Na$^+$\\11p$^+$, 12n};
        \node[shell, minimum size=1.8cm] (K2) at (0,0) {};
        \node[shell, minimum size=3.2cm] (L2) at (0,0) {};
        % Couche M disparue
        \foreach \angle in {90, 210, 330} \node[electron] at (K2.\angle) {};
        \foreach \angle in {10, 55, 100, 145, 190, 235, 280, 325} \node[electron] at (L2.\angle) {};
        % Electron perdu symbolique
        \node[lostelectron] at (0, 2.8) {};
        \draw[popred, dashed, thick] (0,2.3) -- (0,3.3) node[above] {e$^-$ perdu};
        \node[align=center] at (0,-3) {\textbf{Ion sodium Na$^+$}\\(2, 8)\\11 p$^+$, 10 e$^-$};
    \end{scope}

    % ATOME DE CHLORE Cl (2,8,7) - Bas gauche
    \begin{scope}[xshift=-4cm, yshift=-5.5cm]
        \node[nucleus, fill=popgreen!80, align=center] (nucCl) at (0,0){Cl\\17p$^+$, 18n};
        \node[shell, minimum size=1.8cm] (K3) at (0,0) {};
        \node[shell, minimum size=3.2cm] (L3) at (0,0) {};
        \node[shell, minimum size=4.6cm] (M3) at (0,0) {};
        \foreach \angle in {90, 210, 330} \node[electron] at (K3.\angle) {};
        \foreach \angle in {10, 55, 100, 145, 190, 235, 280, 325} \node[electron] at (L3.\angle) {};
        \foreach \angle in {15, 60, 105, 150, 195, 240, 285} \node[electron] at (M3.\angle) {}; % 7 electrons
        \node[align=center] at (0,-3) {\textbf{Atome Cl neutre}\\(2, 8, 7)\\17 p$^+$, 17 e$^-$};
    \end{scope}

    % FLÈCHE TRANSFORMATION Cl
    \draw[arrow] (1.5,-5.5) -- node[above] {\textbf{Gain de 1 e$^-$}} (3.5,-5.5);

    % ION CHLORE Cl- (2,8,8) - Bas droit
    \begin{scope}[xshift=8cm, yshift=-5.5cm]
        \node[nucleus, fill=popgreen!80, align=center] (nucCl2) at (0,0){Cl$^-$\\17p$^+$, 18n};
        \node[shell, minimum size=1.8cm] (K4) at (0,0) {};
        \node[shell, minimum size=3.2cm] (L4) at (0,0) {};
        \node[shell, minimum size=4.6cm] (M4) at (0,0) {};
        \foreach \angle in {90, 210, 330} \node[electron] at (K4.\angle) {};
        \foreach \angle in {10, 55, 100, 145, 190, 235, 280, 325} \node[electron] at (L4.\angle) {};
        \foreach \angle in {15, 60, 105, 150, 195, 240, 285, 330} \node[electron] at (M4.\angle) {}; % 8 electrons
        % Electron gagné symbolique
        \node[electron, fill=popgold] at (M4.330) {};
        \draw[popgold!70!black, dashed, thick] (0,2.3) -- (0,3.3) node[above] {e$^-$ gagné};
        \node[align=center] at (0,-3) {\textbf{Ion chlorure Cl$^-$}\\(2, 8, 8)\\17 p$^+$, 18 e$^-$};
    \end{scope}
\end{tikzpicture}
\legende{Formation des ions \ce{Na+} et \ce{Cl-} à partir des atomes neutres. Le sodium (métal) perd son unique électron de la couche M pour atteindre la configuration stable de l'argon (2,8). Le chlore (non-métal) gagne un électron pour compléter sa couche M (octet) et atteindre aussi la configuration de l'argon. Le noyau (protons + neutrons) reste inchangé.}
\end{popfigure}

\begin{methode}[Déterminer la charge d'un ion simple]
Pour trouver la charge électrique d'un ion à partir de sa composition (nombre de protons $Z$ et nombre d'électrons $N_e$) :
\begin{enumerate}
    \item Identifier le nombre de protons $Z$ (numéro atomique, inchangé).
    \item Compter le nombre d'électrons $N_e$ de l'ion.
    \item Comparer : 
    \begin{itemize}
        \item Si $N_e < Z$ $\rightarrow$ il manque des électrons $\rightarrow$ charge \textbf{positive} (Cation). Valeur de la charge $= Z - N_e$.
        \item Si $N_e > Z$ $\rightarrow$ il y a des électrons en trop $\rightarrow$ charge \textbf{négative} (Anion). Valeur de la charge $= N_e - Z$.
    \end{itemize}
    \item Écrire le symbole de l'élément suivi de la charge en exposant à droite (ex: \ce{Ca^{2+}}, \ce{O^{2-}}).
\end{enumerate}
\end{methode}

\courssub{Écriture des formules ioniques : règles de notation}

L'écriture normalisée d'un ion suit des conventions strictes qu'il faut respecter scrupuleusement pour ne pas perdre de points au BEPC.

\begin{aretenir}[Règles d'écriture d'un ion]
\begin{enumerate}
    \item Le \textbf{symbole de l'élément} (une majuscule, éventuellement une minuscule) s'écrit en premier.
    \item La \textbf{charge} s'écrit en \textbf{exposant à droite} du symbole.
    \item Le \textbf{chiffre} (valeur de la charge) se place \textbf{AVANT} le signe $+$ ou $-$.
    \item Si la charge vaut $1$, on n'écrit pas le chiffre $1$ (on écrit \ce{Na+} et non \ce{Na^{1+}} ; \ce{Cl-} et non \ce{Cl^{1-}}).
\end{enumerate}
\end{aretenir}

\begin{attention}[Pièges classiques à éviter absolument]
\begin{itemize}
    \item \textbf{Ordre chiffre/signe :} On écrit \ce{Ca^{2+}} (deux charges positives) et \textbf{jamais} \ce{Ca^{+2}}. De même \ce{O^{2-}} et non \ce{O^{-2}}.
    \item \textbf{Position de la charge :} En exposant \textbf{à droite} (\ce{Na+}), pas à gauche (\ce{^+Na}), pas en indice (\ce{Na_+}).
    \item \textbf{L'atome n'est pas "cassé" :} Un ion \ce{Na+} n'est pas un atome de sodium "incomplet" ou "cassé". C'est un atome de sodium \textbf{entier} (noyau intact : 11 protons) qui a simplement \textbf{perdu un électron}. Sa masse est quasiment identique à celle de l'atome neutre (la masse de l'électron est négligeable).
    \item \textbf{Confusion atome/ion :} Ne confonds pas l'atome de chlore \ce{Cl} (neutre, gazeux, toxique) et l'ion chlorure \ce{Cl-} (en solution dans ton sang, dans l'eau de mer, indispensable à la vie).
\end{itemize}
\end{attention}

\courssub{Exemple résolu : l'ion calcium}

\begin{exoresolu}[Détermination de la formule de l'ion calcium]
\textbf{Énoncé :} L'atome de calcium a pour numéro atomique $Z = 20$. Il perd deux électrons pour devenir stable. Déterminer la composition de l'ion formé (nombre de protons, nombre d'électrons) et écrire sa formule ionique correcte.

\textbf{Solution détaillée :}
\begin{enumerate}
    \item \textbf{Atome neutre Ca :} $Z = 20$ $\Rightarrow$ \textbf{20 protons}. Neutre $\Rightarrow$ \textbf{20 électrons}. Répartition (Bohr) : (2, 8, 8, 2). Les 2 électrons de la couche M sont les électrons de valence.
    \item \textbf{Formation de l'ion :} L'atome \textbf{perd} ses 2 électrons de valence (couche M) pour atteindre la configuration stable de l'argon (2, 8, 8).
    \item \textbf{Composition de l'ion :}
    \begin{itemize}
        \item Protons : \textbf{20} (noyau inchangé).
        \item Électrons : $20 - 2 = \textbf{18}$.
    \end{itemize}
    \item \textbf{Charge électrique :} $20\ \text{p}^+$ et $18\ \text{e}^-$ $\Rightarrow$ 2 charges positives en excès $\Rightarrow$ charge $= +2$.
    \item \textbf{Formule ionique :} Symbole \ce{Ca}, charge $+2$ en exposant à droite $\rightarrow$ \textbf{\ce{Ca^{2+}}}.
\end{enumerate}
\legende{L'ion calcium \ce{Ca^{2+}} a 20 protons et 18 électrons. C'est un cation.}
\end{exoresolu}

\courssub{Vers les composés ioniques}

Dans le sel de cuisine (\ce{NaCl}), les ions sodium \ce{Na+} et chlorure \ce{Cl-} ne sont pas isolés : ils s'attirent par forces électrostatiques (charges opposées) pour former un \textbf{réseau cristallin} solide. Dans l'eau, ce réseau se brise, les ions se séparent et deviennent mobiles : c'est la \textbf{dissolution ionique}. C'est ce que nous avons mis en évidence par le test de conductivité.

Nous verrons dans la prochaine section (\textit{Ions polyatomiques et composés ioniques}) que des groupes d'atomes peuvent aussi former des ions (comme \ce{NH4+}, \ce{SO4^{2-}}, \ce{NO3-}) et comment écrire les formules des solides ioniques (neutralité électrique globale).

\begin{savaistu}[Les ions dans ton quotidien au Cameroun]
\begin{itemize}
    \item \textbf{Batteries et piles :} La pile de ta lampe torche ou la batterie de la moto-taxi fonctionnent grâce au déplacement d'ions (\ce{Li+}, \ce{H+}, \ce{OH-}, \ce{Pb^{2+}}) à l'intérieur d'un électrolyte (pâte, gel ou liquide). Sans ions, pas de courant portable !
    \item \textbf{Eau du robinet (Camwater / ex-SNEC) :} L'eau du robinet contient des ions \ce{Ca^{2+}}, \ce{Mg^{2+}}, \ce{HCO3-}, \ce{Cl-}. C'est pour cela qu'elle conduit le courant (attention danger si tu touches un fil dénudé les pieds dans l'eau !) et qu'elle laisse du calcaire (dépôt de \ce{CaCO3}) dans les bouilloires.
    \item \textbf{Soudure à l'arc / Chalumeau :} Le gaz de protection (souvent \ce{Ar} argon) s'ionise partiellement pour créer un plasma conducteur permettant l'arc électrique.
    \item \textbf{Corrosion (rouille) :} La rouille sur les tôles des cases ou les vieux fûts d'huile usagée, c'est du fer qui s'oxide : \ce{Fe -> Fe^{2+} -> Fe^{3+}} au contact de l'eau et de l'oxygène. Les ions fer migrent et forment l'oxyde de fer(III) \ce{Fe2O3}.
    \item \textbf{Engrais (NPK) :} Les sacs d'engrais (azote \ce{N}, phosphore \ce{P}, potassium \ce{K}) contiennent des sels ioniques (\ce{NH4NO3}, \ce{KCl}, \ce{Ca3(PO4)2}) qui se dissolvent en ions dans l'eau du sol pour être absorbés par les racines du maïs ou du cacao.
\end{itemize}
\end{savaistu}

\coursec{Ions polyatomiques et composés ioniques}

Dans la section précédente, tu as découvert les \cle{ions simples} : un atome qui gagne ou perd des électrons devient un ion chargé, comme \ce{Na+}, \ce{Cl-} ou \ce{Ca^{2+}}. Mais regarde l'étiquette d'une bouteille d'eau minérale Supermont ou Tangui : tu y lis parfois \ce{HCO3-}, \ce{SO4^{2-}}, \ce{NO3-}... Ces ions-là sont formés de \emph{plusieurs} atomes ! Et lorsque les ions s'assemblent, ils forment des solides bien connus : le sel de cuisine, la craie, le calcaire des roches du Cameroun. C'est ce que nous allons explorer maintenant.

\courssub{Les ions polyatomiques}

\textbf{Une observation pour démarrer.}\par

Dans le jardin de la maison, maman utilise parfois de l'\emph{engrais} pour les plantes. Sur le sac, on lit « contient des nitrates et des sulfates ». Or ces substances ne sont pas des atomes seuls : le mot « nitrate » cache un petit groupe soudé d'atomes d'azote et d'oxygène. Comment un groupe d'atomes peut-il porter une charge électrique ?

\begin{definition}[Ion polyatomique]
Un \cle{ion polyatomique} est un \textbf{groupe d'atomes liés entre eux} qui porte une \textbf{charge électrique globale} (positive ou négative). Le groupe se déplace et réagit \emph{comme un seul bloc} : les atomes restent soudés au cours des transformations chimiques.
\end{definition}

L'idée intuitive est simple : imagine une équipe de football. Chaque joueur est un atome ; l'équipe entière, soudée, porte un seul « résultat » collectif : la charge de l'ion. Par exemple, l'ion sulfate \ce{SO4^{2-}} est formé d'un atome de soufre et de quatre atomes d'oxygène, et l'ensemble porte deux charges négatives.

\begin{aretenir}[Les ions polyatomiques à connaître en 3ème]
\begin{center}
\begin{tabularx}{\linewidth}{|l|c|X|}
\hline
\rowcolor{popblueL} \textbf{Nom de l'ion} & \textbf{Formule} & \textbf{Où le rencontre-t-on ?} \\
\hline
Ion hydroxyde & \ce{OH-} & Soude, produits déboucheurs de canalisations \\
\hline
Ion nitrate & \ce{NO3-} & Engrais, eaux minérales \\
\hline
Ion sulfate & \ce{SO4^{2-}} & Plâtre, engrais, eaux minérales \\
\hline
Ion carbonate & \ce{CO3^{2-}} & Calcaire, craie, coquilles d'escargots \\
\hline
Ion ammonium & \ce{NH4+} & Engrais azotés, produits d'entretien \\
\hline
\end{tabularx}
\end{center}
Remarque importante : l'ion \ce{NH4+} (ammonium) est le \textbf{seul cation polyatomique courant} ; tous les autres ions polyatomiques de cette liste sont des anions (chargés négativement).
\end{aretenir}

\begin{attention}[La charge porte sur tout le groupe]
Dans \ce{SO4^{2-}}, ce n'est \emph{pas} le soufre seul qui porte la charge $2-$ : c'est \textbf{l'ensemble du groupe} (1 soufre + 4 oxygènes) qui porte globalement deux charges négatives. De même, dans \ce{OH-}, le groupe oxygène-hydrogène porte ensemble une charge négative. On ne peut pas « casser » un ion polyatomique en cours de réaction : il reste soudé.
\end{attention}

\begin{savaistu}[Les nitrates et l'agriculture camerounaise]
Les engrais azotés utilisés dans les plantations de maïs, de cacao ou de tomates contiennent souvent des ions nitrate \ce{NO3-} et ammonium \ce{NH4+}. Les plantes absorbent ces ions par leurs racines pour fabriquer leurs protéines. Mais un excès d'engrais peut polluer les nappes phréatiques : c'est pourquoi certaines eaux de forage sont contrôlées pour leur teneur en nitrates.
\end{savaistu}

\courssub{Les composés ioniques : des assemblages réguliers et neutres}

\textbf{Situation concrète.}\par

Le sel de cuisine que grand-mère ajoute à la sauce est un solide blanc, formé de petits cristaux cubiques. Pourtant, il est constitué d'ions \ce{Na+} (chargés positivement) et d'ions \ce{Cl-} (chargés négativement). Alors pourquoi le sel, pris dans la main, ne donne-t-il aucune décharge électrique ? Parce que les charges se compensent exactement !

\begin{definition}[Composé ionique]
Un \cle{composé ionique} est un solide constitué d'un \textbf{assemblage régulier et ordonné} de cations et d'anions, qui s'attirent mutuellement. Cet assemblage forme un \textbf{cristal}. Le composé ionique est toujours \textbf{électriquement neutre}.
\end{definition}

Dans un cristal de sel, les ions ne forment pas de molécules isolées : chaque ion \ce{Na+} est entouré d'ions \ce{Cl-}, et chaque ion \ce{Cl-} est entouré d'ions \ce{Na+}, comme des billes rangées en rangées parfaitement régulières, en alternant les charges. C'est cette attraction entre charges opposées qui donne aux cristaux ioniques leur solidité.

\begin{popfigure}
\begin{center}
\begin{tikzpicture}[scale=0.85]
% Grille 4x4 alternée Na+ (petits, bleu) et Cl- (grands, orange)
\foreach \i in {0,1,2,3}{
  \foreach \j in {0,1,2,3}{
    \pgfmathtruncatemacro{\s}{mod(\i+\j,2)}
    \ifnum\s=0
      \shade[ball color=poporange] (\i*1.6,\j*1.6) circle (0.52);
      \node[white] at (\i*1.6,\j*1.6) {\footnotesize \ce{Cl-}};
    \else
      \shade[ball color=popblue] (\i*1.6,\j*1.6) circle (0.34);
      \node[white] at (\i*1.6,\j*1.6) {\scriptsize \ce{Na+}};
    \fi
  }
}
% liaisons horizontales et verticales (dessinées sous les billes : on repasse des traits légers)
\foreach \i in {0,1,2}{
  \foreach \j in {0,1,2,3}{
    \draw[popmuted, thick, opacity=0.5] (\i*1.6+0.55,\j*1.6) -- (\i*1.6+1.05,\j*1.6);
  }
}
\foreach \i in {0,1,2,3}{
  \foreach \j in {0,1,2}{
    \draw[popmuted, thick, opacity=0.5] (\i*1.6,\j*1.6+0.55) -- (\i*1.6,\j*1.6+1.05);
  }
}
% Légende
\shade[ball color=popblue] (5.6,4.4) circle (0.3);
\node[right, popdark] at (6.0,4.4) {ion sodium \ce{Na+}};
\shade[ball color=poporange] (5.6,3.4) circle (0.45);
\node[right, popdark] at (6.15,3.4) {ion chlorure \ce{Cl-}};
\node[right, popdark, text width=4.6cm] at (6.0,2.2) {Alternance régulière : il y a \textbf{autant} de \ce{Na+} que de \ce{Cl-}.};
\end{tikzpicture}
\end{center}
\legende{Modèle du cristal de chlorure de sodium \ce{NaCl} : les ions \ce{Na+} (bleu) et \ce{Cl-} (orange) alternent régulièrement. Chaque ion positif est entouré d'ions négatifs, ce qui assure la cohésion et la neutralité électrique du cristal.}
\end{popfigure}

\begin{propriete}[Électroneutralité d'un composé ionique — savoir admis]
Tout composé ionique est \textbf{électriquement neutre} : la \textbf{somme des charges positives} des cations est égale à la \textbf{somme des charges négatives} des anions. Cette règle d'\cle{électroneutralité} permet d'écrire la formule de n'importe quel composé ionique.
\end{propriete}

Cette propriété n'est pas démontrée au BEPC : elle est admise, mais tu dois savoir l'\emph{utiliser} pour écrire des formules. C'est un savoir-faire exigible.

\courssub{Écrire la formule d'un composé ionique}

\begin{methode}[Écrire la formule d'un composé ionique (« croisement des charges »)]
Pour écrire la formule d'un composé ionique :
\begin{enumerate}
\item Écrire d'abord le \textbf{cation} (ion positif), puis l'\textbf{anion} (ion négatif), sans leurs charges.
\item Repérer la valeur de chaque charge (sans le signe).
\item \textbf{Croiser} : la charge de l'anion devient l'indice du cation, et la charge du cation devient l'indice de l'anion.
\item Simplifier les indices si possible (on ne les écrit jamais s'ils valent 1).
\item Si l'ion polyatomique doit être pris plusieurs fois, le placer entre \textbf{parenthèses} avec l'indice à l'extérieur.
\item Vérifier : charge totale positive = charge totale négative.
\end{enumerate}
\end{methode}

\begin{exemplebox}[Application de la méthode]
\begin{itemize}
\item \textbf{Chlorure de sodium} : \ce{Na+} (charge 1) et \ce{Cl-} (charge 1) $\Rightarrow$ un ion de chaque : \ce{NaCl}. Vérification : $(+1) + (-1) = 0$.
\item \textbf{Chlorure de calcium} : \ce{Ca^{2+}} (charge 2) et \ce{Cl-} (charge 1) $\Rightarrow$ il faut 2 ions chlorure pour 1 ion calcium : \ce{CaCl2}. Vérification : $(+2) + 2\times(-1) = 0$.
\item \textbf{Sulfate de sodium} : \ce{Na+} (charge 1) et \ce{SO4^{2-}} (charge 2) $\Rightarrow$ il faut 2 ions sodium : \ce{Na2SO4}. Vérification : $2\times(+1) + (-2) = 0$.
\item \textbf{Hydroxyde de calcium} : \ce{Ca^{2+}} (charge 2) et \ce{OH-} (charge 1) $\Rightarrow$ il faut 2 ions hydroxyde : \ce{Ca(OH)2}. Vérification : $(+2) + 2\times(-1) = 0$.
\end{itemize}
\end{exemplebox}

\begin{attention}[Les parenthèses sont indispensables !]
Erreur très fréquente au BEPC : écrire \ce{CaOH2} au lieu de \ce{Ca(OH)2}.
\begin{itemize}
\item \ce{Ca(OH)2} signifie : 1 ion calcium et \textbf{2 ions hydroxyde} \ce{OH-} (donc 2 atomes d'oxygène et 2 atomes d'hydrogène).
\item \ce{CaOH2} voudrait dire : 1 calcium, 1 oxygène et 2 hydrogènes — ce qui est \textbf{faux} et ne correspond à aucun composé réel.
\end{itemize}
Règle : dès qu'un \textbf{ion polyatomique} est pris plusieurs fois, on le met entre parenthèses : \ce{Ca(OH)2}, \ce{Mg(NO3)2}, \ce{(NH4)2SO4}. On ne met jamais de parenthèses pour un ion simple : on écrit \ce{CaCl2}, pas \ce{Ca(Cl)2}.
\end{attention}

\begin{exoresolu}[Formules de composés ioniques]
Énoncé : écris la formule des composés ioniques suivants, puis vérifie l'électroneutralité : a) le chlorure de calcium (\ce{Ca^{2+}} et \ce{Cl-}) ; b) le sulfate de sodium (\ce{Na+} et \ce{SO4^{2-}}) ; c) le carbonate de calcium (\ce{Ca^{2+}} et \ce{CO3^{2-}}), constituant de la craie.
\tcblower
Solution détaillée :

a) \textbf{Chlorure de calcium.} Le cation est \ce{Ca^{2+}} (charge 2), l'anion est \ce{Cl-} (charge 1). Pour équilibrer 2 charges positives, il faut 2 ions chlorure :
\[ \ce{Ca^{2+}} + 2\,\ce{Cl-} \longrightarrow \ce{CaCl2} \]
Vérification : $(+2) + 2\times(-1) = 0$. La formule est \ce{CaCl2}.

b) \textbf{Sulfate de sodium.} Cation \ce{Na+} (charge 1), anion \ce{SO4^{2-}} (charge 2). Il faut 2 ions sodium pour compenser les 2 charges négatives du sulfate :
\[ 2\,\ce{Na+} + \ce{SO4^{2-}} \longrightarrow \ce{Na2SO4} \]
Vérification : $2\times(+1) + (-2) = 0$. La formule est \ce{Na2SO4}. Pas de parenthèses : l'ion sulfate n'est pris qu'une seule fois.

c) \textbf{Carbonate de calcium.} Cation \ce{Ca^{2+}} (charge 2), anion \ce{CO3^{2-}} (charge 2). Les charges sont égales : un ion de chaque suffit :
\[ \ce{Ca^{2+}} + \ce{CO3^{2-}} \longrightarrow \ce{CaCO3} \]
Vérification : $(+2) + (-2) = 0$. La formule est \ce{CaCO3} : c'est le calcaire des roches et la craie du tableau.
\end{exoresolu}

\courssub{Lire une étiquette : les ions dans notre vie quotidienne}

\textbf{Démarche d'investigation : exploitation d'un document.}\par

Prends une bouteille d'eau minérale (Supermont, Tangui, ou une autre eau camerounaise) et lis attentivement l'étiquette « composition ». Tu y trouves un tableau d'ions exprimés en milligrammes par litre. Voici un exemple reconstitué :

\begin{popfigure}
\begin{center}
\begin{tikzpicture}
\node[draw=popdark, rounded corners=4pt, fill=popcream, inner sep=10pt, align=center] (etiq){
\begin{minipage}{0.62\linewidth}
\centering
{\large\bfseries\color{popblue} EAU MINÉRALE « SOURCE DU MONT »}\\[2pt]
{\small Analyse moyenne (en \si{mg/L})}\\[4pt]
\begin{tabular}{|l|c||l|c|}
\hline
\rowcolor{popblueL} \multicolumn{2}{|c||}{\textbf{Cations}} & \multicolumn{2}{c|}{\textbf{Anions}} \\
\hline
Calcium \ce{Ca^{2+}} & 78 & Bicarbonates \ce{HCO3-} & 357 \\
\hline
Magnésium \ce{Mg^{2+}} & 24 & Chlorures \ce{Cl-} & 8 \\
\hline
Sodium \ce{Na+} & 6 & Sulfates \ce{SO4^{2-}} & 12 \\
\hline
Potassium \ce{K+} & 1 & Nitrates \ce{NO3-} & 4 \\
\hline
\end{tabular}
\end{minipage}
};
\end{tikzpicture}
\end{center}
\legende{Document à exploiter : étiquette simplifiée d'une eau minérale. Les cations (\ce{Ca^{2+}}, \ce{Mg^{2+}}, \ce{Na+}, \ce{K+}) et les anions (\ce{HCO3-}, \ce{Cl-}, \ce{SO4^{2-}}, \ce{NO3-}) dissous dans l'eau assurent globalement son électroneutralité.}
\end{popfigure}

\begin{experience}[Questions d'exploitation du document]
\begin{enumerate}
\item \textbf{Observer} : repère sur l'étiquette les ions simples et les ions polyatomiques.
\item \textbf{Classer} : range les ions en cations (charges positives) et anions (charges négatives).
\item \textbf{Interpréter} : l'eau minérale est-elle chargée électriquement ? Non : elle contient à la fois des cations et des anions, et la somme des charges positives compense la somme des charges négatives. L'eau reste \textbf{électriquement neutre}.
\item \textbf{Conclure} : les eaux minérales sont des \textbf{solutions ioniques} : elles contiennent des ions dissous, venus de la dissolution des roches (calcaire \ce{CaCO3}, sels divers) traversées par l'eau souterraine.
\end{enumerate}
\end{experience}

Autre exemple de la vie courante : le \cle{soluté de réhydratation orale} (SRO), distribué dans les centres de santé en cas de diarrhée, contient des ions sodium \ce{Na+}, chlorure \ce{Cl-}, potassium \ce{K+} et du glucose. Ces ions remplacent ceux que le corps a perdus : encore une application directe des composés ioniques !

\begin{savaistu}[Pourquoi les eaux minérales ont-elles des goûts différents ?]
C'est la nature et la quantité des ions dissous qui donnent à chaque eau son goût particulier ! Une eau riche en calcium \ce{Ca^{2+}} et en bicarbonates \ce{HCO3-} a traversé des roches calcaires ; une eau très pauvre en ions est dite « faiblement minéralisée ». Les ions calcium et magnésium sont d'ailleurs utiles à la santé des os et des muscles.
\end{savaistu}

\courssub{Test d'identification des ions chlorure}

Comment vérifier qu'une solution contient des ions chlorure \ce{Cl-} ? Les ions sont invisibles à l'œil nu, mais on peut les « trahir » par une réaction chimique qui produit un \textbf{précipité} (un solide qui apparaît dans le liquide).

\begin{experience}[Test du chlorure au nitrate d'argent]
\textbf{Matériel} : deux tubes à essai, une pipette, une solution de nitrate d'argent (\ce{Ag+} + \ce{NO3-}), une solution à tester (par exemple de l'eau salée), gants et lunettes de protection.

\textbf{Protocole} :
\begin{enumerate}
\item Verser environ \SI{2}{mL} de la solution à tester dans un tube à essai.
\item Ajouter quelques gouttes de solution de nitrate d'argent.
\item Observer.
\end{enumerate}

\textbf{Observation} : il se forme immédiatement un \textbf{précipité blanc} (un trouble blanc laiteux) qui noircit à la lumière.

\textbf{Interprétation} : les ions argent \ce{Ag+} apportés par le nitrate d'argent s'associent aux ions chlorure \ce{Cl-} pour former du chlorure d'argent \ce{AgCl}, un solide blanc \emph{insoluble} :
\[ \ce{Ag+} + \ce{Cl-} \longrightarrow \ce{AgCl} \text{ (précipité blanc)} \]

\textbf{Conclusion} : la formation d'un précipité blanc avec le nitrate d'argent prouve la présence d'ions chlorure \ce{Cl-} dans la solution. L'ion nitrate \ce{NO3-}, lui, ne participe pas à la réaction : on dit que c'est un ion « spectateur ».
\end{experience}

\begin{attention}[Consignes de sécurité]
\begin{itemize}
\item Le \textbf{nitrate d'argent tache la peau} (taches noires tenaces) : porte toujours des \textbf{gants} et des \textbf{lunettes} de protection.
\item Ne jamais \textbf{goûter} un produit de laboratoire, même du sel ou de l'eau salée préparés en TP.
\item Ne jamais verser les restes de produits dans l'évier sans l'accord du professeur.
\item Se laver les mains à la fin de la séance.
\end{itemize}
\end{attention}

\begin{exoresolu}[Qui contient des chlorures ?]
Énoncé : un élève prélève deux échantillons : l'eau A (eau de robinet) et l'eau B (eau dans laquelle on a dissous du sel de cuisine). Il ajoute quelques gouttes de nitrate d'argent dans chaque tube. Dans le tube A, rien ne se passe ; dans le tube B, un précipité blanc apparaît. Interprète ces résultats et écris l'équation de la réaction qui a lieu dans le tube B.
\tcblower
Solution détaillée :

\textbf{Tube A} : aucun précipité ne se forme, donc l'eau de robinet ne contient pas (ou très peu) d'ions chlorure \ce{Cl-} : le test est \textbf{négatif}.

\textbf{Tube B} : le précipité blanc montre la présence d'ions chlorure : le test est \textbf{positif}. C'est logique : le sel de cuisine \ce{NaCl} s'est dissous en libérant des ions \ce{Na+} et \ce{Cl-}.

Les ions argent du réactif réagissent avec les ions chlorure selon :
\[ \ce{Ag+} + \ce{Cl-} \longrightarrow \ce{AgCl} \]
Le chlorure d'argent \ce{AgCl} est insoluble : il forme le précipité blanc observé. Les ions \ce{Na+} et \ce{NO3-} restent en solution : ce sont des ions spectateurs.
\end{exoresolu}

\begin{exercice}[À toi de jouer]
\begin{enumerate}
\item Donne la formule des composés ioniques formés à partir des couples d'ions suivants : a) \ce{K+} et \ce{Cl-} ; b) \ce{Mg^{2+}} et \ce{NO3-} ; c) \ce{NH4+} et \ce{SO4^{2-}}.
\item Sur l'étiquette d'une eau minérale, on lit : calcium \ce{Ca^{2+}}, sodium \ce{Na+}, sulfates \ce{SO4^{2-}}, bicarbonates \ce{HCO3-}. Classe ces ions en cations et anions, puis explique pourquoi cette eau est électriquement neutre.
\item Pourquoi écrit-on \ce{Ca(OH)2} avec des parenthèses, mais \ce{CaCl2} sans parenthèses ?
\end{enumerate}
\end{exercice}

\begin{corrige}[Exercice « À toi de jouer »]
\begin{enumerate}
\item a) \ce{K+} (charge 1) et \ce{Cl-} (charge 1) : un ion de chaque $\Rightarrow$ \ce{KCl}. Vérification : $(+1)+(-1)=0$.

b) \ce{Mg^{2+}} (charge 2) et \ce{NO3-} (charge 1) : il faut 2 ions nitrate, ion polyatomique pris deux fois $\Rightarrow$ parenthèses : \ce{Mg(NO3)2}. Vérification : $(+2)+2\times(-1)=0$.

c) \ce{NH4+} (charge 1) et \ce{SO4^{2-}} (charge 2) : il faut 2 ions ammonium $\Rightarrow$ \ce{(NH4)2SO4}. Vérification : $2\times(+1)+(-2)=0$.
\item Cations (positifs) : \ce{Ca^{2+}} et \ce{Na+}. Anions (négatifs) : \ce{SO4^{2-}} et \ce{HCO3-}. L'eau contient à la fois des charges positives et des charges négatives dont les sommes se compensent exactement : elle est donc électriquement neutre, comme tout composé ou mélange ionique.
\item \ce{OH-} est un ion \emph{polyatomique} pris deux fois : les parenthèses dans \ce{Ca(OH)2} montrent que l'indice 2 s'applique à tout le groupe \ce{OH}. \ce{Cl-} est un ion \emph{simple} : l'indice 2 s'applique directement au chlore, donc \ce{CaCl2} s'écrit sans parenthèses.
\end{enumerate}
\end{corrige}

\begin{aretenir}[L'essentiel de la section]
\begin{itemize}
\item Un \cle{ion polyatomique} est un groupe d'atomes liés portant une charge globale : \ce{SO4^{2-}}, \ce{OH-}, \ce{NO3-}, \ce{CO3^{2-}} (anions) et \ce{NH4+} (seul cation courant).
\item Un \cle{composé ionique} est un cristal régulier de cations et d'anions, toujours \textbf{électriquement neutre} : \ce{NaCl} (sel), \ce{CaCO3} (calcaire, craie).
\item Pour écrire une formule : cation d'abord, anion ensuite, puis on équilibre les charges (\emph{croisement des charges}) ; parenthèses obligatoires si un ion polyatomique est pris plusieurs fois : \ce{CaCl2}, \ce{Na2SO4}, \ce{Ca(OH)2}.
\item Les étiquettes d'eaux minérales et les solutés de réhydratation orale montrent que les ions sont partout autour de nous.
\item Test des ions chlorure : le nitrate d'argent donne un \textbf{précipité blanc} de \ce{AgCl} en présence de \ce{Cl-} (avec gants et lunettes !).
\end{itemize}
\end{aretenir}

\coursec{Synthèse : de l'atome aux matériaux du quotidien}

Nous venons de découvrir que certains atomes gagnent ou perdent des électrons pour devenir des \cle{ions}, et que ces ions s'associent en \cle{composés ioniques} comme le chlorure de sodium \ce{NaCl} ou le nitrate d'ammonium \ce{NH4NO3}. Avant cela, nous avions vu que les atomes peuvent aussi s'assembler en \cle{molécules} neutres comme l'eau \ce{H2O} ou le dioxygène \ce{O2}. Il est temps de rassembler toutes ces pièces du puzzle : cette dernière partie te donne une vision complète, de l'atome invisible jusqu'aux matériaux que tu touches chaque jour au Cameroun.

\courssub{Le grand bilan : trois façons pour les atomes d'exister dans la matière}

\begin{aretenir}[Bilan de la leçon]
Toute la matière qui nous entoure est constituée d'\cle{atomes}. Ces atomes se présentent dans la matière sous trois formes principales :
\begin{itemize}
\item \textbf{Seuls (corps simples monoatomiques)} : les gaz rares existent sous forme d'atomes isolés, comme l'argon \ce{Ar} utilisé dans certaines ampoules ou l'hélium \ce{He} des ballons.
\item \textbf{Assemblés en molécules} : les atomes se lient entre eux pour former des \cle{molécules}, particules électriquement neutres : eau \ce{H2O}, dioxygène \ce{O2}, glucose \ce{C6H12O6}, dioxyde de carbone \ce{CO2}.
\item \textbf{Transformés en ions} : en perdant ou en gagnant un ou plusieurs électrons, les atomes deviennent des \cle{ions} chargés. Les ions positifs (cations) et négatifs (anions) s'attirent et forment des \cle{composés ioniques} électriquement neutres : sel \ce{NaCl}, chaux \ce{Ca(OH)2}, engrais \ce{NH4NO3}.
\end{itemize}
\end{aretenir}

Retiens bien les deux « chemins » possibles à partir de l'atome neutre :

\begin{itemize}
\item \textbf{Chemin 1 (liaison covalente) :} atome $\rightarrow$ \textbf{molécule} (les atomes mettent des électrons en commun pour se lier ; l'ensemble reste neutre).
\item \textbf{Chemin 2 (liaison ionique) :} atome $\rightarrow$ \textbf{ion} (perte ou gain d'électrons) $\rightarrow$ \textbf{composé ionique} (attraction électrostatique entre ions de charges opposées ; l'ensemble est neutre).
\end{itemize}

\begin{popfigure}
\begin{center}
\begin{tikzpicture}[mindmap, grow cyclic, every node/.style={concept}, concept color=popblue!60,
root concept/.append style={font=\large\bfseries, minimum size=2.8cm},
level 1/.append style={level distance=3.6cm, sibling angle=120, font=\small\bfseries},
level 2/.append style={level distance=2.8cm, sibling angle=50, font=\scriptsize}]
\node[root concept]{MATIÈRE}
  child[concept color=poporange!70]{ node {Atomes}
    child[concept color=poporange!40]{ node {Seuls (gaz rares) : \ce{Ar}, \ce{He}} }
    child[concept color=popgreen!60]{ node {Molécules (neutres)}
      child[concept color=popgreen!35]{ node {\ce{H2O}, \ce{O2}, \ce{CO2}} }
      child[concept color=popgreen!35]{ node {Glucose \ce{C6H12O6}} }
    }
    child[concept color=poppurple!60]{ node {Ions (chargés)}
      child[concept color=poppurple!35]{ node {Cations : \ce{Na+}, \ce{Ca^{2+}}, \ce{NH4+}} }
      child[concept color=poppurple!35]{ node {Anions : \ce{Cl-}, \ce{NO3-}, \ce{SO4^{2-}}} }
    }
  };
\end{tikzpicture}
\end{center}
\legende{Carte mentale de synthèse : la matière est faite d'atomes, qui existent seuls (gaz rares), assemblés en molécules neutres, ou transformés en ions chargés formant des composés ioniques.}
\end{popfigure}

\courssub{Tableau récapitulatif : atome, molécule, ion}

Pour bien préparer le BEPC, tu dois savoir distinguer sans hésiter ces trois types de particules. Le tableau suivant résume l'essentiel.

\begin{center}
\begin{tabularx}{\linewidth}{|l|X|X|X|}
\hline
\rowcolor{popblueL} \textbf{Critère} & \textbf{Atome} & \textbf{Molécule} & \textbf{Ion} \\
\hline
Définition & Plus petite particule d'un élément chimique conservant ses propriétés (noyau + électrons) & Assemblage d'au moins deux atomes liés entre eux (liaison covalente) & Atome ou groupe d'atomes ayant perdu ou gagné des électrons \\
\hline
Charge électrique & Nulle (neutre) & Nulle (neutre) & Positive (cation) ou négative (anion) \\
\hline
Formule type & Symbole unique : \ce{Fe}, \ce{O}, \ce{Na} & Plusieurs symboles, pas de charge : \ce{H2O}, \ce{O2}, \ce{CH4} & Charge en exposant : \ce{Na+}, \ce{Cl-}, \ce{NO3-}, \ce{Ca^{2+}} \\
\hline
Exemple vie courante & Fer du clou (réseau d'atomes \ce{Fe}), cuivre des fils électriques & Eau de boisson (\ce{H2O}), air respiré (\ce{O2}, \ce{N2}) & Sel de cuisine (\ce{Na+}, \ce{Cl-}), ions de l'eau minérale (\ce{Ca^{2+}}, \ce{HCO3-}) \\
\hline
\end{tabularx}
\end{center}

\begin{methode}[Identifier le constituant d'une espèce chimique en 3 questions]
Face à une formule chimique au BEPC ou en TP, pose-toi trois questions dans l'ordre :
\begin{enumerate}
\item \textbf{La formule comporte-t-elle une charge ($+$, $-$, $2+$, \dots) en exposant ?} Si oui, c'est un \cle{ion}. Exemple : \ce{SO4^{2-}} est un ion polyatomique (anion).
\item \textbf{Sinon, la formule contient-elle plusieurs symboles d'éléments différents (ou un indice $>1$) ?} Si oui, c'est une \cle{molécule} (corps composé covalent) \textbf{ou} un \cle{composé ionique} écrit sous forme brute (ex. \ce{NaCl}, \ce{CaO}). \emph{Astuce : si un métal et un non-métal sont associés sans charge, c'est un composé ionique.}
\item \textbf{Sinon, n'y a-t-il qu'un seul symbole (éventuellement avec un indice) ?} Alors c'est un \cle{atome} (corps simple monoatomique comme \ce{Fe}, \ce{Ar} ou diatomique comme \ce{O2} — mais \ce{O2} est une molécule !).
\end{enumerate}
\textbf{Vérification finale :} Une molécule et un atome sont neutres ; un ion porte une charge.
\end{methode}

\courssub{De l'atome aux matériaux du quotidien au Cameroun}

La chimie n'est pas une science abstraite : chaque objet de ta journée illustre cette leçon. Observons quelques exemples familiers.

\begin{exemplebox}[Le sel de cuisine : formation d'un composé ionique]
Le sel que maman met dans la sauce est le chlorure de sodium \ce{NaCl}. Il est formé d'ions sodium \ce{Na+} (un atome de sodium qui a perdu un électron) et d'ions chlorure \ce{Cl-} (un atome de chlore qui a gagné un électron).

\begin{center}
\begin{tikzpicture}[scale=0.9, every node/.style={font=\small}]
% Atome Na
\node[draw, circle, minimum size=1.2cm, fill=poporange!20] (Na) at (0,0) {\ce{Na} (11p, 11e-)};
\node[below] at (0,-0.8) {Atome Na neutre};
% Flèche perte e-
\draw[->, thick, poporange] (1.2,0.2) -- (2.5,0.2) node[midway, above] {\cle{- e-}};
% Ion Na+
\node[draw, circle, minimum size=1.0cm, fill=poporange!40] (Na+) at (3.5,0) {\ce{Na+} (11p, 10e-)};
\node[below] at (3.5,-0.7) {Cation \ce{Na+}};

% Atome Cl
\node[draw, circle, minimum size=1.2cm, fill=popgreen!20] (Cl) at (0,-3) {\ce{Cl} (17p, 17e-)};
\node[below] at (0,-3.8) {Atome Cl neutre};
% Flèche gain e-
\draw[->, thick, popgreen] (1.2,-2.8) -- (2.5,-2.8) node[midway, above] {\cle{+ e-}};
% Ion Cl-
\node[draw, circle, minimum size=1.0cm, fill=popgreen!40] (Cl-) at (3.5,-3) {\ce{Cl-} (17p, 18e-)};
\node[below] at (3.5,-3.7) {Anion \ce{Cl-}};

% Réseau
\node[align=center, text width=5cm] at (7.5,-1.5) {
\textbf{Cristal de \ce{NaCl}} :\\
Les ions \ce{Na+} et \ce{Cl-} s'empilent\\
régulièrement par attraction\\
électrostatique.\\
\emph{Neutralité globale :} $N(\ce{Na+}) = N(\ce{Cl-})$.
};
\draw[->, thick, poppurple] (Na+) -- (6.5,-0.5);
\draw[->, thick, poppurple] (Cl-) -- (6.5,-2.5);
\end{tikzpicture}
\end{center}
\legende{Formation des ions \ce{Na+} et \ce{Cl-} et constitution du cristal de sel (schéma simplifié).}
\end{exemplebox}

\begin{exemplebox}[L'eau : une molécule covalente]
L'eau que tu bois est faite de molécules \ce{H2O} : deux atomes d'hydrogène liés à un atome d'oxygène par partage d'électrons. La molécule d'eau est \textbf{électriquement neutre}. Quand on fait bouillir l'eau, ce sont les molécules \ce{H2O} entières qui passent à l'état gazeux (vapeur d'eau) ; les liaisons entre H et O ne se cassent pas.
\end{exemplebox}

\begin{exemplebox}[Les engrais NPK : un cocktail d'ions pour l'agriculture]
Les engrais NPK que les planteurs utilisent dans les champs de maïs, de cacao ou de coton au Cameroun apportent aux plantes des ions indispensables à leur croissance :
\begin{itemize}
\item l'ion nitrate \ce{NO3-} et l'ion ammonium \ce{NH4+} pour l'azote (N) ;
\item l'ion potassium \ce{K+} pour le potassium (K) ;
\item les ions phosphates \ce{H2PO4-} / \ce{HPO4^{2-}} pour le phosphore (P).
\end{itemize}
Dissous dans l'eau du sol, ces ions sont absorbés par les racines : la plante se nourrit littéralement d'ions minéraux !
\end{exemplebox}

\begin{exemplebox}[La chaux : traitement de l'eau et construction]
L'hydroxyde de calcium \ce{Ca(OH)2}, appelé \emph{chaux éteinte}, est un composé ionique formé d'ions calcium \ce{Ca^{2+}} et d'ions hydroxyde \ce{OH-}. Au Cameroun, on l'utilise pour assainir les latrines (élévation du pH), blanchir les murs (badigeon) et traiter l'eau de boisson (floculation des impuretés). C'est encore la chimie des ions au service de la santé et de l'habitat.
\end{exemplebox}

\begin{savaistu}[Le soluté de réhydratation orale (SRO) : des ions qui sauvent des vies]
Lors d'une diarrhée aiguë ou d'un choléra, le corps perd énormément d'eau et de sels minéraux. Le \cle{soluté de réhydratation orale} (SRO ou ORS), distribué gratuitement dans les centres de santé du Cameroun, sauve chaque année des milliers d'enfants. Pourtant, sa composition est d'une simplicité chimique remarquable : un mélange bien dosé d'ions sodium \ce{Na+}, potassium \ce{K+} et chlorure \ce{Cl-} (pour remplacer les pertes), plus du glucose (molécule \ce{C6H12O6}) qui active le transport couplé sodium-glucose dans l'intestin, permettant l'absorption massive d'eau. Comprendre les ions et les molécules, c'est comprendre un des plus grands succès de la médecine de terrain !
\end{savaistu}

\begin{attention}[Ne confonds pas molécule d'eau et ions dissous dans l'eau]
C'est l'erreur la plus fréquente au BEPC !
\begin{itemize}
\item La \textbf{molécule d'eau} \ce{H2O} est \textbf{neutre} : l'eau \textbf{chemiquement pure} (distillée) ne conduit presque pas le courant électrique.
\item L'\textbf{eau du robinet}, de \textbf{source} ou \textbf{minérale} contient en plus des \textbf{ions dissous} : \ce{Ca^{2+}}, \ce{Mg^{2+}}, \ce{Na+}, \ce{HCO3-}, \ce{Cl-}\dots Ce sont \textbf{ces ions libres} qui font que l'eau courante conduit le courant électrique.
\end{itemize}
Dire « l'eau est un ion » ou « \ce{H2O} est chargé » est \textbf{faux}. L'eau est une molécule neutre ; les ions sont des espèces \emph{dissoutes} dans l'eau (solvatées).
\end{attention}

\courssub{Des particules incroyablement petites et nombreuses}

\begin{exoresolu}[Combien d'ions dans un gramme de sel ?]
\textbf{Énoncé.} On pèse \SI{1,0}{g} de sel de cuisine \ce{NaCl}. Sachant qu'un ion sodium a une masse d'environ $3,8 \times 10^{-23}$ g et qu'un ion chlorure a une masse d'environ $5,9 \times 10^{-23}$ g, estime le nombre d'ions \ce{Na+} présents dans ce gramme de sel. On rappelle qu'il y a autant d'ions \ce{Na+} que d'ions \ce{Cl-}.
\tcblower
\textbf{Solution détaillée.}
\begin{enumerate}
\item \textbf{Masse d'un « couple » d'ions \ce{Na+} + \ce{Cl-} (une formule de \ce{NaCl}) :}
\[ m_{\text{couple}} = 3,8 \times 10^{-23} + 5,9 \times 10^{-23} = 9,7 \times 10^{-23} \text{ g} \approx 1,0 \times 10^{-22} \text{ g}. \]
\item \textbf{Nombre de couples d'ions (formules \ce{NaCl}) dans \SI{1,0}{g} :}
\[ N = \frac{m_{\text{total}}}{m_{\text{couple}}} = \frac{1,0}{1,0 \times 10^{-22}} = 1,0 \times 10^{22} \text{ couples}. \]
\item \textbf{Conclusion :} Dans \SI{1,0}{g} de sel, il y a environ $1,0 \times 10^{22}$ ions \ce{Na+} et autant d'ions \ce{Cl-}, soit environ $2,0 \times 10^{22}$ ions au total.
\end{enumerate}
Ce nombre est immense : c'est un 1 suivi de 22 zéros ! Cela montre à quel point les atomes et les ions sont minuscules.
\end{exoresolu}

\begin{savaistu}[Pour aller plus loin : la mole (classe de Première)]
Comment les chimistes comptent-ils ces milliards de milliards de particules sans s'épuiser ? Ils utilisent une unité de « paquet » appelée la \cle{mole} (symbole \cle{mol}), et une grandeur appelée \cle{masse molaire} ($M$, en g/mol). Tu les étudieras en classe de Première. Pour l'instant, retiens seulement ceci : un atome a une masse extraordinairement petite, de l'ordre de $10^{-27}$ kg ($10^{-24}$ g). C'est pourquoi même un grain de sel (\SI{1}{mg}) contient encore environ $10^{19}$ ions !
\end{savaistu}

\courssub{Consignes de sécurité générales en travaux pratiques de chimie}

Toute cette belle chimie se manipule au laboratoire, mais avec des règles strictes. Un bon chimiste est d'abord un chimiste prudent.

\begin{aretenir}[Les règles d'or du TP de chimie]
\begin{enumerate}
\item \textbf{Protège-toi} : porte une blouse en coton (fermée) et des lunettes de protection ; attache les cheveux longs ; pas de sandales (pieds couverts).
\item \textbf{Ne goûte jamais} un produit, même le sel ou le sucre du laboratoire (risque de contamination) ; ne respire pas directement les vapeurs (technique du « éventail » avec la main si le professeur l'autorise).
\item \textbf{Lis toujours l'étiquette} du flacon avant de prélever un produit : nom, formule, pictogrammes de danger (inflammable, corrosif, toxique, danger pour l'environnement).
\item \textbf{Ne verse jamais} un produit non utilisé dans son flacon d'origine (risque de pollution du stock) ; jette les déchets chimiques dans les récipients prévus à cet effet (bac « déchets solides », « déchets liquides organiques », « déchets acides/bases »), \textbf{jamais dans l'évier}.
\item \textbf{Électricité} : lors des montages (électrolyse, conductivité), vérifie le circuit \textbf{avant} de fermer l'interrupteur, ne touche pas les fils dénudés sous tension, signale tout fil qui chauffe anormalement.
\item En cas d'accident (projection dans l'œil ou sur la peau, brûlure, coupure, ingestion), \textbf{préviens immédiatement le professeur} et suis la procédure (douche oculaire, rinçage abondant à l'eau pendant 15 min minimum).
\end{enumerate}
\end{aretenir}

\begin{exercice}[À toi de jouer !]
Pour chaque espèce chimique suivante, indique s'il s'agit d'un atome, d'une molécule, d'un ion simple ou d'un ion polyatomique, puis cite un exemple de la vie courante (au Cameroun si possible) où on la rencontre : \ce{O2} ; \ce{Ca^{2+}} ; \ce{Fe} ; \ce{NO3-} ; \ce{CO2} ; \ce{Na+} ; \ce{Ar}.
\end{exercice}

\begin{corrige}[Exercice : À toi de jouer !]
\begin{itemize}
\item \ce{O2} : \textbf{molécule} (deux atomes d'oxygène liés, neutre) — le dioxygène de l'air que nous respirons, indispensable à la combustion (moto-taxi, groupe électrogène).
\item \ce{Ca^{2+}} : \textbf{ion simple} (un seul atome, chargé positivement, cation) — présent dans l'eau minérale (Volcanic, Tangui), le lait, la chaux \ce{Ca(OH)2}.
\item \ce{Fe} : \textbf{atome} (un seul symbole, neutre) — le fer des clous, tôles, tiges de construction (béton armé), carrosseries de voiture.
\item \ce{NO3-} : \textbf{ion polyatomique} (plusieurs atomes N et O, chargé, anion) — l'ion nitrate des engrais NPK (SONARA/usines locales), des explosifs miniers.
\item \ce{CO2} : \textbf{molécule} (neutre) — le dioxyde de carbone rejeté par la respiration, les moteurs thermiques (moto, voiture), la fermentation (bière, vin de palme).
\item \ce{Na+} : \textbf{ion simple} (cation) — l'ion sodium du sel de cuisine \ce{NaCl}, du soluté de réhydratation orale (SRO), du glutamate (cubes Maggi/Knorr).
\item \ce{Ar} : \textbf{atome} (gaz rare, monoatomique) — l'argon des ampoules à incandescence, du soudage TIG (protection du bain de fusion).
\end{itemize}
\end{corrige}

\begin{aretenir}[L'essentiel de la leçon en cinq phrases]
\begin{enumerate}
\item La matière est constituée d'\cle{atomes} (noyau + électrons), électriquement neutres, de masse $\sim 10^{-27}$ kg.
\item Les atomes s'assemblent en \cle{molécules} neutres (\ce{H2O}, \ce{O2}, \ce{CO2}, \ce{C6H12O6}) par partage d'électrons.
\item En perdant ou gagnant des électrons, les atomes deviennent des \cle{ions} chargés : cations ($+$) et anions ($-$).
\item Les ions de charges opposées forment des \cle{composés ioniques} neutres (\ce{NaCl}, \ce{Ca(OH)2}, \ce{NH4NO3}), présents partout : sel, engrais, chaux, soluté ORS, médicaments.
\item Ces particules sont infiniment petites et donc infiniment nombreuses : \SI{1}{g} de sel contient environ $10^{22}$ ions de chaque sorte.
\end{enumerate}
\end{aretenir}

Tu as maintenant une vision complète des constituants de la matière : de l'atome invisible au grain de sel sur ta table, en passant par l'eau que tu bois et l'air que tu respires, tout est chimie ! Garde cette carte mentale en tête, elle te servira dans toutes les leçons de chimie de l'année — et au BEPC.

\newpage

\coursec{Activité d'intégration}

\courssub{Objectifs de l'activité}

À la fin de cette activité d'intégration, tu seras capable de :

\begin{itemize}
\item classer les espèces chimiques présentes dans une eau minérale en \cle{cations} et \cle{anions}, en justifiant à l'aide du \cle{modèle de Bohr} ;
\item écrire la \cle{formule} d'un \cle{composé ionique} en appliquant la règle d'\cle{électroneutralité} ;
\item vérifier quantitativement qu'une eau minérale est \cle{électriquement neutre} ;
\item proposer un \cle{protocole expérimental} (test de conductivité) prouvant la présence d'ions, avec les \cle{consignes de sécurité} ;
\item expliquer le rôle des ions dans un \cle{soluté de réhydratation orale} (SRO) ;
\item restituer tes résultats sous forme d'une \cle{affiche} claire et rigoureuse.
\end{itemize}

\courssub{Situation}

\textbf{Enquête sur l'eau minérale du quartier.}

Pendant les grandes vacances, le club scientifique du collège de Nkolndongo (Yaoundé) décide d'étudier les eaux de boisson vendues dans le quartier. Les élèves remarquent que certaines personnes préfèrent l'eau minérale en bouteille à l'eau du robinet, car « elle contient des sels minéraux bons pour la santé ». Mais que signifie exactement cette phrase ?

Sur l'étiquette d'une bouteille d'eau minérale camerounaise de \SI{1,5}{L}, les élèves relèvent la composition suivante (teneurs en ions, en milligrammes par litre) :

\begin{center}
\begin{tabularx}{\linewidth}{|l|X|X|}\hline
\rowcolor{popblueL} \textbf{Ion} & \textbf{Formule} & \textbf{Teneur (mg/L)} \\ \hline
Calcium & \ce{Ca^{2+}} & 40 \\ \hline
Magnésium & \ce{Mg^{2+}} & 24 \\ \hline
Sodium & \ce{Na^{+}} & 23 \\ \hline
Potassium & \ce{K^{+}} & 4 \\ \hline
Chlorure & \ce{Cl^{-}} & 35,5 \\ \hline
Sulfate & \ce{SO4^{2-}} & 48 \\ \hline
Hydrogénocarbonate & \ce{HCO3^{-}} & 183 \\ \hline
\end{tabularx}
\end{center}

Le professeur de PCT pose alors le problème suivant : « Cette eau est-elle bien \cle{électriquement neutre} ? Peut-on prouver expérimentalement qu'elle contient des ions ? Et pourquoi le soluté de réhydratation orale, distribué gratuitement dans les centres de santé contre la diarrhée, contient-il justement des ions \ce{Na^{+}}, \ce{K^{+}} et \ce{Cl^{-}} ? »

Par groupes de quatre, vous êtes chargés de mener l'enquête et de présenter vos conclusions sur une \cle{affiche}.

\textbf{Données utiles :} numéros atomiques : $Z(\mathrm{Ca}) = 20$ ; $Z(\mathrm{Cl}) = 17$. Masses molaires atomiques (g/mol) : Ca : 40 ; Mg : 24 ; Na : 23 ; K : 39 ; Cl : 35,5 ; S : 32 ; O : 16 ; H : 1 ; C : 12. On comptera les charges en unités de charge élémentaire ($+1$, $+2$, $-1$, $-2$).

\courssub{Consignes}

\begin{enumerate}
\item \textbf{Classer.} Recopie le tableau de l'étiquette en ajoutant une colonne « cation » ou « anion ». Justifie ton classement à partir de la définition de chaque type d'ion.
\item \textbf{Expliquer avec Bohr.} Dessine le modèle de Bohr simplifié de l'atome de calcium ($Z = 20$, répartition 2-8-8-2) et de l'atome de chlore ($Z = 17$, répartition 2-8-7). Explique pourquoi le calcium forme l'ion \ce{Ca^{2+}} et le chlore l'ion \ce{Cl^{-}}.
\item \textbf{Écrire des formules.} En appliquant la règle d'électroneutralité, écris la formule des trois composés ioniques que l'on obtient en dissolvant : le chlorure de calcium, le sulfate de sodium, l'hydrogénocarbonate de calcium.
\item \textbf{Vérifier l'électroneutralité.} Calcule, pour un litre d'eau, la quantité de matière (en mmol) de chaque ion, puis la somme des charges positives et la somme des charges négatives (en mmol de charges). Conclus.
\item \textbf{Expérimenter.} Propose un protocole (matériel, schéma, étapes) permettant de comparer la conductivité de l'eau minérale et de l'eau distillée. Rédige deux consignes de sécurité.
\item \textbf{Transférer.} Explique, en trois ou quatre phrases, pourquoi un soluté de réhydratation orale contient des ions \ce{Na^{+}}, \ce{K^{+}} et \ce{Cl^{-}}.
\item \textbf{Restituer.} Présente tous tes résultats sur une affiche : schémas de Bohr, tableau complété, calculs, schéma du montage, conclusion.
\end{enumerate}

\courssub{Ressources à mobiliser}

\begin{aretenir}[Boîte à outils de l'activité]
\begin{itemize}
\item Un \cle{cation} est un ion chargé positivement (atome ayant \emph{perdu} des électrons) : \ce{Ca^{2+}}, \ce{Mg^{2+}}, \ce{Na^{+}}, \ce{K^{+}}. Un \cle{anion} est un ion chargé négativement (atome ayant \emph{gagné} des électrons, ou groupement d'atomes chargé négativement) : \ce{Cl^{-}}, \ce{SO4^{2-}}, \ce{HCO3^{-}}.
\item Un atome cherche à acquérir la structure stable en \emph{octet} (8 électrons sur sa couche externe) en perdant ou en gagnant des électrons. Perdre des électrons (négatifs) donne un ion \emph{positif} ; en gagner donne un ion \emph{négatif}.
\item Tout composé ionique et toute solution sont \cle{électriquement neutres} : la somme des charges positives est égale à la somme des charges négatives.
\item Quantité de matière : $n = \dfrac{m}{M}$, avec $m$ la masse et $M$ la masse molaire. Si $m$ est en mg et $M$ en g/mol, alors $n$ s'obtient directement en mmol.
\item Une solution contenant des ions \cle{conduit le courant électrique} ; l'eau distillée (presque sans ions) ne conduit pratiquement pas le courant.
\end{itemize}
\end{aretenir}

\courssub{Démarche et corrigé commenté}

\begin{exoresolu}[Enquête sur l'eau minérale : corrigé complet]
Énoncé : classer les ions de l'étiquette, justifier avec le modèle de Bohr, écrire trois formules de composés ioniques, vérifier l'électroneutralité de l'eau, proposer un test de conductivité et expliquer la composition du SRO.
\tcblower
\textbf{Étape 1 : classement cations / anions.}

Un cation porte des charges $+$ ; un anion porte des charges $-$. On obtient :

\begin{center}
\begin{tabularx}{\linewidth}{|X|X|}\hline
\rowcolor{popblueL} \textbf{Cations} & \textbf{Anions} \\ \hline
\ce{Ca^{2+}}, \ce{Mg^{2+}}, \ce{Na^{+}}, \ce{K^{+}} & \ce{Cl^{-}}, \ce{SO4^{2-}}, \ce{HCO3^{-}} \\ \hline
\end{tabularx}
\end{center}

\textbf{Étape 2 : justification par le modèle de Bohr.}

Calcium ($Z = 20$) : répartition 2-8-8-2. L'atome possède \textbf{2 électrons} sur sa couche externe. Pour obtenir l'octet stable, il \emph{perd} ces 2 électrons : il devient l'ion \ce{Ca^{2+}} (répartition 2-8-8). Le noyau, lui, ne change pas : il garde ses 20 protons ; c'est bien la perte de 2 électrons négatifs qui crée les 2 charges positives.

Chlore ($Z = 17$) : répartition 2-8-7. L'atome possède \textbf{7 électrons} externes ; il lui en manque un pour l'octet. Il \emph{gagne} 1 électron et devient l'ion \ce{Cl^{-}} (répartition 2-8-8).

\begin{popfigure}
\begin{center}
\begin{tikzpicture}[scale=0.85]
% Calcium
\draw[fill=poporangeL] (0,0) circle (0.35) node {\scriptsize Ca};
\draw[popline] (0,0) circle (0.8);
\draw[popline] (0,0) circle (1.3);
\draw[popline] (0,0) circle (1.8);
\draw[popline] (0,0) circle (2.3);
\foreach \a in {90,270} \fill[popblue] (\a:0.8) circle (0.09);
\foreach \a in {0,45,90,135,180,225,270,315} \fill[popblue] (\a:1.3) circle (0.09);
\foreach \a in {22,67,112,157,202,247,292,337} \fill[popblue] (\a:1.8) circle (0.09);
\foreach \a in {90,270} \fill[poporange] (\a:2.3) circle (0.09);
\node[below] at (0,-2.7) {\scriptsize Atome Ca : 2-8-8-\textbf{2}};
% Chlore
\begin{scope}[xshift=7cm]
\draw[fill=popgreenL] (0,0) circle (0.35) node {\scriptsize Cl};
\draw[popline] (0,0) circle (0.8);
\draw[popline] (0,0) circle (1.3);
\draw[popline] (0,0) circle (1.8);
\foreach \a in {90,270} \fill[popblue] (\a:0.8) circle (0.09);
\foreach \a in {0,45,90,135,180,225,270,315} \fill[popblue] (\a:1.3) circle (0.09);
\foreach \a in {0,51,102,153,204,255,306} \fill[popblue] (\a:1.8) circle (0.09);
\node[below] at (0,-2.7) {\scriptsize Atome Cl : 2-8-\textbf{7}};
\end{scope}
\end{tikzpicture}
\end{center}
\legende{Modèles de Bohr simplifiés : les électrons de la couche externe sont en orange (Ca) ; le calcium perd 2 électrons (ion \ce{Ca^{2+}}), le chlore gagne 1 électron (ion \ce{Cl^{-}}).}
\end{popfigure}

\textbf{Étape 3 : formules des composés ioniques.}

Règle : la charge totale positive doit compenser exactement la charge totale négative.
\begin{itemize}
\item Chlorure de calcium : 1 ion \ce{Ca^{2+}} ($+2$) pour 2 ions \ce{Cl^{-}} ($2 \times (-1) = -2$) : \ce{CaCl2}.
\item Sulfate de sodium : 2 ions \ce{Na^{+}} ($2 \times (+1) = +2$) pour 1 ion \ce{SO4^{2-}} ($-2$) : \ce{Na2SO4}.
\item Hydrogénocarbonate de calcium : 1 ion \ce{Ca^{2+}} ($+2$) pour 2 ions \ce{HCO3^{-}} ($2 \times (-1) = -2$) : \ce{Ca(HCO3)2}.
\end{itemize}

\textbf{Étape 4 : vérification de l'électroneutralité (pour 1 L).}

On calcule $n = \dfrac{m}{M}$ en mmol (avec $m$ en mg et $M$ en g/mol), puis les charges apportées. Masses molaires des ions polyatomiques : $M(\ce{SO4}) = 32 + 4 \times 16 = 96$ g/mol ; $M(\ce{HCO3}) = 1 + 12 + 3 \times 16 = 61$ g/mol.

\begin{center}
\begin{tabularx}{\linewidth}{|l|X|X|X|}\hline
\rowcolor{popblueL} \textbf{Ion} & $n$ (mmol/L) & Charge par ion & Charges (mmol/L) \\ \hline
\ce{Ca^{2+}} & $40/40 = 1$ & $+2$ & $+2$ \\ \hline
\ce{Mg^{2+}} & $24/24 = 1$ & $+2$ & $+2$ \\ \hline
\ce{Na^{+}} & $23/23 = 1$ & $+1$ & $+1$ \\ \hline
\ce{K^{+}} & $4/39 \approx 0,1$ & $+1$ & $+0,1$ \\ \hline
\ce{Cl^{-}} & $35,5/35,5 = 1$ & $-1$ & $-1$ \\ \hline
\ce{SO4^{2-}} & $48/96 = 0,5$ & $-2$ & $-1$ \\ \hline
\ce{HCO3^{-}} & $183/61 = 3$ & $-1$ & $-3$ \\ \hline
\end{tabularx}
\end{center}

\begin{align*}
\text{Charges positives} &= 2 + 2 + 1 + 0,1 = 5,1 \ \text{mmol/L}\\
\text{Charges négatives} &= 1 + 1 + 3 = 5,0 \ \text{mmol/L}
\end{align*}

Les deux sommes sont pratiquement égales : le petit écart ($0,1$ mmol/L, soit $2$ \%) provient des arrondis des teneurs affichées sur l'étiquette. L'eau minérale est donc bien globalement \cle{électriquement neutre} : toute solution réelle vérifie cette compensation des charges.

\textbf{Étape 5 : test de conductivité.}

Protocole : réaliser un circuit série comportant un générateur de tension continue de \SI{4,5}{V} (pile plate), une lampe de \SI{3,5}{V} et deux électrodes (tiges de graphite) plongées dans un bécher sans se toucher. Avec l'\textbf{eau distillée}, la lampe ne brille pas : l'eau distillée ne contient pratiquement pas d'ions. Avec l'\textbf{eau minérale}, la lampe brille faiblement : les ions se déplacent et transportent le courant dans la solution. Consignes de sécurité : ne jamais utiliser le courant du secteur (\SI{220}{V}) pour cette expérience ; ne pas toucher les électrodes lorsque le circuit est fermé et essuyer toute eau renversée.

\textbf{Étape 6 : le SRO.} Lors d'une diarrhée, le corps perd beaucoup d'eau et des ions \ce{Na^{+}}, \ce{K^{+}} et \ce{Cl^{-}}. Le SRO remplace ces pertes : les ions sodium permettent à l'eau d'être absorbée par l'intestin, le potassium rééquilibre le fonctionnement des cellules, et les ions chlorure maintiennent la neutralité électrique des liquides du corps. C'est une application directe, qui sauve des vies, de la chimie des ions étudiée dans cette leçon.
\end{exoresolu}

\begin{attention}[Pièges fréquents]
Ne confonds pas le \emph{nombre d'ions} et la \emph{charge totale} : un ion \ce{Ca^{2+}} compte pour \emph{deux} charges positives. Dans \ce{Ca(HCO3)2}, les parenthèses signifient que \emph{tout} le groupe \ce{HCO3} est pris deux fois (et pas seulement l'atome d'hydrogène). Enfin, dans le calcul de $n$, utilise des unités cohérentes : mg avec g/mol donnent des mmol.
\end{attention}

\courssub{Bilan de l'activité}

Cette enquête a permis de relier les trois notions de la leçon dans une situation authentique : l'\cle{atome} (modèle de Bohr) explique la naissance des \cle{ions}, les ions s'associent en \cle{composés ioniques} neutres et donnent à l'eau minérale ses propriétés (conductivité, rôle pour la santé). Tu as aussi appris à lire une étiquette avec un regard scientifique : derrière la mention « riche en sels minéraux » se cachent des cations et des anions dont les charges se compensent. La restitution en affiche développe enfin la compétence communication, essentielle au BEPC et dans la vie citoyenne.

\newpage

\coursec{Méthodes et exercices résolus}

\coursec{Les constituants de la matière}

\textbf{Situation d'accroche.} Dans une maison de Yaoundé, on trouve une marmite en aluminium, une bonbonne de gaz butane, du sel de cuisine et des fils électriques en cuivre. Ces matériaux si différents sont-ils faits des mêmes « briques » ? Toute la matière qui nous entoure — l'eau que tu bois, l'air que tu respires, le métal des moto-taxis — est construite à partir d'entités microscopiques : les \cle{atomes}, les \cle{molécules} et les \cle{ions}. Cette leçon te donne les clés pour les décrire, les nommer et les utiliser, comme l'exige le programme de 3ème et les épreuves du BEPC.

\courssub{Méthode 1 : Décrire la constitution d'un atome}

\begin{definition}[L'atome]
L'\cle{atome} est la plus petite partie d'un corps simple qui conserve les propriétés de ce corps. Il est constitué :
\begin{itemize}
\item d'un \cle{noyau} central, contenant des \cle{protons} (charge positive $+e$) et des \cle{neutrons} (sans charge électrique) ;
\item d'un \cle{nuage électronique} : des \cle{électrons} (charge négative $-e$) qui tournent autour du noyau.
\end{itemize}
Le nombre de protons du noyau est le \cle{numéro atomique}, noté $Z$. Un atome est électriquement \cle{neutre} : il possède autant d'électrons que de protons.
\end{definition}

\begin{methode}[Décrire la constitution d'un atome]
Pour décrire un atome à partir de son symbole et de son numéro atomique $Z$ :
\begin{enumerate}
\item \cle{Identifier le noyau} : il contient les \cle{protons} (charge $+e$) et les \cle{neutrons} (sans charge). Le nombre de protons est égal au numéro atomique $Z$.
\item \cle{Compter les électrons} : un atome est électriquement \cle{neutre}, donc le nombre d'électrons (charge $-e$) est égal au nombre de protons : $N_{\text{électrons}} = Z$.
\item \cle{Répartir les électrons} selon le modèle de Bohr : la couche $K$ reçoit au maximum 2 électrons, la couche $L$ au maximum 8, puis la couche $M$.
\item \cle{Conclure sur la neutralité} : charge totale du nuage électronique $= -Z \cdot e$, charge du noyau $= +Z \cdot e$, somme nulle.
\item Rédiger une phrase de conclusion complète, en citant le nom de l'élément et son symbole.
\end{enumerate}
\textbf{Réflexes :} ordre de grandeur du rayon atomique : $10^{-10}$ m ; le noyau est environ $10^{5}$ fois plus petit que l'atome ; la masse de l'atome est concentrée dans le noyau.
\end{methode}

\begin{popfigure}
\begin{center}
\begin{tikzpicture}[scale=1.0]
% Noyau
\fill[poporange] (0,0) circle (0.35);
\node[white,font=\small\bfseries] at (0,0) {noyau};
% Couches
\draw[popblue,thick] (0,0) circle (0.9);
\draw[popblue,thick] (0,0) circle (1.6);
\draw[popblue,thick] (0,0) circle (2.3);
% Electrons K (2)
\fill[popdark] (0.9,0) circle (0.09);
\fill[popdark] (-0.9,0) circle (0.09);
% Electrons L (8)
\foreach \a in {0,45,90,135,180,225,270,315} {\fill[popdark] (\a:1.6) circle (0.09);}
% Electrons M (3)
\foreach \a in {30,150,270} {\fill[popdark] (\a:2.3) circle (0.09);}
% Labels
\node[popblue,font=\small] at (0.75,0.75) {$K$};
\node[popblue,font=\small] at (1.25,1.25) {$L$};
\node[popblue,font=\small] at (1.75,1.75) {$M$};
\end{tikzpicture}
\end{center}
\legende{Modèle de Bohr de l'atome d'aluminium \ce{Al} ($Z = 13$) : le noyau (13 protons et 14 neutrons) est entouré de 13 électrons répartis sur les couches $K$ (2 électrons), $L$ (8 électrons) et $M$ (3 électrons).}
\end{popfigure}

\begin{exoresolu}[L'atome d'aluminium des marmites]
Les marmites utilisées dans les cuisines camerounaises sont souvent en aluminium. L'atome d'aluminium, de symbole \ce{Al}, possède un numéro atomique $Z = 13$.
\begin{enumerate}
\item Donner la constitution de cet atome (noyau et nuage électronique).
\item Répartir les électrons sur les couches du modèle de Bohr.
\item Justifier la neutralité électrique de cet atome.
\end{enumerate}
\tcblower
\textbf{1. Constitution de l'atome.}

Le numéro atomique $Z = 13$ indique que le noyau de l'atome d'aluminium contient \textbf{13 protons}, chacun portant la charge $+e$. Le noyau contient aussi des neutrons, sans charge électrique (ici 14 neutrons, car le noyau d'aluminium contient au total 27 nucléons : $27 - 13 = 14$).

L'atome étant électriquement neutre, le nuage électronique contient autant d'électrons que le noyau contient de protons, soit \textbf{13 électrons}, chacun portant la charge $-e$.

\textbf{2. Répartition des électrons (modèle de Bohr).}

On remplit les couches dans l'ordre, en respectant les capacités maximales :
\begin{itemize}
\item couche $K$ : 2 électrons au maximum $\rightarrow$ on place 2 électrons ;
\item couche $L$ : 8 électrons au maximum $\rightarrow$ on place 8 électrons ;
\item couche $M$ : il reste $13 - 2 - 8 = 3$ électrons.
\end{itemize}
La répartition s'écrit : $(K)^2 (L)^8 (M)^3$.

\textbf{3. Neutralité électrique.}

Charge du noyau : $Q_{\text{noyau}} = +13\,e$. Charge du nuage électronique : $Q_{\text{nuage}} = -13\,e$.
\[ Q_{\text{atome}} = Q_{\text{noyau}} + Q_{\text{nuage}} = +13\,e - 13\,e = 0. \]
\textbf{Conclusion :} l'atome d'aluminium \ce{Al} ($Z = 13$) est formé d'un noyau contenant 13 protons (et 14 neutrons) et d'un nuage de 13 électrons répartis selon $(K)^2 (L)^8 (M)^3$ ; les charges se compensant exactement, l'atome est électriquement neutre.
\end{exoresolu}

\courssub{Méthode 2 : Écrire et interpréter le symbole d'un élément chimique}

\begin{methode}[Utiliser les symboles des éléments]
\begin{enumerate}
\item \cle{Respecter l'écriture} : le symbole commence toujours par une \textbf{majuscule}, suivie éventuellement d'une minuscule : \ce{H}, \ce{O}, \ce{Na}, \ce{Cl}, \ce{Fe}, \ce{Cu}.
\item \cle{Associer symbole et nom} : certains symboles viennent du latin : \ce{Na} (natrium = sodium), \ce{Fe} (ferrum = fer), \ce{Cu} (cuprum = cuivre), \ce{K} (kalium = potassium), \ce{Au} (aurum = or).
\item \cle{Ne pas confondre} atome et élément : le symbole \ce{O} désigne un atome d'oxygène ; le corps simple dioxygène s'écrit \ce{O2}.
\item \cle{Exploiter le tableau périodique} : chaque case donne le symbole, le numéro atomique $Z$ et des informations sur l'élément.
\end{enumerate}
\textbf{Réflexe :} dans une copie, écris les symboles soigneusement : \ce{Co} est l'atome de cobalt, \ce{CO} est la molécule de monoxyde de carbone (gaz toxique dégagé par les groupes électrogènes mal ventilés !).
\end{methode}

\begin{exoresolu}[Symboles autour de la cuisine]
Dans une cuisine, on trouve : du sel de cuisine (qui contient du sodium et du chlore), une casserole en aluminium, des couverts en fer et des fils électriques en cuivre.
\begin{enumerate}
\item Écrire le symbole de chacun des éléments cités : sodium, chlore, aluminium, fer, cuivre.
\item Un élève écrit le symbole du sodium « NA ». Corriger son erreur et expliquer la règle.
\end{enumerate}
\tcblower
\textbf{1. Symboles des éléments.}
\begin{center}
\begin{tabularx}{\linewidth}{|l|X|X|}\hline
\rowcolor{popblueL} \textbf{Élément} & \textbf{Symbole} & \textbf{Origine du symbole} \\ \hline
Sodium & \ce{Na} & du latin \emph{natrium} \\ \hline
Chlore & \ce{Cl} & première lettre + lettre caractéristique \\ \hline
Aluminium & \ce{Al} & première lettre + lettre caractéristique \\ \hline
Fer & \ce{Fe} & du latin \emph{ferrum} \\ \hline
Cuivre & \ce{Cu} & du latin \emph{cuprum} \\ \hline
\end{tabularx}
\end{center}

\textbf{2. Correction de l'erreur.}

L'écriture « NA » est fausse. La règle internationale impose : la \textbf{première lettre en majuscule}, la \textbf{deuxième en minuscule}. Le symbole correct du sodium est donc \ce{Na}. Écrire « NA » laisserait croire à deux atomes distincts \ce{N} (azote) et \ce{A} (qui n'existe pas).

\textbf{Conclusion :} sodium \ce{Na}, chlore \ce{Cl}, aluminium \ce{Al}, fer \ce{Fe}, cuivre \ce{Cu} ; un symbole s'écrit toujours avec une majuscule suivie, si nécessaire, d'une minuscule.
\end{exoresolu}

\courssub{Méthode 3 : Écrire et interpréter la formule d'une molécule}

\begin{definition}[La molécule]
Une \cle{molécule} est un assemblage d'atomes liés entre eux, électriquement neutre. Sa \cle{formule} indique la nature et le nombre des atomes : l'\cle{indice}, écrit en bas à droite du symbole, précise le nombre d'atomes de cet élément (l'absence d'indice signifie 1 atome).
\end{definition}

\begin{methode}[Lire et écrire une formule moléculaire]
\begin{enumerate}
\item \cle{Lire les indices} : dans \ce{H2O}, l'indice 2 signifie 2 atomes d'hydrogène ; l'absence d'indice après \ce{O} signifie 1 atome d'oxygène. La molécule d'eau contient donc 3 atomes au total.
\item \cle{Distinguer coefficient et indice} : dans $2\,\ce{H2O}$, le coefficient 2 (devant) multiplie toute la molécule : il y a 2 molécules d'eau, soit 4 atomes d'hydrogène et 2 atomes d'oxygène. L'indice 2 (en bas) ne concerne que l'élément qui le précède.
\item \cle{Écrire une formule} à partir du nom : dioxygène = 2 atomes d'oxygène liés $\rightarrow$ \ce{O2} ; dioxyde de carbone = 1 carbone + 2 oxygènes $\rightarrow$ \ce{CO2}.
\item \cle{Représenter} la molécule par un modèle compact : un cercle coloré par atome, les cercles se touchent.
\end{enumerate}
\textbf{Formule à retenir :} nombre total d'atomes = somme des indices (en comptant 1 quand il n'y a pas d'indice).
\end{methode}

\begin{popfigure}
\begin{center}
\begin{tikzpicture}[scale=0.9]
% H2O
\fill[popblue] (0,0) circle (0.45);
\fill[popgreen!70!black] (-0.55,0.45) circle (0.3);
\fill[popgreen!70!black] (0.55,0.45) circle (0.3);
\node[white,font=\small\bfseries] at (0,0) {O};
\node[white,font=\footnotesize\bfseries] at (-0.55,0.45) {H};
\node[white,font=\footnotesize\bfseries] at (0.55,0.45) {H};
\node[font=\small] at (0,-0.85) {\ce{H2O}};
% O2
\fill[popblue] (3,0.2) circle (0.45);
\fill[popblue] (3.85,0.2) circle (0.45);
\node[white,font=\small\bfseries] at (3,0.2) {O};
\node[white,font=\small\bfseries] at (3.85,0.2) {O};
\node[font=\small] at (3.4,-0.85) {\ce{O2}};
% CO2
\fill[popblue] (6,0.2) circle (0.45);
\fill[popdark] (6.85,0.2) circle (0.4);
\fill[popblue] (7.65,0.2) circle (0.45);
\node[white,font=\small\bfseries] at (6,0.2) {O};
\node[white,font=\small\bfseries] at (6.85,0.2) {C};
\node[white,font=\small\bfseries] at (7.65,0.2) {O};
\node[font=\small] at (6.85,-0.85) {\ce{CO2}};
\end{tikzpicture}
\end{center}
\legende{Modèles compacts de trois molécules : l'eau \ce{H2O} (2 atomes d'hydrogène, 1 atome d'oxygène), le dioxygène \ce{O2} (2 atomes d'oxygène) et le dioxyde de carbone \ce{CO2} (1 atome de carbone, 2 atomes d'oxygène).}
\end{popfigure}

\begin{exoresolu}[Le gaz des bonbonnes et l'eau de boisson]
La bonbonne de gaz domestique contient du butane, de formule \ce{C4H10}. L'eau minérale a pour formule \ce{H2O}.
\begin{enumerate}
\item Donner le nom et le nombre de chaque type d'atome présent dans une molécule de butane.
\item Combien d'atomes au total compte une molécule de butane ? Une molécule d'eau ?
\item On considère $3\,\ce{H2O}$. Combien y a-t-il de molécules d'eau ? D'atomes d'hydrogène ? D'atomes d'oxygène ?
\end{enumerate}
\tcblower
\textbf{1. Composition du butane \ce{C4H10}.}

L'indice 4 après \ce{C} indique 4 atomes de \textbf{carbone} ; l'indice 10 après \ce{H} indique 10 atomes d'\textbf{hydrogène}.

\textbf{2. Nombre total d'atomes.}

Butane : $4 + 10 = 14$ atomes. Eau \ce{H2O} : $2 + 1 = 3$ atomes.

\textbf{3. Interprétation de $3\,\ce{H2O}$.}

Le coefficient 3 placé devant la formule multiplie toute la molécule :
\begin{itemize}
\item nombre de molécules d'eau : 3 ;
\item atomes d'hydrogène : $3 \times 2 = 6$ ;
\item atomes d'oxygène : $3 \times 1 = 3$.
\end{itemize}
Vérification : $6 + 3 = 9$ atomes au total, soit bien $3 \times 3$ atomes.

\textbf{Conclusion :} le butane \ce{C4H10} contient 4 atomes de carbone et 10 atomes d'hydrogène (14 atomes) ; l'eau \ce{H2O} en contient 3 ; l'écriture $3\,\ce{H2O}$ représente 3 molécules d'eau, soit 6 atomes d'hydrogène et 3 atomes d'oxygène.
\end{exoresolu}

\courssub{Méthode 4 : Déterminer la formule et la charge d'un ion simple}

\begin{definition}[L'ion]
Un \cle{ion} est un atome (ou un groupe d'atomes) qui a gagné ou perdu un ou plusieurs électrons. Un ion est donc \textbf{chargé électriquement}, contrairement à l'atome qui est neutre.
\end{definition}

\begin{methode}[Passer de l'atome à l'ion simple]
\begin{enumerate}
\item \cle{Perte d'électrons} (charges négatives en moins) $\rightarrow$ ion positif = \cle{cation}. Exemple : \ce{Na -> Na+ + e-}.
\item \cle{Gain d'électrons} $\rightarrow$ ion négatif = \cle{anion}. Exemple : \ce{Cl + e- -> Cl-}.
\item \cle{Écrire la charge} en exposant, à droite du symbole : \ce{Na+}, \ce{Cl-}, \ce{Mg^{2+}}, \ce{O^{2-}}, \ce{Al^{3+}}.
\item \cle{Compter les électrons de l'ion} : $N_{\text{électrons}} = Z - (\text{charge algébrique})$. Exemple : \ce{Mg^{2+}} avec $Z = 12$ possède $12 - 2 = 10$ électrons.
\item \cle{Le noyau ne change jamais} : un ion et son atome ont le même nombre de protons ; c'est le même élément chimique.
\end{enumerate}
\textbf{Réflexe :} les métaux (\ce{Na}, \ce{Mg}, \ce{Al}, \ce{Fe}, \ce{Cu}) forment des cations ; les non-métaux (\ce{Cl}, \ce{O}, \ce{S}) forment des anions.
\end{methode}

\begin{exoresolu}[Les ions du sel de cuisine et de l'eau minérale]
Sur l'étiquette d'une eau minérale camerounaise, on lit : calcium \ce{Ca^{2+}}, magnésium \ce{Mg^{2+}}, chlorure \ce{Cl-}. Données : $Z(\ce{Ca}) = 20$, $Z(\ce{Mg}) = 12$, $Z(\ce{Cl}) = 17$.
\begin{enumerate}
\item L'ion calcium \ce{Ca^{2+}} est-il un cation ou un anion ? Justifier.
\item Déterminer le nombre d'électrons de chacun des trois ions.
\item Comparer le nombre de protons de l'atome de chlore et de l'ion chlorure.
\end{enumerate}
\tcblower
\textbf{1. Nature de l'ion calcium.}

L'ion \ce{Ca^{2+}} porte une charge \textbf{positive} ($2+$) : c'est un \textbf{cation}. L'atome de calcium a perdu 2 électrons ; il possède donc 2 charges négatives de moins que de charges positives, d'où la charge globale $+2e$.

\textbf{2. Nombre d'électrons des ions.}

On applique $N_{\text{électrons}} = Z - (\text{charge algébrique})$ :
\begin{align*}
\ce{Ca^{2+}} : \quad & N = 20 - 2 = 18 \text{ électrons} ;\\
\ce{Mg^{2+}} : \quad & N = 12 - 2 = 10 \text{ électrons} ;\\
\ce{Cl-} : \quad & N = 17 - (-1) = 17 + 1 = 18 \text{ électrons}.
\end{align*}
L'ion chlorure a \textbf{gagné} un électron par rapport à l'atome de chlore, ce qui est cohérent avec sa charge négative.

\textbf{3. Comparaison atome de chlore / ion chlorure.}

La transformation d'un atome en ion ne concerne que les électrons : le noyau est inchangé. L'atome \ce{Cl} et l'ion \ce{Cl-} possèdent donc tous deux \textbf{17 protons}.

\textbf{Conclusion :} \ce{Ca^{2+}} est un cation (perte de 2 électrons) ; \ce{Ca^{2+}} a 18 électrons, \ce{Mg^{2+}} en a 10 et \ce{Cl-} en a 18 ; atome et ion d'un même élément ont toujours le même nombre de protons.
\end{exoresolu}

\courssub{Méthode 5 : Connaître les ions polyatomiques et écrire la formule d'un composé ionique}

\begin{methode}[Ions polyatomiques et composés ioniques]
\begin{enumerate}
\item \cle{Mémoriser les ions polyatomiques usuels} : sulfate \ce{SO4^{2-}}, nitrate \ce{NO3-}, hydroxyde \ce{OH-}, carbonate \ce{CO3^{2-}}, ammonium \ce{NH4+} (seul cation polyatomique courant).
\item \cle{Un ion polyatomique est un groupe d'atomes liés} qui porte une charge globale : dans \ce{SO4^{2-}}, les 4 atomes d'oxygène et l'atome de soufre restent soudés ; c'est l'ensemble qui porte 2 charges négatives.
\item \cle{Écrire un composé ionique} : un solide ionique est électriquement neutre. On combine cations et anions de sorte que la somme des charges soit nulle. Exemple : \ce{Ca^{2+}} + 2 \ce{Cl-} $\rightarrow$ \ce{CaCl2} (car $+2 + 2 \times (-1) = 0$).
\item \cle{Utiliser des parenthèses} si l'ion polyatomique est en plusieurs exemplaires : \ce{Ca^{2+}} + 2 \ce{OH-} $\rightarrow$ \ce{Ca(OH)2}.
\item \cle{Vérifier systématiquement} la neutralité électrique de la formule obtenue.
\end{enumerate}
\end{methode}

\begin{exoresolu}[Le sulfate de cuivre du laboratoire]
Au laboratoire du collège, le professeur utilise du sulfate de cuivre, solide bleu formé des ions cuivre(II) \ce{Cu^{2+}} et sulfate \ce{SO4^{2-}}.
\begin{enumerate}
\item L'ion sulfate est-il un ion simple ou polyatomique ? Justifier.
\item Écrire la formule du sulfate de cuivre en justifiant par la neutralité électrique.
\item Écrire la formule de l'hydroxyde de sodium, composé des ions \ce{Na+} et \ce{OH-}.
\end{enumerate}
\tcblower
\textbf{1. Nature de l'ion sulfate.}

L'ion sulfate \ce{SO4^{2-}} est formé de \textbf{plusieurs atomes liés} (1 atome de soufre et 4 atomes d'oxygène) portant globalement 2 charges négatives : c'est un \textbf{ion polyatomique}.

\textbf{2. Formule du sulfate de cuivre.}

Le composé doit être électriquement neutre. L'ion cuivre(II) porte $+2$ et l'ion sulfate porte $-2$ :
\[ (+2) + (-2) = 0. \]
Un seul ion de chaque sorte suffit. La formule du sulfate de cuivre est donc \ce{CuSO4}.

\textbf{3. Formule de l'hydroxyde de sodium.}

L'ion sodium porte $+1$ et l'ion hydroxyde porte $-1$ :
\[ (+1) + (-1) = 0. \]
La formule est \ce{NaOH} (un seul ion hydroxyde, donc pas de parenthèses).

\textbf{Conclusion :} \ce{SO4^{2-}} est un ion polyatomique ; le sulfate de cuivre a pour formule \ce{CuSO4} et l'hydroxyde de sodium \ce{NaOH}, les deux formules assurant la neutralité électrique du composé.
\end{exoresolu}

\courssub{Méthode 6 : Identifier et interpréter un test de reconnaissance (démarche expérimentale)}

\begin{methode}[Conduire une démarche d'investigation en chimie]
\begin{enumerate}
\item \cle{Observer} et poser le problème (ex. : cette eau de puits contient-elle des ions ?).
\item \cle{Formuler une hypothèse} vérifiable (ex. : « si l'eau conduit le courant, alors elle contient des ions »).
\item \cle{Réaliser l'expérience} en respectant les consignes de sécurité : ne jamais goûter un produit, porter des lunettes, ne pas verser les réactifs dans l'évier, manipuler sous la hotte si nécessaire.
\item \cle{Noter les observations} de façon objective (couleur, précipité, lampe qui s'allume).
\item \cle{Interpréter} : relier l'observation à la présence ou à l'absence de l'espèce chimique cherchée.
\item \cle{Conclure} en validant ou en rejetant l'hypothèse.
\end{enumerate}
\textbf{Réflexes :} l'eau pure ne conduit pas le courant ; une solution contenant des ions conduit le courant (les ions sont des porteurs de charge mobiles). Test des ions chlorure : quelques gouttes de nitrate d'argent $\rightarrow$ précipité blanc de chlorure d'argent \ce{AgCl} qui noircit à la lumière.
\end{methode}

\begin{popfigure}
\begin{center}
\begin{tikzpicture}[scale=0.95]
% Becher
\draw[thick,popdark] (-1.6,-1.8) -- (-1.6,0.6);
\draw[thick,popdark] (1.6,-1.8) -- (1.6,0.6);
\draw[thick,popdark] (-1.6,-1.8) -- (1.6,-1.8);
% Solution
\fill[popblue!25] (-1.55,-1.75) rectangle (1.55,-0.4);
\draw[popblue] (-1.55,-0.4) -- (1.55,-0.4);
\node[font=\small,popdark] at (0,-1.1) {solution ionique};
% Electrodes
\draw[thick,fill=popmuted] (-0.7,-0.9) rectangle (-0.5,1.3);
\draw[thick,fill=popmuted] (0.5,-0.9) rectangle (0.7,1.3);
% Fils vers generateur et lampe
\draw[thick] (-0.6,1.3) -- (-0.6,2.1) -- (-3,2.1) -- (-3,0.8);
\draw[thick] (0.6,1.3) -- (0.6,2.1) -- (3,2.1) -- (3,1.3);
% Generateur
\draw[thick,fill=popcream] (-3.5,0.2) rectangle (-2.5,0.8);
\node[font=\small] at (-3,0.5) {$+$ $-$};
\node[font=\small,below] at (-3,0.15) {générateur};
% Lampe
\draw[thick] (3,1.3) -- (3,0.9);
\draw[thick,fill=popgold!40] (3,0.55) circle (0.35);
\draw[thick] (2.83,0.38) -- (3.17,0.72);
\draw[thick] (2.83,0.72) -- (3.17,0.38);
\draw[thick] (3,0.2) -- (3,-0.6) -- (-3,-0.6) -- (-3,0.2);
\node[font=\small,right] at (3.45,0.55) {lampe};
% Ions dans la solution
\node[poporange,font=\footnotesize\bfseries] at (-1.1,-0.7) {$+$};
\node[popgreen!60!black,font=\footnotesize\bfseries] at (1.1,-1.4) {$-$};
\node[poporange,font=\footnotesize\bfseries] at (0.2,-1.5) {$+$};
\node[popgreen!60!black,font=\footnotesize\bfseries] at (-0.3,-0.6) {$-$};
\node[poporange,font=\footnotesize\bfseries] at (1.2,-0.6) {$+$};
\node[popgreen!60!black,font=\footnotesize\bfseries] at (-1.2,-1.5) {$-$};
\end{tikzpicture}
\end{center}
\legende{Schéma du montage de test de conductivité : un générateur de très basse tension, une lampe témoin et deux électrodes plongées dans la solution. Si la lampe brille, la solution contient des ions mobiles ($+$ et $-$).}
\end{popfigure}

\begin{exoresolu}[L'eau du robinet conduit-elle le courant ?]
Un groupe d'élèves de 3ème veut savoir si l'eau du robinet du collège contient des ions. Ils réalisent le circuit suivant : un générateur, une lampe et deux électrodes plongées dans un bécher d'eau du robinet. La lampe brille faiblement. Avec de l'eau distillée, la lampe reste éteinte.
\begin{enumerate}
\item Formuler l'hypothèse testée par les élèves.
\item Interpréter les deux observations.
\item Conclure et citer une précaution de sécurité pour cette expérience.
\end{enumerate}
\tcblower
\textbf{1. Hypothèse.}

« Si l'eau du robinet contient des ions, alors elle conduit le courant électrique et la lampe brille. »

\textbf{2. Interprétation.}

Avec l'eau du robinet, la lampe brille : le courant traverse la solution, donc des \textbf{porteurs de charge mobiles}, c'est-à-dire des \textbf{ions}, sont présents (calcium, magnésium, chlorure, etc.). La brillance est faible car la concentration en ions est faible.

Avec l'eau distillée, la lampe ne brille pas : l'eau distillée est presque pure, elle ne contient pratiquement pas d'ions ; or les molécules d'eau \ce{H2O}, électriquement neutres, ne transportent pas le courant.

\textbf{3. Conclusion et sécurité.}

L'hypothèse est \textbf{validée} : l'eau du robinet contient des ions, contrairement à l'eau distillée. Précaution : ne jamais approcher le montage électrique les mains mouillées et utiliser une \textbf{très basse tension} (générateur de \SI{6}{V} maximum) pour éviter tout danger.

\textbf{Conclusion générale :} la conduction du courant par une solution est la preuve de la présence d'ions ; c'est pourquoi l'eau ordinaire est conductrice alors que l'eau pure ne l'est pas.
\end{exoresolu}

\courssub{Méthode 7 : Relier atome, molécule, ion et matériau du quotidien}

\begin{methode}[Raisonner de l'échelle microscopique à l'échelle macroscopique]
\begin{enumerate}
\item \cle{Identifier l'espèce chimique} présente dans le matériau : atome (cuivre \ce{Cu} d'un fil électrique), molécule (eau \ce{H2O}, butane \ce{C4H10}) ou ions (sel \ce{Na+} et \ce{Cl-}).
\item \cle{Distinguer les échelles} : un atome mesure environ $10^{-10}$ m ; un grain de sel visible contient des milliards de milliards d'ions.
\item \cle{Relier structure et propriété} : les métaux conduisent le courant grâce à leurs électrons mobiles ; les solutions ioniques conduisent grâce aux ions mobiles ; les solides ioniques sont durs mais cassants car leurs ions sont ordonnés.
\item \cle{Conclure} en reliant le matériau à son usage quotidien (fils en cuivre pour l'électricité, sel pour l'assaisonnement, butane pour la cuisson).
\end{enumerate}
\end{methode}

\begin{exoresolu}[Trois matériaux de la maison]
Dans une maison à Yaoundé, on trouve : des fils électriques en cuivre, du sel de cuisine et une bonbonne de butane.
\begin{enumerate}
\item Pour chaque matériau, préciser si l'espèce chimique constitutive est un atome, une molécule ou des ions.
\item Expliquer pourquoi le fil de cuivre conduit le courant alors que le sel sec ne le conduit pas.
\item Pourquoi l'eau salée conduit-elle le courant ?
\end{enumerate}
\tcblower
\textbf{1. Nature des espèces chimiques.}
\begin{itemize}
\item Cuivre du fil : ensemble d'\textbf{atomes} de cuivre \ce{Cu} (corps simple métallique).
\item Sel de cuisine : ensemble d'\textbf{ions} sodium \ce{Na+} et chlorure \ce{Cl-} régulièrement ordonnés (composé ionique \ce{NaCl}).
\item Butane de la bonbonne : ensemble de \textbf{molécules} \ce{C4H10} (composé moléculaire).
\end{itemize}

\textbf{2. Conduction du cuivre et non-conduction du sel sec.}

Dans le métal cuivre, des électrons sont \textbf{mobiles} : mis en mouvement par un générateur, ils transportent le courant. Dans le sel sec, les ions \ce{Na+} et \ce{Cl-} existent mais sont \textbf{immobilisés} dans le cristal : sans porteurs de charge mobiles, le courant ne passe pas.

\textbf{3. Conduction de l'eau salée.}

Dissous dans l'eau, les ions \ce{Na+} et \ce{Cl-} se séparent et deviennent \textbf{mobiles} : ils transportent le courant. C'est pourquoi l'eau salée est conductrice alors que le sel sec et l'eau pure ne le sont pas.

\textbf{Conclusion :} la conduction électrique exige des porteurs de charge \emph{mobiles} : électrons libres dans les métaux, ions en solution ; la connaissance des constituants microscopiques explique les usages des matériaux de la vie quotidienne.
\end{exoresolu}

\begin{aretenir}[L'essentiel de la leçon]
\begin{itemize}
\item L'\cle{atome} = noyau (protons $+e$ et neutrons) + électrons ($-e$) ; il est électriquement neutre : nombre d'électrons = nombre de protons = $Z$.
\item Le \cle{symbole} d'un élément s'écrit avec une majuscule, parfois suivie d'une minuscule (\ce{Na}, \ce{Cl}, \ce{Fe}).
\item Une \cle{molécule} est un assemblage neutre d'atomes liés ; sa formule indique la nature et le nombre des atomes (\ce{H2O}, \ce{C4H10}).
\item Un \cle{ion} est un atome (ou groupe d'atomes) ayant gagné ou perdu des électrons : \cle{cation} positif (\ce{Na+}, \ce{Ca^{2+}}), \cle{anion} négatif (\ce{Cl-}, \ce{SO4^{2-}}). Le noyau, lui, ne change jamais.
\item Un \cle{composé ionique} est électriquement neutre : la somme des charges des ions est nulle (\ce{NaCl}, \ce{CaCl2}, \ce{CuSO4}).
\item Une solution contenant des ions \textbf{mobiles} conduit le courant électrique ; l'eau pure et le sel sec ne conduisent pas.
\end{itemize}
\end{aretenir}

\begin{savaistu}[Le sel, l'iode et la santé au Cameroun]
Au Cameroun, le sel de cuisine vendu dans les marchés est \emph{iodé} : on y ajoute de petites quantités d'ions iodure \ce{I-} pour prévenir les troubles de la thyroïde, notamment le goitre, autrefois fréquents dans certaines régions. C'est un bel exemple d'utilisation des connaissances sur les ions pour protéger la santé publique. Par ailleurs, le monoxyde de carbone \ce{CO}, gaz invisible et inodore dégagé par les groupes électrogènes et les charbons de bois mal ventilés, est mortel : une molécule microscopique peut ainsi expliquer un danger bien réel, d'où l'importance d'aérer les pièces.
\end{savaistu}

\begin{attention}[Les 5 erreurs qui coûtent des points au BEPC]
\begin{enumerate}
\item \textbf{Confondre atome, molécule et ion.} \ce{O} est un atome, \ce{O2} une molécule, \ce{O^{2-}} un ion. Écrire « la molécule de sodium \ce{Na} » ou « l'atome d'eau » est une erreur de vocabulaire sanctionnée : le sodium existe sous forme d'atomes (ou d'ions \ce{Na+}), l'eau sous forme de molécules \ce{H2O}.
\item \textbf{Mal écrire les symboles et les charges.} Une majuscule obligatoire en tête du symbole (\ce{Na} et non « NA » ou « na »). La charge d'un ion s'écrit en \emph{exposant} avec le nombre avant le signe : \ce{Mg^{2+}} et non « Mg+2 ». Ne pas confondre \ce{Co} (cobalt) et \ce{CO} (monoxyde de carbone).
\item \textbf{Croire que l'ion a un noyau différent de l'atome.} Passer de \ce{Cl} à \ce{Cl-} ne change que le nombre d'électrons ; le nombre de protons ($Z = 17$) reste identique. Dire « l'ion chlorure a 18 protons » est faux : il a 17 protons et 18 électrons.
\item \textbf{Oublier la neutralité électrique dans les formules ioniques.} Le chlorure de calcium s'écrit \ce{CaCl2} (il faut 2 ions \ce{Cl-} pour 1 ion \ce{Ca^{2+}}), et non « CaCl ». Toujours vérifier que la somme des charges est nulle, et mettre des parenthèses pour les ions polyatomiques multiples : \ce{Ca(OH)2} et non « CaOH2 ».
\item \textbf{Confondre coefficient et indice, et bâcler la conclusion.} Dans $3\,\ce{H2O}$, il y a 6 atomes d'hydrogène (et non 2) : le coefficient multiplie toute la molécule. Enfin, toute question « Justifier » ou « Conclure » exige une \textbf{phrase complète} avec la justification (calcul de charge, observation expérimentale), pas un simple « oui » ou « non ».
\end{enumerate}
\end{attention}

\newpage

\coursec{Exercices}

\courssub{Je vérifie mes connaissances}

\begin{exercice}[QCM : l'atome et ses constituants]
Pour chaque question, choisis la (ou les) bonne(s) réponse(s).
\begin{enumerate}
\item Le noyau d'un atome contient :
\begin{enumerate}
\item des protons et des électrons ;
\item des protons et des neutrons ;
\item uniquement des électrons.
\end{enumerate}
\item L'électron est une particule :
\begin{enumerate}
\item chargée positivement ;
\item chargée négativement ;
\item sans charge électrique.
\end{enumerate}
\item Un atome est électriquement neutre parce que :
\begin{enumerate}
\item il ne contient aucune particule chargée ;
\item il contient autant de protons que d'électrons ;
\item il contient autant de neutrons que d'électrons.
\end{enumerate}
\item Dans le modèle de Bohr, les électrons :
\begin{enumerate}
\item sont répartis sur des couches autour du noyau ;
\item sont mélangés au noyau ;
\item sont immobiles au centre de l'atome.
\end{enumerate}
\end{enumerate}
\end{exercice}

\begin{exercice}[Vrai ou faux, en justifiant]
Réponds par \textbf{Vrai} ou \textbf{Faux}, puis justifie ta réponse en une phrase.
\begin{enumerate}
\item La formule \ce{H2O} signifie que la molécule d'eau contient deux atomes d'oxygène et un atome d'hydrogène.
\item L'ion \ce{Na+} est un atome de sodium qui a perdu un électron.
\item L'ion \ce{Cl-} est un cation.
\item L'ion sulfate \ce{SO4^{2-}} est un ion polyatomique.
\item Un composé ionique comme le sel de cuisine \ce{NaCl} est électriquement neutre.
\end{enumerate}
\end{exercice}

\begin{exercice}[Définitions à compléter]
Recopie et complète les phrases suivantes avec les mots ou expressions : \emph{noyau}, \emph{électrons}, \emph{molécule}, \emph{anion}, \emph{cation}, \emph{neutre}.
\begin{enumerate}
\item L'atome est constitué d'un \ldots{} central (protons et neutrons) autour duquel gravitent les \ldots{}.
\item Une \ldots{} est un édifice formé de plusieurs atomes liés entre eux.
\item Un ion chargé positivement s'appelle un \ldots{} ; un ion chargé négativement s'appelle un \ldots{}.
\item Tout composé ionique est électriquement \ldots{}.
\end{enumerate}
\end{exercice}

\begin{exercice}[Calculs directs : composition d'atomes]
\begin{enumerate}
\item L'atome d'aluminium possède 13 protons. Combien possède-t-il d'électrons ? Justifie.
\item Le noyau de l'atome de carbone contient 6 protons et 6 neutrons. Quel est son nombre total de nucléons (particules du noyau) ?
\item L'atome d'oxygène possède 8 électrons. En appliquant la règle de remplissage des couches (2 électrons au maximum sur la première couche, 8 sur la deuxième), donne la répartition de ses électrons sur les couches.
\end{enumerate}
\end{exercice}

\courssub{Je m'entraîne}

\begin{exercice}[Lecture de formules chimiques]
Pour chacune des formules suivantes, indique le nom et le nombre d'atomes de chaque élément présent dans la molécule :
\[\ce{CH4} \qquad \ce{C6H12O6} \qquad \ce{Ca(OH)2} \qquad \ce{Al2(SO4)3}\]
\emph{Application :} le glucose \ce{C6H12O6} est le sucre contenu dans les solutés de réhydratation orale utilisés contre la déshydratation.
\end{exercice}

\begin{exercice}[Électrolyse de l'eau]
Au cours de l'électrolyse de l'eau réalisée au laboratoire, on recueille \SI{20}{mL} d'un premier gaz à une électrode et \SI{10}{mL} d'un second gaz à l'autre électrode.
\begin{enumerate}
\item Sachant que le dihydrogène se forme en volume double de celui du dioxygène, identifier chaque gaz.
\item Proposer un test de reconnaissance pour chacun des deux gaz (décrire le geste et l'observation attendue).
\item En déduire la formule de la molécule d'eau et écrire l'équation de la réaction : \ce{2H2O -> 2H2 + O2}. Vérifier la conservation des atomes.
\end{enumerate}
\end{exercice}

\begin{exercice}[Électroneutralité des composés ioniques]
On dispose des cations \ce{Na+}, \ce{Ca^{2+}}, \ce{Al^{3+}} et des anions \ce{Cl-}, \ce{O^{2-}}, \ce{SO4^{2-}}.
\begin{enumerate}
\item Rappeler la règle d'électroneutralité d'un composé ionique.
\item Écrire les formules des composés ioniques formés par :
\begin{enumerate}
\item le sodium avec le chlore (sel de cuisine) ;
\item le calcium avec le chlore ;
\item l'aluminium avec l'oxygène (alumine) ;
\item le calcium avec le sulfate ;
\item l'aluminium avec le sulfate.
\end{enumerate}
\end{enumerate}
\end{exercice}

\begin{exercice}[QCM de sécurité au laboratoire]
Pour chaque situation de TP, choisis le bon comportement et justifie ton choix.
\begin{enumerate}
\item Pour tester la conductivité d'une solution ionique, je :
\begin{enumerate}
\item plonge les électrodes dans la solution sans me laver les mains mouillées ;
\item vérifie d'abord que mes mains sont sèches et que le montage est hors tension avant de le modifier ;
\item goûte la solution pour vérifier qu'elle conduit le courant.
\end{enumerate}
\item Pour manipuler une solution de nitrate d'argent (toxique et tachante), je :
\begin{enumerate}
\item porte une blouse, des lunettes et des gants ;
\item la verse dans l'évier après usage ;
\item la manipule à mains nues car elle est incolore.
\end{enumerate}
\item En cas de projection d'un produit chimique dans l'œil, je :
\begin{enumerate}
\item frotte l'œil avec les mains ;
\item rince abondamment à l'eau et j'appelle le professeur ;
\item attends la fin de la séance.
\end{enumerate}
\end{enumerate}
\end{exercice}

\courssub{Je me prépare au Bac}

\begin{exercice}[Type BEPC : l'atome de magnésium]
Le magnésium est un métal utilisé dans les alliages légers et présent dans certains médicaments contre les brûlures d'estomac vendus en pharmacie à Yaoundé. L'atome de magnésium, de symbole \ce{Mg}, possède \textbf{12 protons} et \textbf{12 neutrons}.
\begin{enumerate}
\item Donner la constitution complète de cet atome : nombre de protons, de neutrons et d'électrons. Justifier le nombre d'électrons.
\item Calculer le nombre de nucléons de cet atome.
\item En appliquant la règle de remplissage des couches (2 électrons au maximum sur la première couche, 8 sur la deuxième), donner la répartition électronique de l'atome de magnésium.
\item Dessiner le modèle de Bohr de cet atome (noyau légendé, couches et électrons représentés par des points).
\item Le magnésium a tendance à perdre les 2 électrons de sa couche externe. Écrire la formule de l'ion qu'il forme et justifier sa charge. S'agit-il d'un cation ou d'un anion ?
\item Écrire la formule du composé ionique formé entre l'ion magnésium et l'ion chlorure \ce{Cl-}, en justifiant par l'électroneutralité.
\end{enumerate}
\end{exercice}

\begin{exercice}[Type BEPC : l'expérience de Rutherford]
En 1911, Ernest Rutherford bombarde une mince feuille d'or avec des particules alpha (chargées positivement). Les observations sont les suivantes :
\begin{itemize}
\item la grande majorité des particules traverse la feuille sans être déviée ;
\item quelques particules sont légèrement déviées ;
\item une particule sur environ 20 000 rebondit vers l'arrière.
\end{itemize}
\begin{enumerate}
\item Schématiser l'expérience : source de particules alpha, feuille d'or, et les trois types de trajectoires observées (légende numérotée).
\item Que signifie le fait que la majorité des particules traverse la feuille sans déviation ? En déduire une première conclusion sur la structure de l'atome.
\item Expliquer pourquoi certaines particules, chargées positivement, rebondissent. En déduire une deuxième conclusion concernant le noyau (taille, charge, masse).
\item L'atome d'or possède 79 protons. Combien possède-t-il d'électrons ? Justifier.
\item L'ion or(III) a pour formule \ce{Au^{3+}}. Dire si cet atome d'or a perdu ou gagné des électrons, et combien.
\end{enumerate}
\end{exercice}

\begin{exercice}[Type BEPC : étude d'un sachet de soluté de réhydratation orale]
Pour lutter contre la déshydratation due au choléra ou aux diarrhées, les centres de santé du Cameroun utilisent des sachets de soluté de réhydratation orale (SRO) à dissoudre dans l'eau. L'étiquette d'un sachet indique la présence des espèces chimiques suivantes après dissolution : ions sodium \ce{Na+}, ions potassium \ce{K+}, ions chlorure \ce{Cl-}, ions citrate \ce{C6H5O7^{3-}} et glucose \ce{C6H12O6}.
\begin{enumerate}
\item Définir les termes \emph{cation} et \emph{anion}.
\item Classer les ions présents dans le soluté en cations et en anions.
\item L'ion citrate est-il un ion simple ou un ion polyatomique ? Justifier.
\item Le glucose \ce{C6H12O6} est-il un ion ou une molécule ? Justifier. Dénombrer les atomes de chaque élément dans cette molécule.
\item Expliquer pourquoi la solution obtenue conduit le courant électrique, alors que l'eau pure ne la conduit pratiquement pas.
\item Le glucose favorise l'absorption du sodium par l'organisme. Proposer une explication simple de l'intérêt de mélanger glucose et sels dans le traitement de la déshydratation.
\item Citer deux règles de sécurité à respecter lors de la préparation du soluté (qualité de l'eau, hygiène des mains et des récipients).
\end{enumerate}
\end{exercice}

\newpage

\coursec{Corrigés des exercices}

\begin{corrige}[Exercice 1 — QCM : l'atome et ses constituants]
\begin{enumerate}
\item \textbf{Réponse b.} Le noyau d'un atome contient des \cle{protons} (chargés positivement) et des \cle{neutrons} (sans charge). Les électrons, eux, gravitent autour du noyau : les réponses a et c sont donc fausses.
\item \textbf{Réponse b.} L'\cle{électron} est une particule chargée \textbf{négativement} (charge $-e$). C'est le proton qui est chargé positivement et le neutron qui est sans charge.
\item \textbf{Réponse b.} Un atome est électriquement neutre parce qu'il contient \textbf{autant de protons (charges $+$) que d'électrons (charges $-$)} : les charges se compensent exactement. La réponse a est fausse car l'atome contient bien des particules chargées ; la réponse c est fausse car le nombre de neutrons n'intervient pas dans la neutralité électrique.
\item \textbf{Réponse a.} Dans le \cle{modèle de Bohr}, les électrons sont répartis sur des \textbf{couches} (ou niveaux d'énergie) autour du noyau, en mouvement permanent. Ils ne sont ni mélangés au noyau, ni immobiles.
\end{enumerate}
\end{corrige}

\begin{corrige}[Exercice 2 — Vrai ou faux, en justifiant]
\begin{enumerate}
\item \textbf{Faux.} La formule \ce{H2O} signifie que la molécule d'eau contient \textbf{deux atomes d'hydrogène et un atome d'oxygène} (l'indice 2 se rapporte à l'élément H qui le précède ; l'absence d'indice après O signifie un seul atome d'oxygène).
\item \textbf{Vrai.} L'atome de sodium possède 11 protons et 11 électrons ; en perdant un électron (charge $-$), il lui reste 11 protons et 10 électrons, soit une charge globale $+1$ : c'est l'ion \ce{Na+}.
\item \textbf{Faux.} L'ion \ce{Cl-} est chargé négativement : c'est un \textbf{anion}. Un cation est un ion chargé positivement (comme \ce{Na+}).
\item \textbf{Vrai.} L'ion sulfate \ce{SO4^{2-}} est formé de \textbf{plusieurs atomes} (1 atome de soufre et 4 atomes d'oxygène) portant une charge globale $2-$ : c'est bien un ion polyatomique.
\item \textbf{Vrai.} Dans \ce{NaCl}, chaque cation \ce{Na+} (charge $+1$) compense exactement un anion \ce{Cl-} (charge $-1$) : la somme des charges est nulle, le composé ionique est électriquement neutre.
\end{enumerate}
\end{corrige}

\begin{corrige}[Exercice 3 — Définitions à compléter]
\begin{enumerate}
\item L'atome est constitué d'un \textbf{noyau} central (protons et neutrons) autour duquel gravitent les \textbf{électrons}.
\item Une \textbf{molécule} est un édifice formé de plusieurs atomes liés entre eux.
\item Un ion chargé positivement s'appelle un \textbf{cation} ; un ion chargé négativement s'appelle un \textbf{anion}.
\item Tout composé ionique est électriquement \textbf{neutre}.
\end{enumerate}
\end{corrige}

\begin{corrige}[Exercice 4 — Calculs directs : composition d'atomes]
\begin{enumerate}
\item L'atome d'aluminium possède \textbf{13 électrons}. Justification : un atome est électriquement neutre, donc il contient autant d'électrons que de protons, soit 13.
\item Les nucléons sont les particules du noyau : protons + neutrons.
\[ N = 6 + 6 = 12 \ \text{nucléons.} \]
\item L'atome d'oxygène possède 8 électrons. On remplit d'abord la première couche (2 électrons au maximum), puis la deuxième :
\[ 8 = 2 + 6 \]
Répartition : \textbf{2 électrons sur la première couche et 6 électrons sur la deuxième couche}.
\end{enumerate}
\end{corrige}

\begin{corrige}[Exercice 5 — Lecture de formules chimiques]
Règle de lecture : l'indice placé en bas à droite d'un symbole multiplie le nombre d'atomes de cet élément ; un indice placé après une parenthèse multiplie \emph{tous} les atomes de la parenthèse.
\begin{itemize}
\item \ce{CH4} (méthane, gaz des marais) : 1 atome de \textbf{carbone} et 4 atomes d'\textbf{hydrogène}.
\item \ce{C6H12O6} (glucose) : 6 atomes de \textbf{carbone}, 12 atomes d'\textbf{hydrogène} et 6 atomes d'\textbf{oxygène}, soit 24 atomes au total.
\item \ce{Ca(OH)2} (hydroxyde de calcium) : l'indice 2 s'applique à toute la parenthèse $(OH)$. Donc : 1 atome de \textbf{calcium}, $2 \times 1 = 2$ atomes d'\textbf{oxygène} et $2 \times 1 = 2$ atomes d'\textbf{hydrogène}.
\item \ce{Al2(SO4)3} (sulfate d'aluminium) : 2 atomes d'\textbf{aluminium} ; l'indice 3 multiplie la parenthèse $(SO_4)$ : $3 \times 1 = 3$ atomes de \textbf{soufre} et $3 \times 4 = 12$ atomes d'\textbf{oxygène}.
\end{itemize}
\end{corrige}

\begin{attention}[Piège fréquent : parenthèses dans les formules]
Dans \ce{Ca(OH)2}, il n'y a pas « un atome d'oxygène et deux atomes d'hydrogène » : l'indice 2 après la parenthèse multiplie \textbf{à la fois} O et H. On obtient 2 O et 2 H. De même, \ce{Al2(SO4)3} contient 3 atomes de soufre et 12 atomes d'oxygène, pas 4 ni 3.
\end{attention}

\begin{corrige}[Exercice 6 — Électrolyse de l'eau]
\begin{enumerate}
\item Le dihydrogène \ce{H2} se forme en volume double du dioxygène \ce{O2}. Donc :
\begin{itemize}
\item le gaz recueilli en plus grand volume (\SI{20}{mL}) est le \textbf{dihydrogène} \ce{H2} (à la cathode, électrode reliée au pôle $-$) ;
\item le gaz recueilli en plus petit volume (\SI{10}{mL}) est le \textbf{dioxygène} \ce{O2} (à l'anode, électrode reliée au pôle $+$).
\end{itemize}
On vérifie : $20 = 2 \times 10$, le rapport des volumes est bien de 2.
\item Tests de reconnaissance :
\begin{itemize}
\item \textbf{Dihydrogène} : on approche une allumette enflammée de l'ouverture du tube ; on entend une \textbf{petite détonation} (un « aboiement »).
\item \textbf{Dioxygène} : on présente une \textbf{bûchette incandescente} (pointe rougie, sans flamme) à l'ouverture du tube ; elle \textbf{se rallume} et brûle vivement.
\end{itemize}
\item L'électrolyse de l'eau produit uniquement du dihydrogène et du dioxygène : la molécule d'eau ne contient donc que des atomes d'hydrogène et d'oxygène. Comme le volume de dihydrogène est le double de celui de dioxygène, il y a deux fois plus d'atomes d'hydrogène que d'atomes d'oxygène : la formule de l'eau est \ce{H2O}.
Équation de la réaction :
\[ \ce{2H2O -> 2H2 + O2} \]
Vérification de la conservation des atomes :
\begin{itemize}
\item hydrogène : $2 \times 2 = 4$ atomes à gauche ; $2 \times 2 = 4$ atomes à droite ;
\item oxygène : $2 \times 1 = 2$ atomes à gauche ; $1 \times 2 = 2$ atomes à droite.
\end{itemize}
Le nombre d'atomes de chaque élément est le même avant et après : l'équation est bien équilibrée (loi de conservation de la matière).
\end{enumerate}
\end{corrige}

\begin{corrige}[Exercice 7 — Électroneutralité des composés ioniques]
\begin{enumerate}
\item \textbf{Règle d'électroneutralité} : un composé ionique est électriquement neutre ; la somme des charges positives des cations doit être égale (en valeur absolue) à la somme des charges négatives des anions. On choisit donc les nombres d'ions de sorte que les charges se compensent exactement.
\item Formules des composés :
\begin{enumerate}
\item Sodium + chlore : \ce{Na+} ($+1$) avec \ce{Cl-} ($-1$) : un ion de chaque suffit. Formule : \ce{NaCl}.
\item Calcium + chlore : \ce{Ca^{2+}} ($+2$) avec \ce{Cl-} ($-1$) : il faut 2 ions chlorure pour 1 ion calcium ($+2$ et $2 \times (-1) = -2$). Formule : \ce{CaCl2}.
\item Aluminium + oxygène : \ce{Al^{3+}} ($+3$) avec \ce{O^{2-}} ($-2$) : le plus petit commun multiple de 3 et 2 est 6 ; il faut 2 ions aluminium ($2 \times (+3) = +6$) et 3 ions oxygène ($3 \times (-2) = -6$). Formule : \ce{Al2O3}.
\item Calcium + sulfate : \ce{Ca^{2+}} ($+2$) avec \ce{SO4^{2-}} ($-2$) : les charges sont égales et opposées, un ion de chaque suffit. Formule : \ce{CaSO4}.
\item Aluminium + sulfate : \ce{Al^{3+}} ($+3$) avec \ce{SO4^{2-}} ($-2$) : il faut 2 ions aluminium ($+6$) et 3 ions sulfate ($-6$). Formule : \ce{Al2(SO4)3} (parenthèses obligatoires car l'ion sulfate est polyatomique et apparaît 3 fois).
\end{enumerate}
\end{enumerate}
\end{corrige}

\begin{attention}[Points de vigilance : écriture des formules ioniques]
Ne jamais écrire les charges dans la formule finale du composé (\ce{NaCl} et non \ce{Na+Cl-}). Mettre des parenthèses autour d'un ion polyatomique dès qu'il est en plusieurs exemplaires : \ce{Al2(SO4)3}, pas \ce{Al2SO43}.
\end{attention}

\begin{corrige}[Exercice 8 — QCM de sécurité au laboratoire]
\begin{enumerate}
\item \textbf{Réponse b.} Avant toute manipulation d'un montage électrique, on vérifie que les mains sont sèches (l'eau, surtout salée, conduit le courant et augmente le risque d'électrisation) et que le montage est hors tension avant de le modifier. Goûter une solution (réponse c) est absolument interdit au laboratoire.
\item \textbf{Réponse a.} Le nitrate d'argent est toxique et tache la peau : on porte obligatoirement une \textbf{blouse}, des \textbf{lunettes de protection} et des \textbf{gants}. On ne verse jamais un produit chimique dans l'évier (réponse b) : on le récupère dans un bidon de récupération prévu à cet effet. « Incolore » ne signifie pas « inoffensif » (réponse c).
\item \textbf{Réponse b.} En cas de projection dans l'œil, on rince \textbf{abondamment et immédiatement à l'eau} (pendant plusieurs minutes, paupières ouvertes) et on \textbf{appelle le professeur}. Frotter l'œil aggrave la lésion ; attendre la fin de la séance laisse le produit agir : chaque seconde compte.
\end{enumerate}
\end{corrige}

\begin{corrige}[Exercice 9 — Type BEPC : l'atome de magnésium]
\begin{enumerate}
\item Constitution de l'atome de magnésium :
\begin{itemize}
\item protons : 12 (donné) ;
\item neutrons : 12 (donné) ;
\item électrons : \textbf{12}, car un atome est électriquement neutre : il possède autant d'électrons (charges $-$) que de protons (charges $+$).
\end{itemize}
\item Nombre de nucléons (particules du noyau) :
\[ N = \text{protons} + \text{neutrons} = 12 + 12 = 24 \ \text{nucléons.} \]
\item Répartition des 12 électrons : on remplit la première couche (2 maximum), puis la deuxième (8 maximum), puis la troisième avec le reste :
\[ 12 = 2 + 8 + 2 \]
Répartition : \textbf{2 électrons sur la 1\iere{} couche, 8 sur la 2\ieme{}, 2 sur la 3\ieme{}} (couche externe).
\item Modèle de Bohr de l'atome de magnésium :
\begin{popfigure}
\begin{center}
\begin{tikzpicture}[scale=0.9]
\fill[poporange] (0,0) circle (0.45);
\node[white,font=\footnotesize\bfseries] at (0,0) {noyau};
\draw[popblue,thick] (0,0) circle (1.0);
\draw[popblue,thick] (0,0) circle (1.8);
\draw[popblue,thick] (0,0) circle (2.6);
\foreach \a in {0,180} \fill[popink] (\a:1.0) circle (0.09);
\foreach \a in {0,45,90,135,180,225,270,315} \fill[popink] (\a:1.8) circle (0.09);
\foreach \a in {90,270} \fill[popink] (\a:2.6) circle (0.09);
\node[popdark,font=\footnotesize] at (3.6,0.6) {12 protons};
\node[popdark,font=\footnotesize] at (3.6,0.1) {12 neutrons};
\draw[->,popmuted] (2.9,0.35) -- (0.55,0.15);
\end{tikzpicture}
\end{center}
\legende{Modèle de Bohr de l'atome de magnésium \ce{Mg} : noyau (12 protons, 12 neutrons) et 12 électrons répartis 2 ; 8 ; 2 sur trois couches.}
\end{popfigure}
\item En perdant ses 2 électrons externes (charges $-$), l'atome conserve 12 protons (charges $+$) mais n'a plus que 10 électrons : la charge globale est $12 - 10 = +2$. L'ion formé est \ce{Mg^{2+}}. Chargé positivement, c'est un \textbf{cation}.
\item L'ion magnésium \ce{Mg^{2+}} porte la charge $+2$ et l'ion chlorure \ce{Cl-} la charge $-1$. Par électroneutralité, il faut \textbf{2 ions chlorure pour 1 ion magnésium} : $(+2) + 2 \times (-1) = 0$. Formule du chlorure de magnésium : \ce{MgCl2}.
\end{enumerate}
\textbf{Barème indicatif (sur 10)} : constitution justifiée (2 pts) ; nucléons (1 pt) ; répartition électronique (2 pts) ; schéma légendé (2 pts) ; ion \ce{Mg^{2+}} justifié, cation (2 pts) ; formule \ce{MgCl2} justifiée (1 pt).

\textbf{Points de vigilance} : la justification « atome neutre $\Rightarrow$ autant d'électrons que de protons » est exigée ; le schéma doit comporter une légende (noyau, couches, électrons) ; la charge de l'ion doit être justifiée par le calcul $12 - 10 = +2$, pas seulement énoncée.
\end{corrige}

\begin{corrige}[Exercice 10 — Type BEPC : l'expérience de Rutherford]
\begin{enumerate}
\item Schéma de l'expérience :
\begin{popfigure}
\begin{center}
\begin{tikzpicture}[scale=1.0]
\draw[fill=popmuted!40] (-4.5,-0.4) rectangle (-3.5,0.4);
\node[font=\footnotesize,popdark] at (-4.0,-0.8) {1. Source $\alpha$};
\draw[fill=popgold!60,popgold] (0,-1.6) rectangle (0.18,1.6);
\node[font=\footnotesize,popdark] at (1.15,-1.85) {2. Feuille d'or};
\draw[->,thick,popblue] (-3.4,0) -- (4.0,0);
\draw[->,thick,popgreen] (-3.4,0.7) -- (-0.1,0.7) -- (4.0,1.1);
\draw[->,thick,poporange] (-3.4,-0.7) -- (-0.05,-0.7) -- (-2.2,-1.5);
\node[font=\footnotesize,popblue] at (3.2,-0.35) {3. Non déviée};
\node[font=\footnotesize,popgreen] at (3.4,1.35) {4. Déviée};
\node[font=\footnotesize,poporange] at (-2.4,-1.75) {5. Renvoyée};
\end{tikzpicture}
\end{center}
\legende{Expérience de Rutherford : 1. source de particules alpha ; 2. feuille d'or ; 3. trajectoire non déviée (majorité) ; 4. trajectoire déviée ; 5. particule renvoyée (très rare).}
\end{popfigure}
\item Si la grande majorité des particules traverse la feuille sans déviation, c'est qu'elles ne rencontrent aucun obstacle sur leur trajet. \textbf{Conclusion 1} : l'atome est essentiellement constitué de \textbf{vide} ; la matière a une structure lacunaire.
\item Les particules alpha sont chargées positivement. Si certaines rebondissent, c'est qu'elles ont été violemment repoussées par un « obstacle » lui aussi chargé \textbf{positivement} (les charges de même signe se repoussent). Comme ce rebond est très rare (1 particule sur 20 000), cet obstacle est \textbf{très petit} devant la taille de l'atome ; et pour renvoyer une particule lancée à grande vitesse, il doit concentrer \textbf{presque toute la masse} de l'atome. \textbf{Conclusion 2} : l'atome possède un \textbf{noyau} très petit, chargé positivement, qui concentre l'essentiel de la masse.
\item L'atome d'or possède \textbf{79 électrons}, car un atome est électriquement neutre : nombre d'électrons = nombre de protons = 79.
\item L'ion \ce{Au^{3+}} porte 3 charges positives de plus que l'atome neutre : l'atome d'or a donc \textbf{perdu 3 électrons} (il lui reste 76 électrons pour 79 protons : $79 - 76 = +3$).
\end{enumerate}
\textbf{Barème indicatif (sur 10)} : schéma légendé avec les trois trajectoires (2 pts) ; conclusion « atome lacunaire » (2 pts) ; explication du rebond et conclusion sur le noyau, taille, charge, masse (3 pts) ; 79 électrons justifiés (1,5 pt) ; perte de 3 électrons (1,5 pt).

\textbf{Points de vigilance} : il faut explicitement relier chaque observation à sa conclusion (démarche d'interprétation) ; pour le rebond, citer la répulsion entre charges de même signe ; pour la question 4, la justification par la neutralité électrique de l'atome est indispensable.
\end{corrige}

\begin{corrige}[Exercice 11 — Type BEPC : le soluté de réhydratation orale]
\begin{enumerate}
\item Un \cle{cation} est un ion chargé \textbf{positivement} (atome ou groupe d'atomes ayant perdu un ou plusieurs électrons). Un \cle{anion} est un ion chargé \textbf{négativement} (atome ou groupe d'atomes ayant gagné un ou plusieurs électrons).
\item Classement des ions du soluté :
\begin{itemize}
\item \textbf{Cations} (charges $+$) : ion sodium \ce{Na+} et ion potassium \ce{K+} ;
\item \textbf{Anions} (charges $-$) : ion chlorure \ce{Cl-} et ion citrate \ce{C6H5O7^{3-}}.
\end{itemize}
\item L'ion citrate \ce{C6H5O7^{3-}} est un ion \textbf{polyatomique} : il est formé de \textbf{plusieurs atomes} liés entre eux (6 atomes de carbone, 5 d'hydrogène, 7 d'oxygène) portant globalement la charge $3-$.
\item Le glucose \ce{C6H12O6} est une \textbf{molécule} : sa formule ne porte \textbf{aucune charge} électrique, c'est un édifice neutre formé d'atomes liés. Composition : 6 atomes de carbone, 12 atomes d'hydrogène et 6 atomes d'oxygène (24 atomes au total).
\item Une solution ne conduit le courant que si elle contient des \textbf{porteurs de charge mobiles}. La solution de SRO contient des ions (\ce{Na+}, \ce{K+}, \ce{Cl-}, \ce{C6H5O7^{3-}}) libres de se déplacer dans l'eau : elle conduit le courant électrique. L'eau pure ne contient pratiquement pas d'ions : elle conduit très mal le courant.
\item La déshydratation fait perdre à l'organisme à la fois de l'\textbf{eau} et des \textbf{sels} (ions sodium notamment). Le glucose favorise l'absorption du sodium par l'intestin, et l'eau suit le sodium absorbé : le mélange glucose + sels permet donc de faire rentrer simultanément l'eau et les sels dans l'organisme, ce qui réhydrate efficacement le malade.
\item Deux règles de sécurité :
\begin{itemize}
\item utiliser une \textbf{eau propre et sûre} (eau bouillie puis refroidie, ou eau traitée) pour dissoudre le sachet ;
\item se \textbf{laver les mains} et utiliser des \textbf{récipients propres} pour préparer et conserver le soluté (afin de ne pas introduire de nouveaux microbes).
\end{itemize}
\end{enumerate}
\textbf{Barème indicatif (sur 10)} : définitions cation/anion (1,5 pt) ; classement (1,5 pt) ; ion polyatomique justifié (1 pt) ; molécule justifiée et dénombrement (2 pts) ; conduction expliquée par les ions mobiles (2 pts) ; intérêt du mélange glucose-sels (1 pt) ; règles de sécurité (1 pt).

\textbf{Points de vigilance} : distinguer clairement molécule (neutre) et ion (chargé) ; la justification « polyatomique » doit mentionner la présence de plusieurs atomes ; pour la conduction électrique, le mot-clé attendu est « ions \textbf{mobiles} en solution ».
\end{corrige}

\newpage

\coursec{Fiche bilan}

\begin{popfigure}
\begin{tikzpicture}[mindmap, grow cyclic, every node/.style={concept, font=\small},
  root concept/.append style={concept color=popblue, font=\bfseries, text=white, minimum size=2.6cm},
  level 1/.append style={level distance=3.4cm, sibling angle=60, font=\footnotesize},
  level 2/.append style={level distance=2.2cm, sibling angle=45, font=\scriptsize}]
\node[root concept]{Constituants de la matière}
  child[concept color=poporange]{ node {L'atome}
    child[concept color=poporangeL]{ node {Noyau : protons $+$, neutrons} }
    child[concept color=poporangeL]{ node {\'Electrons $-$ en couches} }
  }
  child[concept color=popgreen]{ node {Symbole et Bohr}
    child[concept color=popgreenL]{ node {\ce{^{A}_{Z}X} : $Z$ protons} }
    child[concept color=popgreenL]{ node {Couches K, L, M} }
  }
  child[concept color=poppurple]{ node {Molécules}
    child[concept color=poppurpleL]{ node {Assemblage d'atomes} }
    child[concept color=poppurpleL]{ node {Formule : \ce{H2O}, \ce{CO2}} }
  }
  child[concept color=poppink]{ node {Ions simples}
    child[concept color=poppinkL]{ node {Cations : \ce{Na+}, \ce{Ca^{2+}}} }
    child[concept color=poppinkL]{ node {Anions : \ce{Cl-}, \ce{O^{2-}}} }
  }
  child[concept color=popgold]{ node {Ions polyatomiques}
    child[concept color=popgoldL]{ node {\ce{SO4^{2-}}, \ce{OH-}, \ce{NH4+}} }
  }
  child[concept color=popmuted]{ node {Composés ioniques}
    child[concept color=popmuted!60]{ node {Neutres : \ce{NaCl}, \ce{CaSO4}} }
  };
\end{tikzpicture}
\legende{Carte mentale de la leçon : de l'atome aux composés ioniques.}
\end{popfigure}

\begin{bilan}
\textbf{Définitions clés}
\begin{itemize}
  \item \cle{Atome} : plus petite partie d'un corps simple, constituée d'un \cle{noyau} (protons chargés $+$ et neutrons neutres) et d'\cle{électrons} (charge $-$) qui gravitent autour.
  \item L'atome est électriquement \cle{neutre} : nombre d'électrons = nombre de protons.
  \item \cle{Molécule} : assemblage électriquement neutre d'atomes liés entre eux (ex. \ce{H2O}, \ce{O2}, \ce{CO2}).
  \item \cle{Ion} : atome ou groupe d'atomes ayant gagné ou perdu un ou plusieurs électrons. \cle{Cation} : ion positif (perte d'électrons) ; \cle{anion} : ion négatif (gain d'électrons).
  \item \cle{Composé ionique} : solide formé de cations et d'anions, électriquement neutre au total.
\end{itemize}

\textbf{Formules et représentations à connaître par c\oe ur}
\begin{center}
\begin{tabularx}{\linewidth}{|l|X|X|}
\hline
\rowcolor{popblueL} \textbf{Notion} & \textbf{Écriture} & \textbf{Exemple} \\
\hline
Noyau d'un atome & \ce{^{A}_{Z}X} ($Z$ protons, $A$ nucléons) & \ce{^{12}_{6}C} : 6 protons, 6 neutrons \\
\hline
Modèle de Bohr & Couches K ($2\,e^-$), L ($8\,e^-$), M ($18\,e^-$ max) & Na ($Z=11$) : (K)$^2$(L)$^8$(M)$^1$ \\
\hline
Ions simples & \ce{X^{n+}} ou \ce{X^{n-}} & \ce{Na+}, \ce{Mg^{2+}}, \ce{Cl-}, \ce{O^{2-}} \\
\hline
Ions polyatomiques & groupe d'atomes chargé & \ce{OH-}, \ce{SO4^{2-}}, \ce{NO3-}, \ce{NH4+}, \ce{CO3^{2-}} \\
\hline
Composé ionique & neutralité électrique & \ce{NaCl}, \ce{CaCl2}, \ce{Na2SO4} \\
\hline
\end{tabularx}
\end{center}

\textbf{Méthodes express}
\begin{itemize}
  \item Pour trouver la structure électronique : remplir K, puis L, puis M, dans la limite de $Z$ électrons.
  \item Pour écrire un composé ionique : ajuster les nombres d'ions pour que la charge totale soit nulle (ex. \ce{Ca^{2+}} + $2\,\ce{Cl-}$ $\rightarrow$ \ce{CaCl2}).
  \item Pour identifier un ion : solution de nitrate d'argent pour \ce{Cl-} (précipité blanc), soude pour \ce{Cu^{2+}} (précipité bleu).
\end{itemize}
\end{bilan}

\begin{aretenir}[Checklist avant le BEPC]
\begin{itemize}
  \item[$\square$] Je sais décrire la constitution d'un atome (noyau, protons, neutrons, électrons).
  \item[$\square$] Je sais que l'atome est électriquement neutre.
  \item[$\square$] Je sais lire le symbole \ce{^{A}_{Z}X} et trouver le nombre de protons, neutrons, électrons.
  \item[$\square$] Je sais dessiner le modèle de Bohr d'un atome simple ($Z \leq 20$).
  \item[$\square$] Je sais distinguer atome, molécule et ion.
  \item[$\square$] Je sais écrire et lire une formule moléculaire (\ce{H2O}, \ce{CO2}, \ce{O2}).
  \item[$\square$] Je connais les ions simples usuels : \ce{Na+}, \ce{K+}, \ce{Ca^{2+}}, \ce{Mg^{2+}}, \ce{Cl-}, \ce{O^{2-}}.
  \item[$\square$] Je connais les ions polyatomiques usuels : \ce{OH-}, \ce{SO4^{2-}}, \ce{NO3-}, \ce{NH4+}.
  \item[$\square$] Je sais écrire la formule d'un composé ionique neutre (\ce{NaCl}, \ce{CaCl2}).
  \item[$\square$] Je sais citer un test d'identification d'ions en solution.
  \item[$\square$] Je respecte les consignes de sécurité au laboratoire (blouse, pas de goût, gants pour les produits corrosifs).
\end{itemize}
\end{aretenir}

\findecours

\end{document}
